\documentclass[11pt,openany]{book}

% -----------------------------------------------------------------------------
% Polynomial-Logarithmic Transseries
% A self-contained treatment of arithmetic, composition, and series reversal.
% -----------------------------------------------------------------------------

\usepackage[T1]{fontenc}
\usepackage[utf8]{inputenc}
\usepackage{microtype}
\usepackage{geometry}
\geometry{
  paperwidth=7.25in,
  paperheight=10.25in,
  inner=0.86in,
  outer=0.72in,
  top=0.78in,
  bottom=0.84in,
  headsep=0.22in,
  footskip=0.38in
}
\usepackage{amsmath,amssymb,amsthm,mathtools}
% Canonical mathematical notation for active FabiusFunction LaTeX documents.
%
% This file is intentionally a plain TeX input rather than a LaTeX package.
% Every standalone document loads it by an explicit path relative to that
% document's directory; included fragments inherit the definitions from their
% root document.  Keep this file independent of document layout and metadata.
%
% Contract:
%   * one command has one mathematical meaning throughout the active corpus;
%   * names encode semantic roles rather than a document-local choice of letter;
%   * visually similar but differently normalized objects have different print;
%   * every command below is defined normatively in the notation catalogue.
%
% The active corpus excludes docs/archive/.  Historical sources may retain
% their original notation only inside an explicitly labelled quotation.

\makeatletter
\@ifundefined{mathbb}{%
  \PackageError{fabius-notation}{The amssymb package must be loaded first}%
    {Load amsmath and amssymb before inputting fabius-notation.tex.}%
}{}
\@ifundefined{operatorname}{%
  \PackageError{fabius-notation}{The amsmath package must be loaded first}%
    {Load amsmath before inputting fabius-notation.tex.}%
}{}
\makeatother

% Version stamp used by the catalogue and by static audits.
\newcommand{\FabiusNotationVersion}{2026-08-31}

% Number systems.  NaturalNumbers is explicitly the nonnegative convention.
\newcommand{\NaturalNumbers}{\mathbb{N}_{0}}
\newcommand{\PositiveIntegers}{\mathbb{N}_{>0}}
\newcommand{\IntegerNumbers}{\mathbb{Z}}
\newcommand{\RationalNumbers}{\mathbb{Q}}
\newcommand{\RealNumbers}{\mathbb{R}}
\newcommand{\ComplexNumbers}{\mathbb{C}}
\newcommand{\ExtendedRealNumbers}{\overline{\mathbb{R}}}

% The three principal Fabius--Rvachev functions.  Subscripts are deliberate:
% a bare F and a calligraphic F have several unrelated uses in the corpus.
\newcommand{\FabiusBounded}{F_{\mathrm{bd}}}
\newcommand{\FabiusGlobal}{F_{\mathrm{gl}}}
\newcommand{\FabiusQuantile}{Q_{F}}
\newcommand{\FabiusClampedQuantile}{Q_{F,\mathrm{cl}}}
\newcommand{\RvachevUp}{\operatorname{up}}

% Constants, elementary operators, and delimiters.
\newcommand{\EulerE}{\mathrm{e}}
\newcommand{\ImaginaryUnit}{\mathrm{i}}
\newcommand{\Differential}{\mathop{}\!\mathrm{d}}
\newcommand{\AbsoluteValue}[1]{\left\lvert #1\right\rvert}
\newcommand{\Norm}[1]{\left\lVert #1\right\rVert}
\newcommand{\Floor}[1]{\left\lfloor #1\right\rfloor}
\newcommand{\Ceiling}[1]{\left\lceil #1\right\rceil}
\newcommand{\FractionalPart}[1]{\left\{#1\right\}}
\newcommand{\PositivePart}[1]{\left(#1\right)_{+}}
\newcommand{\IndicatorOf}[1]{\mathbf{1}_{#1}}
\newcommand{\SetOf}[1]{\left\{#1\right\}}
\newcommand{\InnerProduct}[1]{\left\langle #1\right\rangle}
\newcommand{\CardinalityOf}[1]{\left\lvert #1\right\rvert}
\newcommand{\DefinitionEquals}{\mathrel{\coloneqq}}
\newcommand{\Sign}{\operatorname{sgn}}
\newcommand{\IdentityOperator}{\operatorname{id}}
\newcommand{\SpanOperator}{\operatorname{span}}
\newcommand{\TraceOperator}{\operatorname{tr}}
\newcommand{\RankOperator}{\operatorname{rank}}
\newcommand{\DiagonalOperator}{\operatorname{diag}}
\newcommand{\ResidueOperator}{\operatorname*{Res}}
\newcommand{\LeastCommonMultiple}{\operatorname{lcm}}
\newcommand{\OrderOperator}{\operatorname{ord}}
\newcommand{\PrincipalLogarithm}{\operatorname{Log}}
\newcommand{\FallingFactorial}[2]{\left(#1\right)^{\underline{#2}}}
\newcommand{\RisingFactorial}[2]{\left(#1\right)^{\overline{#2}}}

% Support, topology, convolution, and metric notation.
\newcommand{\SupportOperator}{\operatorname{supp}}
\newcommand{\SupportOf}[1]{\SupportOperator\!\left(#1\right)}
\newcommand{\TopologicalSupportOf}[1]{\operatorname{tsupp}\!\left(#1\right)}
\newcommand{\Convolution}{\mathbin{*}}
\newcommand{\MetricDistance}{\operatorname{dist}}
\newcommand{\TotalVariationDistance}{d_{\mathrm{TV}}}

% Probability.  Equality and convergence in law are not metric distance.
\newcommand{\Probability}{\mathbb{P}}
\newcommand{\Expectation}{\mathbb{E}}
\newcommand{\Variance}{\operatorname{Var}}
\newcommand{\Covariance}{\operatorname{Cov}}
\newcommand{\LawOperator}{\operatorname{Law}}
\newcommand{\LawOf}[1]{\LawOperator\!\left(#1\right)}
\newcommand{\EqualInLaw}{\mathrel{\overset{\mathrm{law}}{=}}}
\newcommand{\ConvergesInLaw}{\mathrel{\xrightarrow{\mathrm{law}}}}

% Fourier and integral transforms.  The printed labels expose normalization.
% FourierTwoPi uses exp(-2*pi*i*x*xi); FourierAngular uses exp(-i*x*omega).
\newcommand{\FourierTwoPi}[1]{\widehat{#1}_{\,2\pi}}
\newcommand{\FourierAngular}[1]{\widehat{#1}_{\,\mathrm{ang}}}
\newcommand{\LaplaceTransformSymbol}{\mathcal{L}}
\newcommand{\MellinTransformSymbol}{\mathcal{M}}
\newcommand{\LaplaceTransformOf}[1]{\LaplaceTransformSymbol\!\left[#1\right]}
\newcommand{\MellinTransformOf}[1]{\MellinTransformSymbol\!\left[#1\right]}
\newcommand{\MomentGeneratingFunction}[1]{M_{#1}}
\newcommand{\CharacteristicFunction}[1]{\varphi_{#1}}

% The two standard sinc functions must never share one symbol.
\newcommand{\SincRad}{\operatorname{sinc}_{\mathrm{rad}}}
\newcommand{\SincPi}{\operatorname{sinc}_{\pi}}

% Binary arithmetic and Thue--Morse notation.
\newcommand{\BinaryDigitSum}{\operatorname{wt}_{2}}
\newcommand{\ThueMorseSignSymbol}{\varepsilon}
\newcommand{\ThueMorseSign}[1]{\varepsilon_{#1}}
\newcommand{\PadicValuation}[1]{\nu_{#1}}
\newcommand{\TwoAdicValuation}{\nu_{2}}

% q-series notation.  Every base is an explicit argument.
\newcommand{\QPochhammer}[3]{\left(#1;#2\right)_{#3}}
\newcommand{\GaussianBinomial}[3]{%
  \genfrac{[}{]}{0pt}{}{#1}{#2}_{#3}%
}
\newcommand{\QInteger}[2]{\left[#1\right]_{#2}}

% Named special functions.
\newcommand{\Polylogarithm}{\operatorname{Li}}
\newcommand{\LambertW}{W}
\newcommand{\GammaFunction}{\Gamma}
\newcommand{\RiemannZeta}{\zeta}

% Asymptotics.  BigO and LittleO are operator tokens for existing formula
% grammar; prefer the At forms so that the limiting regime is printed.
\newcommand{\BigO}{\mathcal{O}}
\newcommand{\LittleO}{o}
\newcommand{\BigOAt}[2]{\mathcal{O}_{#2}\!\left(#1\right)}
\newcommand{\LittleOAt}[2]{o_{#2}\!\left(#1\right)}

\usepackage{newtxtext,newtxmath}
\usepackage{bm}
\usepackage{array,booktabs,longtable,tabularx}
\usepackage{enumitem}
\usepackage{xcolor}
\usepackage[most]{tcolorbox}
\usepackage{listings}
\usepackage{fancyhdr}
\usepackage{titlesec}
\usepackage{bookmark}
\usepackage{xurl}
\usepackage{hyperref}
\usepackage[capitalise,nameinlink,noabbrev]{cleveref}
\usepackage{needspace}

\definecolor{DeepBlue}{HTML}{173A5E}
\definecolor{MediumBlue}{HTML}{2F6690}
\definecolor{PaleBlue}{HTML}{EAF2F8}
\definecolor{DeepTeal}{HTML}{1D5C63}
\definecolor{PaleTeal}{HTML}{E9F5F3}
\definecolor{DeepGold}{HTML}{7A5B13}
\definecolor{PaleGold}{HTML}{FBF5E6}
\definecolor{DeepRed}{HTML}{7A2E2E}
\definecolor{PaleRed}{HTML}{FBECEC}
\definecolor{SoftGray}{HTML}{F4F5F6}
\definecolor{DarkGray}{HTML}{3D4650}
\definecolor{CodeGreen}{HTML}{2E6B3E}

\hypersetup{
  colorlinks=true,
  linkcolor=DeepBlue,
  citecolor=DeepTeal,
  urlcolor=MediumBlue,
  pdftitle={Polynomial-Logarithmic Transseries: Arithmetic, Composition, and Series Reversal},
  pdfauthor={OpenAI},
  pdfsubject={Polynomial-logarithmic asymptotic series and Lambert W},
  pdfkeywords={transseries, polynomial-logarithmic series, Lambert W, asymptotic inversion, composition, series reversion, Stirling numbers}
}

\setcounter{tocdepth}{2}
\setcounter{secnumdepth}{3}

\pagestyle{fancy}
\fancyhf{}
\fancyhead[LE]{\small\scshape\nouppercase{\leftmark}}
\fancyhead[RO]{\small\scshape\nouppercase{\rightmark}}
\fancyfoot[LE,RO]{\thepage}
\renewcommand{\headrulewidth}{0.4pt}
\renewcommand{\footrulewidth}{0pt}
\setlength{\headheight}{14pt}

\titleformat{\chapter}[display]
  {\normalfont\bfseries\color{DeepBlue}}
  {\filleft\Large\chaptername\ \thechapter}
  {0.5ex}
  {\titlerule\vspace{1.0ex}\filright\Huge}
  [\vspace{0.6ex}\titlerule]
\titleformat{\section}{\Large\bfseries\color{DeepBlue}}{\thesection}{0.7em}{}
\titleformat{\subsection}{\large\bfseries\color{DeepTeal}}{\thesubsection}{0.7em}{}
\titleformat{\subsubsection}{\normalsize\bfseries\color{DarkGray}}{\thesubsubsection}{0.7em}{}

\setlist[itemize]{leftmargin=1.4em,itemsep=0.25ex,topsep=0.5ex}
\setlist[enumerate]{leftmargin=1.8em,itemsep=0.35ex,topsep=0.5ex}
\setlength{\parindent}{1.15em}
\setlength{\parskip}{0.15ex}
\raggedbottom
\allowdisplaybreaks[2]

\theoremstyle{definition}
\newtheorem{definition}{Definition}[chapter]
\newtheorem{example}{Example}[chapter]
\newtheorem{algorithm}{Algorithm}[chapter]
\newtheorem{exercise}{Exercise}[chapter]
\theoremstyle{plain}
\newtheorem{theorem}{Theorem}[chapter]
\newtheorem{proposition}{Proposition}[chapter]
\newtheorem{lemma}{Lemma}[chapter]
\newtheorem{corollary}{Corollary}[chapter]
\theoremstyle{remark}
\newtheorem{remark}{Remark}[chapter]

\newtcolorbox{intuitionbox}[1][]{
  enhanced,breakable,colback=PaleBlue,colframe=MediumBlue,
  boxrule=0.6pt,arc=1.2mm,left=1.5mm,right=1.5mm,top=1.1mm,bottom=1.1mm,
  title={Intuition},fonttitle=\bfseries,#1
}
\newtcolorbox{checkpointbox}[1][]{
  enhanced,breakable,colback=PaleTeal,colframe=DeepTeal,
  boxrule=0.6pt,arc=1.2mm,left=1.5mm,right=1.5mm,top=1.1mm,bottom=1.1mm,
  title={Checkpoint},fonttitle=\bfseries,#1
}
\newtcolorbox{warningbox}[1][]{
  enhanced,breakable,colback=PaleRed,colframe=DeepRed,
  boxrule=0.6pt,arc=1.2mm,left=1.5mm,right=1.5mm,top=1.1mm,bottom=1.1mm,
  title={A common trap},fonttitle=\bfseries,#1
}
\newtcolorbox{summarybox}[1][]{
  enhanced,breakable,colback=PaleGold,colframe=DeepGold,
  boxrule=0.6pt,arc=1.2mm,left=1.5mm,right=1.5mm,top=1.1mm,bottom=1.1mm,
  title={Chapter summary},fonttitle=\bfseries,#1
}
\newtcolorbox{workedbox}[1][]{
  enhanced,breakable,colback=SoftGray,colframe=DarkGray,
  boxrule=0.5pt,arc=1.2mm,left=1.5mm,right=1.5mm,top=1.1mm,bottom=1.1mm,
  title={Worked calculation},fonttitle=\bfseries,#1
}
\newtcolorbox{scopebox}[1][]{
  enhanced,breakable,colback=white,colframe=DeepGold,
  boxrule=0.8pt,arc=1.2mm,left=1.5mm,right=1.5mm,top=1.1mm,bottom=1.1mm,
  title={Scope and closure},fonttitle=\bfseries,#1
}

\lstdefinestyle{pythonstyle}{
  language=Python,
  basicstyle=\ttfamily\small,
  keywordstyle=\bfseries\color{DeepBlue},
  commentstyle=\itshape\color{CodeGreen},
  stringstyle=\color{DeepRed},
  columns=fullflexible,
  frame=single,
  rulecolor=\color{DarkGray},
  backgroundcolor=\color{SoftGray},
  showstringspaces=false,
  keepspaces=true,
  xleftmargin=0.5em,
  xrightmargin=0.5em,
  aboveskip=0.8em,
  belowskip=0.8em,
  breaklines=true
}
\lstdefinestyle{pseudostyle}{
  basicstyle=\ttfamily\small,
  columns=fullflexible,
  frame=single,
  rulecolor=\color{DarkGray},
  backgroundcolor=\color{SoftGray},
  showstringspaces=false,
  keepspaces=true,
  xleftmargin=0.5em,
  xrightmargin=0.5em,
  aboveskip=0.8em,
  belowskip=0.8em,
  breaklines=true
}

% Mathematical notation.
\newcommand{\K}{\mathbb{K}}
\newcommand{\val}{\nu}
\newcommand{\rev}[1]{\CompositionalInverseOf{#1}}
\newcommand{\recip}[1]{\MultiplicativeInverseOf{#1}}
\newcommand{\stirlingone}[2]{\genfrac{[}{]}{0pt}{}{#1}{#2}}
\newcommand{\Dop}{\mathrm{D}}
\newcommand{\Deltaop}{\boldsymbol{\Delta}}
\newcommand{\normalexp}{\mathopen{:}\exp\mathclose{:}}
\newcommand{\compose}{\mathbin{\circ}}
\newcommand{\asym}{\sim}
\newcommand{\precdom}{\prec}
\newcommand{\succdom}{\succ}
% LOCAL-NON-FOURIER-HAT: \wideeq formal coefficientwise identity relation
\newcommand{\wideeq}{\mathrel{\widehat{=}}}
% LOCAL-NON-FOURIER-HAT: \DulacFormalSeriesOf formal Dulac-series expansion
\newcommand{\DulacFormalSeriesOf}[1]{\widehat{#1}}

\crefname{definition}{definition}{definitions}
\crefname{example}{example}{examples}
\crefname{algorithm}{algorithm}{algorithms}
\crefname{exercise}{exercise}{exercises}
\crefname{theorem}{theorem}{theorems}
\crefname{proposition}{proposition}{propositions}
\crefname{lemma}{lemma}{lemmas}
\crefname{corollary}{corollary}{corollaries}

\begin{document}

\frontmatter

\begin{titlepage}
\thispagestyle{empty}
\pagecolor{PaleBlue}
\color{DeepBlue}
\vspace*{1.05in}
\begin{center}
  {\rule{0.88\textwidth}{1.1pt}}\\[1.0em]
  {\fontsize{32}{37}\selectfont\bfseries
   Polynomial--Logarithmic\\[0.12em] Transseries}\\[0.9em]
  {\Large Arithmetic, Division, Composition, and Series Reversal}\\[1.0em]
  {\rule{0.64\textwidth}{0.7pt}}\\[1.8em]
  {\large A constructive guide centered on the asymptotic inversion\\
   behind Lambert's \(W\) function}\\[3.2em]
  \begin{tcolorbox}[width=0.82\textwidth,colback=white,colframe=DeepTeal,
    boxrule=0.8pt,arc=1.5mm,left=3mm,right=3mm,top=2mm,bottom=2mm]
  \centering
  From coefficient convolution and reciprocal recurrences to normal-ordered
  composition, Bell-polynomial bookkeeping, Newton reversion, Stirling-number
  formulas, branch-sensitive Lambert expansions, and implementable algorithms.
  \end{tcolorbox}
  \vfill
  {\large OpenAI}\\[0.5em]
  {\normalsize August 2026}
\end{center}
\color{black}
\nopagecolor
\end{titlepage}

\chapter*{Abstract}
\addcontentsline{toc}{chapter}{Abstract}
Polynomial--logarithmic transseries are expansions whose coefficients are
polynomials in a slow logarithmic variable and whose successive levels are
powers of a faster variable.  The model example is
\[
 W(z)=X-\log X+\frac{\log X}{X}
      +\frac{\log X(\log X-2)}{2X^2}+\cdots,
 \qquad X=\PrincipalLogarithm z+2\pi \ImaginaryUnit k,
\]
which inverts \(w\mapsto w\EulerE^w\) on the \(k\)-th branch.  These expansions sit
between ordinary formal power series and the full logarithmic--exponential
transseries field.  They are elementary enough for explicit coefficient
algorithms, yet rich enough to display the essential phenomena of asymptotic
scales, logarithmic coefficient growth, failure of naive closure, and
compositional inversion.

This article develops a self-contained calculus for such series.  It begins
with the ring \(\K[L]((X^{-1}))\), explains its valuation and asymptotic meaning,
and carefully distinguishes it from the larger field
\(\K(L)((X^{-1}))\).  It then treats multiplication, powers, reciprocal and
quotient recurrences, differentiation, logarithms and exponentials, regular and
near-identity composition, Bell-polynomial coefficient extraction, and formal
series reversal.  Universal composition and inverse formulas are displayed
through fourth order.  Newton algorithms are proved to double formal precision.
The Lambert \(W\) expansion is derived twice---by triangular coefficient
matching and by an explicit Stirling-cycle-number formula---and is extended to
all complex branches, the real branch \(W_{-1}\), and equations of the form
\(y+\alpha\log y=X\).  A final implementation chapter turns the theory into
precise data structures, pseudocode, verification tests, and a compact SymPy
prototype.

The emphasis throughout is constructive: every infinite operation is reduced
to a finite calculation at a requested truncation order, every closure claim is
stated with its hypotheses, and every important formula is accompanied by a
residual check.

\chapter*{Reader's guide}
\addcontentsline{toc}{chapter}{Reader's guide}
The basic prerequisites are calculus, ordinary power series, and comfort with
formal algebra.  Familiarity with asymptotic notation is helpful but not
required.  Readers mainly interested in computation may read
Chapters~\ref{ch:coordinates}, \ref{ch:multiplication}, \ref{ch:division},
\ref{ch:composition}, \ref{ch:reversion}, and \ref{ch:lambert} in order, then
jump to Chapter~\ref{ch:implementation}.  Readers interested in the boundary
between this elementary calculus and general transseries should also read
Chapters~\ref{ch:closure} and \ref{ch:dulac}.

Two notational cautions are worth stating immediately.
\begin{enumerate}
\item The symbol \(A^{-1}\) is ambiguous.  We write \(\recip{A}=1/A\) for a
multiplicative inverse and \(\rev{A}\) for a compositional inverse.
\item The variable \(L\) always means \(\log X\) unless a chapter explicitly
switches to a small variable.  It is not an independent variable under
composition: changing \(X\) changes \(L\) as well.
\end{enumerate}

\begin{scopebox}
The phrase \emph{polynomial--logarithmic transseries} is used here for a
practical, explicitly computable subclass of power--log transseries.  The core
class has integer powers of \(X^{-1}\) and polynomial coefficients in
\(L=\log X\).  Later chapters allow rational coefficients in \(L\), generalized
real exponents, and finitely many iterated logarithms.  Full transseries permit
much more, including exponentials of large transseries; the present article
uses that larger theory as context but does not require it.
\end{scopebox}

\tableofcontents

\mainmatter

\chapter{Why polynomial--logarithmic transseries appear}\label{ch:motivation}

\section{The first logarithm forced by inversion}
Suppose that a large positive quantity \(Y\) is related to another large
quantity \(X\) by
\begin{equation}\label{eq:first-implicit}
  Y+\log Y=X.
\end{equation}
If the logarithm were absent, the answer would simply be \(Y=X\).  Since
\(\log Y\) is much smaller than \(Y\), a first correction is
\[
  Y\approx X-\log X.
\]
But this is not exact.  Substituting it into the left side of
\cref{eq:first-implicit} gives
\[
 X-\log X+\log(X-\log X)
 =X+\log\!\left(1-\frac{\log X}{X}\right).
\]
Expanding the last logarithm produces
\[
 -\frac{\log X}{X}
 -\frac{(\log X)^2}{2X^2}
 -\frac{(\log X)^3}{3X^3}-\cdots.
\]
The correction to \(X-\log X\) therefore involves powers of \(X^{-1}\), and
their coefficients involve powers of \(\log X\).  Continuing the calculation
gives
\begin{equation}\label{eq:W-preview}
\begin{split}
Y={}&X-L+\frac{L}{X}
 +\frac{L(L-2)}{2X^2}
 +\frac{L(2L^2-9L+6)}{6X^3}\\
&+\frac{L(3L^3-22L^2+36L-12)}{12X^4}+\cdots,
\qquad L=\log X.
\end{split}
\end{equation}
Equation~\eqref{eq:first-implicit} is exactly the logarithmic form of the
Lambert equation \(Y\EulerE^Y=z\), with \(X=\log z\).  Thus
\cref{eq:W-preview} is the large-argument expansion of Lambert's \(W\)
function.

\begin{intuitionbox}
An inverse expansion remembers every asymptotic scale that was needed to undo
the original function.  The dominant term \(Y\) in \(Y+\log Y\) creates the
scale \(X\); the smaller term \(\log Y\) creates the scale \(\log X\); repeated
correction creates powers of \(X^{-1}\) carrying polynomial dependence on
\(\log X\).
\end{intuitionbox}

\section{A scale between power series and full transseries}
Ordinary power series use one small monomial, say \(t\):
\[
 a_0+a_1t+a_2t^2+\cdots.
\]
At infinity we may take \(t=X^{-1}\).  A polynomial--logarithmic series allows
the coefficient of each \(t^n\) to be a polynomial in the slower variable
\(L=\log X\):
\begin{equation}\label{eq:basic-pl-form}
 A(X)=\sum_{n\ge n_0}a_n(L)X^{-n},
 \qquad a_n(L)\in\K[L].
\end{equation}
Examples include
\[
 X-L+\frac{L}{X},\qquad
 1+\frac{L^2-3L+1}{X^2},\qquad
 X^{3/2}\left(1+\frac{2L-1}{X}+\cdots\right).
\]
This is already a transseries in the modest sense that it combines more than
one asymptotic scale.  For every fixed integer \(m\),
\[
 X^{-1}\prec L^{-m}\prec 1\prec L^m\prec X
 \qquad (X\to+\infty),
\]
where \(f\prec g\) means \(f/g\to0\).  Thus one step in the power of \(X\)
dominates every finite change in the power of \(L\).

Full logarithmic--exponential transseries allow arbitrary well-based sums of
monomials built from powers, exponentials, logarithms, and iteration.  They are
the right global setting for closure under broad classes of operations
\cite{edgar-beginners,edgar-composition,vdhoeven-book,adh-book}.  The smaller
class \eqref{eq:basic-pl-form} has a different advantage: each coefficient can
be manipulated by finite polynomial algebra, and the truncation order in
\(X^{-1}\) supplies a simple valuation.

\section{Four recurring sources}
Polynomial--logarithmic expansions arise in several closely related ways.

\subsection{Asymptotic functional inversion}
If
\[
 f(y)=y+\alpha\log y+\frac{p(\log y)}{y}+\cdots,
\]
then its inverse usually has the form
\[
 \rev{f}(x)=x-\alpha\log x
 +\sum_{n\ge1}\frac{p_n(\log x)}{x^n}.
\]
This includes Lambert \(W\), generalized Lambert equations, asymptotic
inverses used in prime-number estimates, and inverses of many exp--log
functions.  Symbolic algorithms for broader exp--log inverse problems were
developed by Salvy, Shackell, Richardson, and van der Hoeven
\cite{salvy-fast,salvy-shackell,rssv-exp-log}.

\subsection{Implicit equations}
Even when no named inverse function is available, an equation such as
\[
 y+\log y+\frac{\log y}{y}=x
\]
can be solved coefficient by coefficient.  The formal residual determines one
new polynomial at each order.  This triangular mechanism is one of the main
themes of this article.

\subsection{Iteration and normal forms}
Power--log series also occur in the formal classification of maps near fixed
points and in Dulac expansions associated with planar dynamical systems.
Composition and compositional inversion are then intrinsic rather than
auxiliary operations.  In this literature one often works near \(x=0^+\), with
monomials \(x^\alpha(\log(1/x))^m\), and sometimes must admit iterated
logarithms \cite{mardesic-powerlog,mardesic-dulac,peran-logarithmic}.

\subsection{Asymptotic solutions of differential equations}
Differential equations whose dominant balance is algebraic or logarithmic can
produce coefficients that are polynomial in \(\log X\).  Resonance is a common
mechanism: when the linear operator that normally determines a coefficient
loses invertibility on a particular exponent, a logarithm enters.  Although
our main focus is algebra and inversion, the differentiation rules developed
below are exactly those used in such calculations.

\section{Formal equality and analytic meaning}
A formal series is an algebraic object.  The statement
\[
 A\wideeq B
\]
will mean that the two formal coefficient arrays agree, not that two infinite
numerical sums converge.  We write
\[
 A=B+\FormalBigOAt{t^N}{t},\qquad t=X^{-1},
\]
when the coefficients of \(t^n\) agree for every \(n<N\).  This is a
\emph{formal} big-O notation.

An actual function \(f(X)\) has a polynomial--logarithmic asymptotic expansion
if its successive remainders are controlled in the corresponding ordered
scale.  One convenient formulation is the following.  Suppose
\(a_n\in\K[L]\) and let \(d_N=\max_{n\le N}\deg a_n\).  We say that
\[
 f(X)\sim\sum_{n\ge n_0}a_n(\log X)X^{-n}
\]
if for every \(N\ge n_0\) there exists an integer \(M_N\) such that
\begin{equation}\label{eq:analytic-remainder}
 f(X)-\sum_{n=n_0}^{N}a_n(\log X)X^{-n}
 =\BigOAt{X^{-N-1}(\log X)^{M_N}}{X\to\infty,\ X\in S}
\end{equation}
through the sector or ray under consideration.  Stronger definitions enumerate
all individual monomials \(X^{-n}L^m\), but
\cref{eq:analytic-remainder} is well suited to blockwise computation.

\begin{warningbox}
Formal correctness and analytic validity are separate questions.  A formal
inverse can exist even when the corresponding analytic branch is absent on a
given domain.  Conversely, an analytic inverse may have different formal
expansions in different sectors because logarithm branches change.  Our
algebraic operations are branch-neutral until a specific analytic
interpretation is imposed.
\end{warningbox}

\section{The two-variable illusion}
It is tempting to regard \(t=X^{-1}\) and \(L=\log X\) as independent formal
variables.  For addition and multiplication this causes no harm.  Under
differentiation or composition it is wrong:
\[
 \frac{\Differential L}{\Differential X}=\frac1X=t,
 \qquad
 \log G(X)\ne L\quad\text{unless }G(X)=X.
\]
The relation between the fast and slow variables is the source of the shifted
differential operators that will appear later.  In particular,
\begin{equation}\label{eq:derivative-preview}
 \frac{\Differential}{\Differential X}\left(a(L)X^{-n}\right)
 =\bigl(a'(L)-n a(L)\bigr)X^{-n-1}.
\end{equation}
The same coupling is responsible for the derivative terms in all universal
composition and reversal formulas.

\section{A roadmap of the algebra}
The main operations have different closure behavior.
\begin{center}
\begin{tabularx}{0.95\textwidth}{@{}>{\bfseries}l X X@{}}
\toprule
Operation & Typical input hypothesis & Natural output class \\
\midrule
Addition, multiplication & \(\K[L]((t))\) & \(\K[L]((t))\) \\
Differentiation & \(\K[L]((t))\) & \(t\K[L]((t))\) up to the leading power \\
Reciprocal & leading coefficient a nonzero constant & \(\K[L]((t))\) \\
General reciprocal & leading coefficient a nonzero polynomial & \(\K(L)((t))\) \\
Near-identity composition & inner series \(X+\LittleO(X)\) of controlled form & usually \(\K[L]((t))\) \\
General composition & inner leading term contains \(L^\beta\) & may require \(\log L\) \\
Near-identity reversal & \(F(X)=X+\LittleO(X)\) with polynomial--log blocks & \(\K[L]((t))\) \\
\bottomrule
\end{tabularx}
\end{center}
The word ``usually'' in the composition row matters.  Closure is proved for the
specific classes used in this article; unrestricted formal substitution can
fail even in much larger transseries settings unless support conditions are
imposed.

\begin{summarybox}
Polynomial--logarithmic transseries arise naturally when a leading algebraic
scale is corrected by a logarithm and then inverted or composed.  The basic
normal form is a Laurent series in \(t=X^{-1}\) with coefficients in
\(\K[L]\), where \(L=\log X\).  Formal truncation is t-adic, while analytic
meaning requires a separate remainder statement.  The variables \(t\) and
\(L\) are coupled under differentiation and composition.
\end{summarybox}

\chapter{Coordinates, rings, valuations, and normal forms}\label{ch:coordinates}

\section{The large-variable coordinate system}
Throughout the core chapters, \(X\) tends to infinity in a region where a
branch of \(\log X\) has been fixed.  Set
\begin{equation}\label{eq:qL-def}
 t=X^{-1},\qquad L=\log X.
\end{equation}
The coefficient field \(\K\) may be \(\RationalNumbers\), \(\RealNumbers\), or \(\ComplexNumbers\).  For symbolic
work it can also be a rational function field in external parameters.

\begin{definition}[Polynomial--logarithmic Laurent ring]\label{def:PL-ring}
The polynomial--logarithmic Laurent ring at infinity is
\begin{equation}\label{eq:PL-ring}
 \PolynomialLogLaurentRing{\K}=\K[L]((t))
 =\left\{\sum_{n\ge n_0}a_n(L)t^n:
 n_0\in\IntegerNumbers,\ a_n\in\K[L]\right\}.
\end{equation}
Only finitely many negative powers of \(t\) are allowed, but infinitely many
nonnegative powers may occur.
\end{definition}
Thus \(X-L+L/X\) is represented as \(t^{-1}-L+Lt\).  The notation
\(\K[L]((t))\) emphasizes that the outer formal-series variable is \(t\), while
\(L\) belongs to the coefficient ring.

\section{Valuation and truncation}
For nonzero
\[
 A=\sum_{n\ge n_0}a_n(L)t^n,
 \qquad a_{n_0}\ne0,
\]
define
\begin{equation}\label{eq:qvaluation}
 \val_t(A)=n_0,
 \qquad \val_t(0)=+\infty.
\end{equation}
The t-adic valuation satisfies
\[
 \val_t(AB)=\val_t(A)+\val_t(B),
 \qquad
 \val_t(A+B)\ge\min\{\val_t(A),\val_t(B)\}.
\]
The valuation-tail notation
\[
 A=B+\FormalBigOAt{t^N}{t}
\]
means \(\val_t(A-B)\ge N\).  The truncation operator is
\begin{equation}\label{eq:truncation}
 [A]_{<N}=\sum_{n_0\le n<N}a_n(L)t^n.
\end{equation}
Every algorithm below takes \(N\) as input and performs only finitely many
coefficient operations to compute \([A]_{<N}\).

\begin{proposition}[t-adic completeness]\label{prop:qcomplete}
Let \((A_j)\) be a sequence in \(\PolynomialLogLaurentRing{\K}\) such that for every \(N\), the
truncations \([A_j]_{<N}\) eventually stabilize.  Then there is a unique
\(A\in\PolynomialLogLaurentRing{\K}\) whose truncation at every order is that stable value.
\end{proposition}
\begin{proof}
For each integer \(n\), choose \(N>n\).  Eventual stabilization below outer
\(t\)-order \(N\)
fixes a unique coefficient \(a_n(L)\).  These coefficients are compatible as
\(N\) varies.  The resulting formal Laurent series has a lower bound on its
support because the sequence is already Laurent-bounded at some initial
precision.  Uniqueness follows coefficientwise.
\end{proof}
This elementary completeness is what makes Newton iteration a formal algorithm:
one does not sum numerical values; one stabilizes more and more coefficients.

\section{The internal logarithmic order}
Within a fixed t-block, powers of \(L\) are ordered by degree.  For monomials
\(t^nL^m=X^{-n}(\log X)^m\), the asymptotic order at \(+\infty\) is
lexicographic:
\begin{equation}\label{eq:lex-order}
 t^{n_1}L^{m_1}\succ t^{n_2}L^{m_2}
 \quad\Longleftrightarrow\quad
 \begin{cases}
 n_1<n_2,\quad\text{or}\\
 n_1=n_2\text{ and }m_1>m_2.
 \end{cases}
\end{equation}
Indeed, every power of \(L\) is smaller than every positive power of \(X\).
The block notation \(a_n(L)t^n\) deliberately groups finitely many adjacent
monomials.

Cancellation can lower the degree inside a block or remove the block entirely.
For example,
\[
 (L^2+L)t^3+(-L^2+2)t^3=(L+2)t^3.
\]
An implementation must normalize polynomial coefficients before deciding
whether a t-block vanishes.

\section{Why the polynomial ring is not a field}
The coefficient ring \(\K[L]\) has only the nonzero constants as units.  Hence
\(\PolynomialLogLaurentRing{\K}\) is a ring but not a field.

\begin{proposition}[{Units in \(\K[L]((t))\)}]\label{prop:units}
Let
\[
 A=t^r\bigl(a_0(L)+a_1(L)t+a_2(L)t^2+\cdots\bigr),
 \qquad a_0\ne0.
\]
Then \(A\) has a reciprocal in \(\PolynomialLogLaurentRing{\K}\) if and only if \(a_0(L)\) is a
nonzero constant.
\end{proposition}
\begin{proof}
The monomial \(t^r\) is invertible, so it suffices to consider \(r=0\).  If
\(AB=1\), the constant t-coefficient gives \(a_0(L)b_0(L)=1\) in
\(\K[L]\).  Thus \(a_0\) must be a unit, hence a nonzero constant.  Conversely,
if \(a_0\in\K^\times\), the triangular reciprocal recurrence in
\cref{ch:division} constructs all \(b_n\in\K[L]\).
\end{proof}

The natural field extension is
\begin{equation}\label{eq:rational-log-field}
 \RationalLogLaurentField{\K}=\K(L)((t)).
\end{equation}
Every nonzero element of \(\RationalLogLaurentField{\K}\) is invertible because \(\K(L)\) is a
field.  A smaller extension, \(\K[L,L^{-1}]((t))\), suffices when denominators
are monomials in \(L\), but not for a factor such as \(L+1\).

\begin{workedbox}
The series \(A=L+t\) has no reciprocal in \(\PolynomialLogLaurentRing{\K}\).  In \(\RationalLogLaurentField{\K}\),
\[
 \frac1{L+t}
 =\frac1L\frac1{1+t/L}
 =\frac1L-\frac{t}{L^2}+\frac{t^2}{L^3}-\cdots.
\]
The failure is not caused by the infinite t-expansion; it is already present in
the leading coefficient \(L\).
\end{workedbox}

\section{Affine and monomial prefactors}
Many useful expansions do not begin at t-order zero.  Write a nonzero series as
\begin{equation}\label{eq:prefactor-form}
 A=cX^\rho L^\beta(1+U),
 \qquad U\prec1,
\end{equation}
whenever such a factorization is available.  In the strict ring
\(\PolynomialLogLaurentRing{\K}\), \(\rho\in\IntegerNumbers\), \(\beta\in\NaturalNumbers\), and the factor \(L^\beta\) may
prevent invertibility.  In generalized power--log classes, \(\rho,\beta\) may be
real or complex.

The factorization \eqref{eq:prefactor-form} separates three tasks:
\begin{enumerate}
\item invert or raise the monomial \(cX^\rho L^\beta\) explicitly;
\item apply an ordinary binomial, logarithm, or exponential series to \(1+U\);
\item truncate t-adically, using that \(U\) has positive valuation.
\end{enumerate}
This is the polynomial--logarithmic analogue of the canonical multiplicative
decomposition used in full transseries.

\section{Degree slopes and finite bookkeeping}
The Lambert polynomials satisfy \(\deg \LambertPolynomial{n}=n\).  More generally, many
calculations preserve a linear degree bound
\begin{equation}\label{eq:slope-bound}
 \deg a_n\le \sigma n+C.
\end{equation}
Such a bound is not needed for formal multiplication, but it translates a
t-adic remainder into a useful analytic estimate and limits coefficient size.

\begin{proposition}[Slope under multiplication]\label{prop:slope-product}
Suppose
\[
 \deg a_n\le A+\alpha n,
 \qquad
 \deg b_n\le B+\beta n
\]
for \(n\ge0\).  If \(C=AB=\sum c_nt^n\), then
\[
 \deg c_n\le A+B+\max(\alpha,\beta)n.
\]
\end{proposition}
\begin{proof}
Each summand \(a_i b_{n-i}\) in \(c_n\) has degree at most
\[
 A+B+\alpha i+\beta(n-i)
 \le A+B+\max(\alpha,\beta)n.
\]
Taking a finite sum cannot increase this bound.
\end{proof}

\section{Generalized exponent supports}
For some problems integer powers of \(t\) are too restrictive.  Let
\(S\subset\RealNumbers\) be a well-ordered set with the property that every real bound
contains only finitely many elements of \(S\), or more generally assume the
usual Hahn-series support conditions.  Then one may consider
\begin{equation}\label{eq:generalized-powerlog}
 A(X)=\sum_{\lambda\in S}a_\lambda(L)X^{-\lambda}.
\end{equation}
Multiplication is defined by
\[
 (AB)_\nu=\sum_{\lambda+\mu=\nu}a_\lambda b_\mu,
\]
provided each coefficient receives only finitely many contributions and the
sum support remains well ordered.  Finite-rank grids
\[
 S\subset\lambda_0+\NaturalNumbers\omega_1+\cdots+\NaturalNumbers\omega_r,
 \qquad \omega_j>0,
\]
are particularly convenient.  The integer t-adic theory developed below then
extends with multi-indices or an ordered grid.

\begin{warningbox}
A formal sum over arbitrary real exponents is not automatically meaningful.
Without a support condition, a coefficient of a product may receive infinitely
many contributions, or the support may contain an infinite descending chain in
the asymptotic order.  ``Write all powers that seem useful'' is not a
definition of a transseries algebra.
\end{warningbox}

\section{Changing to a small variable}
Set \(x=X^{-1}\to0^+\).  Then
\[
 L=\log X=\log(1/x)=-\log x.
\]
The core expansion becomes
\begin{equation}\label{eq:small-variable-form}
 \sum_{n\ge n_0}p_n\!\left(\log\frac1x\right)x^n.
\end{equation}
This is the convention common in power--log and Dulac series.  Both coordinate
systems encode the same algebra, but signs inside logarithmic polynomials
change if one uses \(\log x\) rather than \(\log(1/x)\).  We postpone a fuller
discussion to \cref{ch:dulac}.

\begin{summarybox}
The core ring is \(\PolynomialLogLaurentRing{\K}=\K[L]((t))\) with \(t=X^{-1}\) and
\(L=\log X\).  Its t-adic valuation makes truncation and convergence purely
coefficientwise.  It is not a field: a series is a unit only when its first
nonzero t-coefficient is a nonzero constant.  General division naturally leads
to \(\K(L)((t))\).  Generalized real exponents are possible, but only with
support conditions that preserve local finiteness.
\end{summarybox}

\chapter{Multiplication and differential calculus}\label{ch:multiplication}

\section{Addition and normalization}
Let
\[
 A=\sum_{n\ge r}a_n(L)t^n,
 \qquad
 B=\sum_{n\ge s}b_n(L)t^n.
\]
After inserting zero coefficients where necessary,
\[
 A+B=\sum_{n\ge\min(r,s)}\bigl(a_n(L)+b_n(L)\bigr)t^n.
\]
One should expand, collect, or otherwise normalize each polynomial coefficient
before testing it for zero.  This matters because cancellation may raise the
t-valuation.

\begin{example}
If
\[
 A=t^{-1}-L+(L^2-1)t,
 \qquad
 B=-t^{-1}+L+(2-L^2)t+3t^2,
\]
then
\[
 A+B=t+3t^2,
\]
so two leading t-blocks and the quadratic logarithmic terms all cancel.
\end{example}

\section{Cauchy multiplication}
Multiplication uses the same convolution as ordinary power series, but the
coefficient arithmetic takes place in \(\K[L]\).

\begin{proposition}[Product formula]\label{prop:product-formula}
For
\[
 A=\sum_{i\ge r}a_i(L)t^i,
 \qquad
 B=\sum_{j\ge s}b_j(L)t^j,
\]
the product is
\begin{equation}\label{eq:product-convolution}
 AB=\sum_{n\ge r+s}c_n(L)t^n,
 \qquad
 c_n(L)=\sum_{i+j=n}a_i(L)b_j(L).
\end{equation}
For each fixed \(n\), the sum is finite.
\end{proposition}
\begin{proof}
Because \(i\ge r\) and \(j\ge s\), the equation \(i+j=n\) restricts
\(i\) to the finite interval \(r\le i\le n-s\).  Distributivity then gives
\cref{eq:product-convolution}.  The result is again a Laurent series with
polynomial coefficients.
\end{proof}

To compute modulo \(t^N\), discard every pair with \(i+j\ge N\).  No higher
coefficient can influence a lower one.

\begin{algorithm}[Truncated multiplication]\label{alg:mul}
Given sparse maps \(i\mapsto a_i(L)\) and \(j\mapsto b_j(L)\), and a cutoff
\(N\):
\begin{enumerate}
\item initialize \(c_n=0\) for all required \(n<N\);
\item for each stored \(i\) and \(j\) with \(i+j<N\), add \(a_i b_j\) to
      \(c_{i+j}\);
\item normalize the polynomials and delete zero blocks.
\end{enumerate}
The output is \([AB]_{<N}\).
\end{algorithm}

\begin{workedbox}
Take
\[
 A=1+(L-1)t+(L^2+2)t^2+\FormalBigOAt{t^3}{t},
\]
\[
 B=1-(L+2)t+(3L-1)t^2+\FormalBigOAt{t^3}{t}.
\]
The coefficient of \(t\) is
\[
 (L-1)-(L+2)=-3.
\]
The coefficient of \(t^2\) is
\[
 (L^2+2)+(3L-1)-(L-1)(L+2)=2L+3.
\]
Therefore
\[
 AB=1-3t+(2L+3)t^2+\FormalBigOAt{t^3}{t}.
\]
The leading logarithmic squares cancel even though neither factor has constant
coefficients.
\end{workedbox}

\section{Powers of a normalized series}
If \(U\in t\K[L][[t]]\), then \(U^m\in t^m\K[L][[t]]\).  Consequently, for a
fixed cutoff \(t^N\), only \(m<N\) can contribute to a series in powers of
\(U\).  This simple observation makes formal binomial, exponential, and
logarithmic series finite at every requested precision.

For an integer \(r\ge0\), \((1+U)^r\) is an ordinary finite polynomial.  For
arbitrary \(\alpha\in\K\), define
\begin{equation}\label{eq:formal-binomial}
 (1+U)^\alpha
 =\sum_{m\ge0}\binom{\alpha}{m}U^m,
 \qquad
 \binom{\alpha}{m}=
 \frac{\alpha(\alpha-1)\cdots(\alpha-m+1)}{m!}.
\end{equation}
Modulo \(t^N\), the sum stops at \(m=N-1\).

\begin{example}[Square root]
Let
\[
 A=1+2Lq+(L^2+1)t^2+\FormalBigOAt{t^3}{t}.
\]
Using \((1+U)^{1/2}=1+U/2-U^2/8+\cdots\),
\[
 A^{1/2}=1+Lq+\frac12t^2+\FormalBigOAt{t^3}{t},
\]
because the \(L^2t^2\) contributions cancel.
\end{example}

\section{Differentiation and the shifted derivative}
Let \(\Dop=\Differential/\Differential L\).  Since \(L=\log X\),
\[
 \frac{\Differential L}{\Differential X}=X^{-1}.
\]
For a single block,
\begin{equation}\label{eq:block-derivative}
 \frac{\Differential}{\Differential X}\bigl(a_n(L)X^{-n}\bigr)
 =\bigl(\Dop a_n-na_n\bigr)X^{-n-1}.
\end{equation}
It is convenient to define
\begin{equation}\label{eq:Delta-n}
 \Delta_n=\Dop-n.
\end{equation}
Then
\begin{equation}\label{eq:derivative-series}
 A'(X)=\sum_n \Delta_n a_n(L)\,X^{-n-1}.
\end{equation}
The index shift is as important as the coefficient derivative.

\begin{proof}[Derivation of \cref{eq:block-derivative}]
The product and chain rules give
\[
 \frac{\Differential}{\Differential X}\left(a_n(\log X)X^{-n}\right)
 =a_n'(L)X^{-1}X^{-n}-na_n(L)X^{-n-1}.
\]
Collecting the common power proves the formula.
\end{proof}

\begin{example}
For
\[
 A=X-L+\frac{L^2-3L}{X},
\]
we obtain
\[
 A'=1-\frac1X
 +\frac{(2L-3)-(L^2-3L)}{X^2}
 =1-\frac1X+\frac{-L^2+5L-3}{X^2}.
\]
\end{example}

\begin{warningbox}
If a computer algebra system treats \(L\) as a symbol independent of \(X\), it
will incorrectly differentiate \(a(L)X^{-n}\) as merely
\(-na(L)X^{-n-1}\).  Either substitute \(L=\log X\) before differentiating or
implement the shifted operator \(\Delta_n\) explicitly.
\end{warningbox}

\section{Repeated derivatives}
One derivative changes both the power index and the shifted operator.  Applying
\cref{eq:block-derivative} repeatedly gives
\begin{equation}\label{eq:repeated-derivative}
 \frac{\Differential^r}{\Differential X^r}\bigl(a_n(L)X^{-n}\bigr)
 =\left(\prod_{j=0}^{r-1}(\Dop-n-j)\right)a_n(L)\,X^{-n-r}.
\end{equation}
The factors commute because each is a polynomial in \(\Dop\).  For example,
\[
 \frac{\Differential^2}{\Differential X^2}\bigl(a_n(L)X^{-n}\bigr)
 =\bigl(\Dop-n-1\bigr)\bigl(\Dop-n\bigr)a_n(L)X^{-n-2}.
\]
This formula is useful in Taylor expansions during composition.

\section{Formal antiderivatives}
Although integration is not a central operation here, it clarifies why
polynomial coefficients are stable under differential calculus.  To integrate
\(p(L)X^{-n}\), seek \(Q(L)X^{-n+1}\).  Then
\[
 \frac{\Differential}{\Differential X}\left(Q(L)X^{-n+1}\right)
 =\bigl(\Dop-(n-1)\bigr)Q(L)X^{-n}.
\]
For \(n\ne1\), the operator \(\Dop-(n-1)\) is invertible on \(\K[L]\), with
finite inverse
\begin{equation}\label{eq:operator-inverse-polynomials}
 \bigl(\Dop-a\bigr)^{-1}p
 =-\sum_{j=0}^{\deg p}\frac{p^{(j)}}{a^{j+1}},
 \qquad a\ne0.
\end{equation}
Indeed, the sum telescopes after applying \(\Dop-a\).  When \(n=1\),
\[
 \int p(L)X^{-1}\Differential X=\int p(L)\Differential L,
\]
which is again polynomial in \(L\).  Thus termwise antiderivatives remain
polynomial--logarithmic, up to an arbitrary constant.

\section{Differential equations for coefficient blocks}
Suppose a formal equation produces
\[
 (\Dop-n)p_n(L)=r_n(L)
\]
for a known polynomial \(r_n\).  If \(n\ne0\), \cref{eq:operator-inverse-polynomials}
solves uniquely for \(p_n\).  If \(n=0\), the equation is
\(p_0'=r_0\), which determines \(p_0\) up to a constant.  Such resonant
constants are often fixed by a normalization condition or by matching another
asymptotic region.

This blockwise inversion is the simplest instance of a general principle:
formal asymptotic equations become triangular because the dominant monomial
separates successive t-orders.

\section{Products of many factors and Bell polynomials}
Repeated multiplication can be organized by partitions.  Let
\[
 U(t)=\sum_{j\ge1}u_jt^j.
\]
The coefficient of \(t^n\) in \(U^r\) is
\begin{equation}\label{eq:ordinary-compositions-coeff}
 [t^n]U^r
 =\sum_{j_1+\cdots+j_r=n}u_{j_1}\cdots u_{j_r},
\end{equation}
where all \(j_i\ge1\).  Equivalently, in terms of the partial exponential Bell
polynomial \(\ExponentialPartialBellPolynomial{n}{r}\),
\begin{equation}\label{eq:bell-coeff}
 [t^n]U^r
 =\frac{r!}{n!}\,
 \ExponentialPartialBellPolynomial{n}{r}\bigl(1!u_1,2!u_2,\ldots,(n-r+1)!u_{n-r+1}\bigr).
\end{equation}
We will use this identity to expand \(\log(1+U)\), exponentials, and composition
systematically.

\section{Complexity at fixed truncation}
With dense t-coefficients, schoolbook multiplication through order \(N\) uses
\(\BigO(N^2)\) polynomial products.  If the logarithmic degree is \(\BigO(N)\), naive
polynomial multiplication adds another quadratic factor in the worst case.
For the moderate orders typical in hand asymptotics, this is acceptable and
transparent.  At high orders one may use fast convolution in \(t\), fast
polynomial multiplication in \(L\), or relaxed power-series algorithms.

The valuation is more important than asymptotic complexity: a correct
implementation should never compute terms that cannot affect the requested
precision.

\begin{summarybox}
Multiplication is Cauchy convolution in the t-index and ordinary polynomial
multiplication in \(L\).  Positive t-valuation makes binomial and other
power-series substitutions finite at fixed precision.  Differentiation is
controlled by the shifted operator \(\Delta_n=\Dop-n\), because
\(L=\log X\) is coupled to \(X\).  Bell polynomials package the coefficient
combinatorics of repeated products.
\end{summarybox}

\chapter{Reciprocals, division, logarithms, and exponentials}\label{ch:division}

\section{Multiplicative inversion is a triangular problem}
Let
\begin{equation}\label{eq:A-recip-general}
 A=a_0+a_1t+a_2t^2+\cdots,
 \qquad a_0\in\K^\times,
\end{equation}
and seek
\[
 \recip{A}=B=b_0+b_1t+b_2t^2+\cdots
\]
such that \(AB=1\).  Comparing t-coefficients gives
\[
 a_0b_0=1,
\]
and, for \(n\ge1\),
\[
 a_0b_n+\sum_{k=1}^n a_kb_{n-k}=0.
\]
Therefore
\begin{equation}\label{eq:reciprocal-recurrence}
 b_0=a_0^{-1},
 \qquad
 b_n=-a_0^{-1}\sum_{k=1}^n a_kb_{n-k}.
\end{equation}
Every \(b_n\) depends only on earlier \(b_j\), so the calculation is
triangular.

\begin{proposition}[Existence and uniqueness of the reciprocal]\label{prop:reciprocal}
If \(a_0\in\K^\times\), the recurrence \eqref{eq:reciprocal-recurrence}
defines the unique \(B\in\K[L][[t]]\) satisfying \(AB=1\).
\end{proposition}
\begin{proof}
Existence follows by induction: once \(b_0,\ldots,b_{n-1}\) are known, the
coefficient equation at t-order \(n\) has the unique solution shown in
\cref{eq:reciprocal-recurrence}.  The resulting product has constant
coefficient one and every positive t-coefficient zero.  If \(B\) and
\(\widetilde B\) are two reciprocals, then
\(B=B(A\widetilde B)=(BA)\widetilde B=\widetilde B\).
\end{proof}

A Laurent prefactor causes no difficulty.  If
\[
 A=t^r\widetilde A,
 \qquad \val_t(\widetilde A)=0,
\]
then
\[
 \recip{A}=t^{-r}\recip{\widetilde A}.
\]

\begin{workedbox}
Let
\[
 A=1+Lq+(L^2+1)t^2+\FormalBigOAt{t^4}{t},
\]
with the t-cubic coefficient taken to be zero.  The recurrence gives
\[
 b_0=1,
 \qquad b_1=-L,
\]
\[
 b_2=-\bigl(Lb_1+(L^2+1)b_0\bigr)=-1,
\]
and
\[
 b_3=-\bigl(Lb_2+(L^2+1)b_1\bigr)=L^3+2L.
\]
Hence
\[
 \frac1A=1-Lq-t^2+(L^3+2L)t^3+\FormalBigOAt{t^4}{t}.
\]
Direct multiplication is the quickest verification.
\end{workedbox}

\section{Geometric normalization}
When \(A=a_0(1+U)\) with \(U\in t\K[L][[t]]\), the recurrence is equivalent to
\begin{equation}\label{eq:geometric-recip}
 \frac1A=a_0^{-1}\sum_{m\ge0}(-U)^m.
\end{equation}
At precision \(t^N\), only \(m<N\) contributes.  The recurrence is usually
better for sequential coefficient generation; the geometric formula is often
better for hand calculations and proofs.

\section{Direct quotient recurrence}
To divide \(C\) by \(A\), one may compute \(C\recip{A}\), but a direct
recurrence avoids storing the entire reciprocal.  Suppose
\[
 A=\sum_{n\ge0}a_nt^n,
 \qquad
 C=\sum_{n\ge0}c_nt^n,
 \qquad
 Q=\frac CA=\sum_{n\ge0}d_nt^n,
\]
with \(a_0\in\K^\times\).  From \(AQ=C\),
\begin{equation}\label{eq:quotient-recurrence}
 d_0=\frac{c_0}{a_0},
 \qquad
 d_n=\frac1{a_0}
 \left(c_n-\sum_{k=1}^n a_kd_{n-k}\right).
\end{equation}
This works without change when the coefficients lie in \(\K(L)\) rather than
\(\K[L]\).

\begin{example}
With
\[
 C=1+(2L+1)t+L^2t^2,
 \qquad
 A=1+(L-1)t+t^2,
\]
we find
\[
 d_0=1,
 \quad
 d_1=(2L+1)-(L-1)=L+2,
\]
\[
 d_2=L^2-(L-1)(L+2)-1=-L+1.
\]
Thus
\[
 \frac CA=1+(L+2)t+(1-L)t^2+\FormalBigOAt{t^3}{t}.
\]
\end{example}

\section{What changes when the leading coefficient depends on \texorpdfstring{\(L\)}{L}}
If \(a_0(L)\) is a nonconstant polynomial, the same recurrence still works in
\(\K(L)[[t]]\):
\[
 b_0=\frac1{a_0(L)},
 \qquad
 b_n=-\frac1{a_0(L)}\sum_{k=1}^na_k(L)b_{n-k}(L).
\]
The output generally has rational, not polynomial, logarithmic coefficients.
For instance,
\[
 \frac1{L+t}=L^{-1}-L^{-2}t+L^{-3}t^2-\cdots.
\]

There are three reasonable responses in symbolic work:
\begin{enumerate}
\item enlarge the coefficient field to \(\K(L)\);
\item factor the logarithmic monomial and work in a class allowing negative
powers of \(L\);
\item keep the expression factored rather than expanding it.
\end{enumerate}
The right choice depends on the intended subsequent operations.  Expanding into
\(L^{-1}\) can be asymptotically useful, but it changes the monomial hierarchy.

\section{Newton inversion}
The recurrence computes one coefficient at a time.  Newton's method can double
the number of correct coefficients at each step.  Given an approximation \(B\)
to \(A^{-1}\), set
\begin{equation}\label{eq:newton-recip}
 B_{\mathrm{new}}=B(2-AB).
\end{equation}

\begin{proposition}[Precision doubling for reciprocals]\label{prop:newton-recip}
If
\[
 AB=1-E,
 \qquad \val_t(E)\ge N,
\]
then
\[
 AB_{\mathrm{new}}=1-E^2,
 \qquad \val_t(1-AB_{\mathrm{new}})\ge2N.
\]
\end{proposition}
\begin{proof}
Substitute \(AB=1-E\):
\[
 AB_{\mathrm{new}}
 =(AB)(2-AB)=(1-E)(1+E)=1-E^2.
\]
The valuation is multiplicative, so \(\val_t(E^2)\ge2N\).
\end{proof}

A practical implementation begins with \(B_0=a_0^{-1}\) and truncates after
each Newton step to the desired doubled precision.  Fast multiplication then
makes Newton inversion asymptotically preferable at high order.

\section{Formal logarithms}
If \(U\in t\K[L][[t]]\), define
\begin{equation}\label{eq:formal-log1p}
 \log(1+U)=\sum_{m\ge1}\frac{(-1)^{m+1}}mU^m.
\end{equation}
This lies in \(t\K[L][[t]]\).  For
\[
 A=cX^\rho(1+U),
 \qquad c\ne0,
\]
a chosen logarithm branch gives
\begin{equation}\label{eq:formal-log-general}
 \log A=\log c+\rho L+\log(1+U).
\end{equation}
If a factor \(L^\beta\) is present, an additional term \(\beta\log L\) appears.
That term lies outside the one-log-level coefficient ring and is the first clear
signal that iterated logarithms may be required.

\begin{workedbox}
Let
\[
 A=1+Lq+(2-L)t^2+\FormalBigOAt{t^3}{t}.
\]
Then \(U=Lt+(2-L)t^2\), and
\[
 \log A=U-\frac{U^2}{2}+\FormalBigOAt{t^3}{t}
 =Lq+\left(2-L-\frac{L^2}{2}\right)t^2+\FormalBigOAt{t^3}{t}.
\]
\end{workedbox}

\section{Formal exponentials}
If \(V\in t\K[L][[t]]\), define
\begin{equation}\label{eq:formal-exp-small}
 \exp V=\sum_{m\ge0}\frac{V^m}{m!}.
\end{equation}
It lies in \(1+t\K[L][[t]]\).  At fixed t-precision, only finitely many powers
of \(V\) contribute.  The formal identities
\[
 \exp(\log(1+U))=1+U,
 \qquad
 \log(\exp V)=V
\]
follow from the corresponding identities in the complete t-adic algebra.

An arbitrary polynomial in \(L\) cannot be exponentiated while staying inside
\(\PolynomialLogLaurentRing{\K}\).  For example,
\[
 \exp(L)=X=t^{-1}
\]
is allowed after recognizing the monomial, but
\[
 \exp(L^2)=X^{\log X}
\]
is a genuinely new transmonomial.  Full transseries handle such exponentials
by enlarging the monomial group.

\section{Logarithmic differentiation}
For a unit \(A\),
\begin{equation}\label{eq:logderivative}
 \frac{A'}A=(\log A)'.
\end{equation}
If \(A=cX^\rho(1+U)\), then
\[
 \frac{A'}A=\frac\rho X+\frac{U'}{1+U}.
\]
The logarithmic derivative often has higher t-valuation than \(A'\) itself and
is a natural measure of how rapidly a transmonomial changes.  In the general
theory of composition, comparisons between a perturbation and logarithmic
derivatives of monomials distinguish shallow, moderate, and deep substitutions
\cite{edgar-composition}.

\section{General powers}
For a normalized unit \(A=a_0(1+U)\) and \(\alpha\in\K\), define
\begin{equation}\label{eq:general-power}
 A^\alpha=a_0^\alpha(1+U)^\alpha,
\end{equation}
with compatible choices of \(a_0^\alpha\) and logarithm branch.  Equivalently,
\(A^\alpha=\exp(\alpha\log A)\).  The binomial formula
\cref{eq:formal-binomial} gives a direct coefficient algorithm.

For a general prefactor \(cX^\rho L^\beta(1+U)\),
\[
 A^\alpha=c^\alpha X^{\alpha\rho}L^{\alpha\beta}(1+U)^\alpha.
\]
Closure may require generalized powers of \(X\) or \(L\), but the small-unit
part remains t-adic.

\section{Residual-based verification}
Every reciprocal or quotient should be checked by multiplication.  If
\(B_N\) is intended as the reciprocal of \(A\) through order \(t^{N-1}\), verify
\begin{equation}\label{eq:recip-residual-test}
 AB_N-1=\FormalBigOAt{t^N}{t}.
\end{equation}
For \(Q_N=C/A\), verify
\[
 AQ_N-C=\FormalBigOAt{t^N}{t}.
\]
These tests detect off-by-one truncation errors, failures to normalize
polynomials, and accidental use of \(L\) as an independent variable.

\begin{summarybox}
Reciprocal and quotient coefficients satisfy triangular convolution
recurrences.  Polynomial coefficients are preserved only when the leading
coefficient is a nonzero constant; otherwise the natural coefficient field is
\(\K(L)\).  Newton's update \(B\mapsto B(2-AB)\) squares the residual and
doubles t-adic precision.  Logarithms and exponentials are defined on small
t-adic perturbations, while prefactors determine whether new monomials such as
\(\log L\) or \(\exp(L^2)\) must be admitted.
\end{summarybox}

\chapter{Composition}\label{ch:composition}

\section{Why composition is the delicate operation}
Addition and multiplication keep \(L\) fixed.  Composition does not.  If an
outer series is written in a variable \(Y\), then substitution \(Y=G(X)\)
changes both of its building blocks:
\[
 Y^{-1}\longmapsto G(X)^{-1},
 \qquad
 \log Y\longmapsto\log G(X).
\]
The second change is responsible for derivatives of coefficient polynomials;
the first is responsible for shifted derivatives.  Moreover, an unrestricted
substitution can introduce \(\log\log X\), new real powers, or even ill-based
supports.  We therefore begin with a composition class whose closure can be
proved directly.

\section{Regular inner arguments}
Let the outer series be
\begin{equation}\label{eq:outer-composition}
 F(Y)=\sum_{n\ge n_0}f_n(\log Y)Y^{-n},
 \qquad f_n\in\K[L].
\end{equation}
Assume that the inner argument has the form
\begin{equation}\label{eq:regular-inner}
 G(X)=\lambda X(1+u(X)),
 \qquad \lambda\in\K^\times,
 \qquad u\in t\K[L][[t]].
\end{equation}
For an analytic interpretation, choose compatible logarithm branches and a
region in which \(G(X)\) stays in the domain of the outer branch.  Formally set
\begin{equation}\label{eq:delta-def}
 \delta=\log(1+u)\in t\K[L][[t]].
\end{equation}
Then
\begin{equation}\label{eq:logG-and-power}
 \log G=L+\log\lambda+\delta,
 \qquad
 G^{-n}=\lambda^{-n}X^{-n}\EulerE^{-n\delta}.
\end{equation}

\section{The normal-ordered composition formula}
Let \(\Dop=\Differential/\Differential L\).  Because \(\delta\) itself depends on \(L\), the
operators ``multiply by \(\delta\)'' and \(\Dop\) do not commute.  We therefore
introduce an explicit normal-ordering convention.

\begin{definition}[Normal-ordered shifted exponential]\label{def:normal-ordered}
For a polynomial \(f(L)\), a formal small series \(\delta\), and an integer
\(n\), define
\begin{equation}\label{eq:normal-ordered-def}
 \mathopen{:}\EulerE^{\delta(\Dop-n)}\mathclose{:}\,f
 :=\sum_{r\ge0}\frac{\delta^r}{r!}(\Dop-n)^rf.
\end{equation}
The colons mean that every power of \(\delta\) is placed to the left and is not
differentiated by \(\Dop\).
\end{definition}
Even though \((\Dop-n)^rf\) need not vanish for large \(r\), the series is
t-adically meaningful because \(\val_t(\delta^r)\ge r\).

\begin{theorem}[Regular composition formula]\label{thm:regular-composition}
Under \cref{eq:outer-composition,eq:regular-inner},
\begin{equation}\label{eq:regular-composition-formula}
 (F\compose G)(X)
 =\sum_{n\ge n_0}\lambda^{-n}X^{-n}
 \mathopen{:}\EulerE^{\delta(\Dop-n)}\mathclose{:}\,
 f_n(L+\log\lambda).
\end{equation}
At every finite t-precision, the right side requires only finitely many terms.
\end{theorem}
\begin{proof}
For one outer block,
\[
 f_n(\log G)G^{-n}
 =\lambda^{-n}X^{-n}
 f_n(L+\log\lambda+\delta)\EulerE^{-n\delta}.
\]
Taylor-expand the polynomial coefficient:
\[
 f_n(L+\log\lambda+\delta)
 =\sum_{j\ge0}\frac{\delta^j}{j!}\Dop^j
 f_n(L+\log\lambda).
\]
Multiplying by
\(\EulerE^{-n\delta}=\sum_{k\ge0}(-n)^k\delta^k/k!\) and collecting the total
power \(r=j+k\) gives
\[
 \sum_{r\ge0}\frac{\delta^r}{r!}(\Dop-n)^r
 f_n(L+\log\lambda),
\]
by the binomial theorem.  This is precisely the normal-ordered expression in
\cref{eq:regular-composition-formula}.  To compute modulo \(t^N\), only outer
indices \(n<N\) and powers \(r<N-n\) can contribute, up to the finite negative
part of the Laurent series.
\end{proof}

\begin{warningbox}
Without the normal-ordering colons, the expression
\(\exp(\delta(\Dop-n))\) is ambiguous and generally wrong.  If interpreted as
an ordinary differential-operator exponential, repeated application of
\(\delta\Dop\) differentiates \(\delta\).  Taylor substitution requires
\(\Dop\) to act only on the outer coefficient polynomial.
\end{warningbox}

\section{Coefficient extraction with Bell polynomials}
Write
\[
 \delta(t)=\sum_{j\ge1}\delta_j(L)t^j.
\]
By \cref{eq:bell-coeff},
\[
 [t^m]\delta^r
 =\frac{r!}{m!}\ExponentialPartialBellPolynomial{m}{r}
 \bigl(1!\delta_1,2!\delta_2,\ldots\bigr).
\]
Consequently, the coefficient of t-order \(N\) in
\cref{eq:regular-composition-formula} is
\begin{equation}\label{eq:composition-bell}
\begin{split}
 [t^N](F\compose G)
 ={}&\sum_{n=n_0}^{N}\lambda^{-n}
 \frac1{(N-n)!}
 \sum_{r=0}^{N-n}
 \ExponentialPartialBellPolynomial{N-n}{r}
 \bigl(1!\delta_1,2!\delta_2,\ldots\bigr)\\
 &\hspace{8em}\times
 (\Dop-n)^r f_n(L+\log\lambda),
\end{split}
\end{equation}
with \(\ExponentialPartialBellPolynomial{0}{0}=1\) and \(\ExponentialPartialBellPolynomial{m}{0}=0\) for \(m>0\).  This is a closed,
finite formula for every coefficient.

\begin{algorithm}[Regular composition to t-order \(N\)]\label{alg:regular-composition}
\leavevmode
\begin{enumerate}
\item Normalize the inner series as \(G=\lambda X(1+u)\).
\item Compute \(\delta=[\log(1+u)]_{<N}\).
\item For every outer block \(f_n(\log Y)Y^{-n}\) that can contribute, compute
      the shifted derivatives \((\Dop-n)^rf_n(L+\log\lambda)\).
\item Multiply them by \(\delta^r/r!\), the factor
      \(\lambda^{-n}t^n\), and truncate after every multiplication.
\item Sum and normalize the coefficient polynomials.
\end{enumerate}
Bell polynomials may replace repeated powers of \(\delta\) when many
coefficients are required.
\end{algorithm}

\section{A first worked composition}
Take
\[
 F(Y)=\log Y+\frac{(\log Y)^2}{Y},
 \qquad
 G(X)=2X\left(1+\frac{L}{X}\right).
\]
Here \(\lambda=2\), \(u=Lt\), and
\[
 \delta=\log(1+Lq)=Lq-\frac{L^2}{2}t^2+\FormalBigOAt{t^3}{t}.
\]
Directly,
\[
 \log G=L+\log2+\delta.
\]
Also,
\[
 G^{-1}=\frac{t}{2}\EulerE^{-\delta}
 =\frac t2\left(1-Lq+\frac{3L^2}{2}t^2+\FormalBigOAt{t^3}{t}\right).
\]
Therefore, through t-order two,
\begin{align*}
 F(G)
={}&L+\log2+Lt-\frac{L^2}{2}t^2\\
&+\frac t2\bigl(L+\log2+\delta\bigr)^2\EulerE^{-\delta}
 +\FormalBigOAt{t^3}{t}\\
={}&L+\log2
 +\left[L+\frac12(L+\log2)^2\right]t\\
&+\left[-\frac{L^2}{2}
 +L(L+\log2)-\frac L2(L+\log2)^2\right]t^2
 +\FormalBigOAt{t^3}{t}.
\end{align*}
The shifted-operator formula produces the same result without separately
expanding the logarithmic coefficient and the reciprocal power.

\section{Near-identity inner arguments}
The most important case for series reversal is
\begin{equation}\label{eq:near-identity-inner}
 G(X)=X+A(L)+\frac{B(L)}X+\frac{C(L)}{X^2}
      +\frac{D(L)}{X^3}+\FormalBigOAt{t^4}{t}.
\end{equation}
Factoring out \(X\) gives
\[
 G=X\bigl(1+At+Bt^2+Ct^3+Dt^4+\cdots\bigr).
\]
Thus
\begin{equation}\label{eq:delta-near-identity}
\begin{split}
 \delta=\log(G/X)
 ={}&At+\left(B-\frac{A^2}{2}\right)t^2
 +\left(C-AB+\frac{A^3}{3}\right)t^3\\
 &+\left(D-AC-\frac{B^2}{2}+A^2B-\frac{A^4}{4}\right)t^4
 +\FormalBigOAt{t^5}{t}.
\end{split}
\end{equation}
Every capital letter here is a polynomial in \(L\), and the products are
ordinary products at the same value of \(L\).

\section{Universal near-identity composition coefficients}
Let the outer series be
\begin{equation}\label{eq:outer-near-identity}
 F(Y)=Y+a(\log Y)+\frac{b(\log Y)}Y
 +\frac{c(\log Y)}{Y^2}+\frac{d(\log Y)}{Y^3}
 +\FormalBigOAt{Y^{-4}}{Y^{-1}},
\end{equation}
and let \(G\) be as in \cref{eq:near-identity-inner}.  Write
\[
 F(G(X))=X+h_0(L)+\frac{h_1(L)}X+\frac{h_2(L)}{X^2}
 +\frac{h_3(L)}{X^3}+\FormalBigOAt{t^4}{t}.
\]
Primes below mean derivatives with respect to \(L\), and all uncapitalized
functions are evaluated at \(L\).

\begin{theorem}[Composition through fourth block]\label{thm:near-composition-coeffs}
The coefficients are
\begin{align}
 h_0={}&A+a,\label{eq:h0}\\
 h_1={}&Aa'+B+b,\label{eq:h1}\\
 h_2={}&-\frac12A^2a'+\frac12A^2a''-Ab+Ab'
       +Ba'+C+c,\label{eq:h2}\\
 h_3={}&\frac13A^3a'-\frac12A^3a''+\frac16A^3a'''
       +A^2b-\frac32A^2b'+\frac12A^2b''\notag\\
 &{}-ABa'+ABa''-2Ac+Ac'-Bb+Bb'+Ca'+D+d.
 \label{eq:h3}
\end{align}
\end{theorem}
\begin{proof}
Use \(G=X(1+u)\) with
\(u=At+Bt^2+Ct^3+Dt^4+\cdots\).  Expand
\(\delta=\log(1+u)\) by \cref{eq:delta-near-identity}, then Taylor-expand
\(a(L+\delta),b(L+\delta),\ldots\), and multiply by
\(G^{-j}=t^j(1+u)^{-j}\).  At t-order three, one needs derivatives through
\(a'''\), \(b''\), and \(c'\).  Collecting powers gives
\cref{eq:h0,eq:h1,eq:h2,eq:h3}.  The calculation is finite and is also an
immediate specialization of \cref{eq:regular-composition-formula}.
\end{proof}

These formulas display an important asymmetry: derivatives fall on the outer
coefficient functions \(a,b,c,d\), not on the inner functions \(A,B,C,D\).
This is because the outer logarithmic argument is Taylor-expanded around
\(L\); the inner coefficients merely supply the increment.

\begin{workedbox}[title={One step toward Lambert \(W\)}]
Let
\[
 F(Y)=Y+\log Y,
 \qquad
 G(X)=X-L+\frac{L}{X}.
\]
In \cref{thm:near-composition-coeffs}, the outer coefficients are
\(a(L)=L\) and \(b=c=d=0\), while \(A=-L\), \(B=L\), and \(C=D=0\).  Hence
\[
 h_0=-L+L=0,
 \qquad
 h_1=(-L)(1)+L=0,
\]
and
\[
 h_2=-\frac{L^2}{2}+L.
\]
Thus
\[
 F(G)=X+\left(L-\frac{L^2}{2}\right)X^{-2}+\FormalBigOAt{t^3}{t}.
\]
Adding
\[
 C(L)=\frac{L^2}{2}-L=\frac{L(L-2)}2
\]
to the coefficient of \(X^{-2}\) in \(G\) cancels this residual.  This is the
second Lambert polynomial.
\end{workedbox}

\section{Formal Taylor composition}
For a near-identity substitution \(G=X+H\) with \(H/X\in t\K[L][[t]]\), one
may also write
\begin{equation}\label{eq:formal-taylor-composition}
 F(X+H)=\sum_{r\ge0}\frac{F^{(r)}(X)}{r!}H^r.
\end{equation}
This is t-adically meaningful for polynomial--logarithmic \(F\): each
derivative of a block raises its t-order by one, while \(H\) has t-valuation at
least zero, so the r-th Taylor term gains at least r t-orders relative to the
same outer block.

\begin{proposition}[Taylor formula in the near-identity class]\label{prop:taylor-composition}
If \(F\in\PolynomialLogLaurentRing{\K}\) and \(H/X\in t\K[L][[t]]\), then
\cref{eq:formal-taylor-composition} agrees coefficientwise with direct
substitution \(F\compose(X+H)\).
\end{proposition}
\begin{proof}
For a single block \(f_n(\log X)X^{-n}\), the claim is the ordinary formal
Taylor identity after replacing \(X\) by an indeterminate and remembering the
chain rule \cref{eq:repeated-derivative}.  Modulo \(t^N\), only finitely many
outer blocks and derivative orders contribute, so the identity extends by
finite summation.  Passing to all \(N\) proves coefficientwise equality.
\end{proof}

The Taylor form is conceptually simple and leads directly to Newton reversion.
The normal-ordered form is often more efficient because it packages the effect
of the logarithmic change and the reciprocal power in one shifted operator.

\section{Associativity, identity, and the chain rule}
Within the regular or near-identity composition classes above,
\[
 (F\compose G)\compose H=F\compose(G\compose H),
 \qquad
 F\compose X=X\compose F=F.
\]
These statements may be checked at every finite t-precision, where all
substitutions reduce to ordinary finite algebra.  Differentiation satisfies
\begin{equation}\label{eq:formal-chain-rule}
 (F\compose G)'=(F'\compose G)G'.
\end{equation}
Again, the proof holds blockwise and then extends by t-adic completeness.
These properties turn suitable near-identity series into a group under
composition once compositional inverses are constructed.

\section{Scaling and power-leading inner arguments}
The same method works for
\begin{equation}\label{eq:power-leading-inner}
 G(X)=\lambda X^p(1+u),
 \qquad p>0.
\end{equation}
Now
\[
 \log G=pL+\log\lambda+\delta,
 \qquad
 G^{-n}=\lambda^{-n}X^{-pn}\EulerE^{-n\delta}.
\]
Thus
\begin{equation}\label{eq:power-leading-composition}
 F\compose G
 =\sum_n\lambda^{-n}X^{-pn}
 \mathopen{:}\EulerE^{\delta(\Dop-n)}\mathclose{:}
 f_n(pL+\log\lambda),
\end{equation}
where \(\Dop\) differentiates the displayed argument.  If \(p\) is not an
integer, generalized t-exponents are required, but no new logarithmic level
appears.

\section{The first closure failure: a logarithmic prefactor}
Suppose instead that
\begin{equation}\label{eq:log-prefactor-inner}
 G(X)=\lambda X^pL^\beta(1+u).
\end{equation}
Then
\begin{equation}\label{eq:nested-log-created}
 \log G=pL+\beta\log L+\log\lambda+\delta.
\end{equation}
If any outer coefficient depends nontrivially on \(\log Y\), substitution
creates \(\log L=\log\log X\).  The one-level ring \(\K[L]((t))\) is therefore
not closed under this composition unless \(\beta=0\) or the outer logarithmic
dependence happens to cancel.

\begin{example}
Take \(F(Y)=\log Y\) and \(G(X)=XL\).  Then
\[
 F(G)=\log X+\log\log X=L+\log L,
\]
which cannot be expressed with polynomial coefficients in \(L\) alone.
\end{example}

The correct remedy is not to suppress the new term but to enlarge the
coefficient algebra to include \(L_2=\log L\), and, if necessary, further
iterates \(L_3=\log L_2\), and so on.  This is one route from elementary
power--log series to logarithmic transseries.

\section{Why arbitrary infinite substitution can fail}
Even if each individual monomial can be substituted, an infinite sum may fail
to define a valid series.  Two pathologies are typical:
\begin{enumerate}
\item infinitely many distinct outer monomials may collapse onto the same
      output monomial, producing an undefined infinite coefficient sum;
\item the substituted support may cease to be well ordered in the asymptotic
      direction.
\end{enumerate}
The t-adic class with \(G\sim\lambda X\) avoids these problems because, below
any cutoff, only finitely many outer t-blocks and Taylor orders contribute.
General transseries composition requires more sophisticated support theorems;
Edgar's treatment gives a detailed account \cite{edgar-composition}.

\section{A composition checklist}
Before composing, answer the following questions.
\begin{enumerate}
\item Is the inner argument asymptotically large in the intended analytic
      region?
\item Can it be factored as \(\lambda X^p(1+u)\) with \(u\prec1\)?
\item Does \(\log G\) introduce \(\log L\) or a further logarithmic level?
\item Do the output t-exponents remain in the chosen support grid?
\item At the requested precision, which outer blocks and Taylor orders can
      actually contribute?
\item Has the result been checked by an independent expansion or by a
      functional residual?
\end{enumerate}

\begin{summarybox}
For \(G=\lambda X(1+u)\), composition is governed by
\[
 F\compose G
 =\sum_n\lambda^{-n}X^{-n}
 \mathopen{:}\EulerE^{\delta(\Dop-n)}\mathclose{:}
 f_n(L+\log\lambda),
 \qquad \delta=\log(1+u).
\]
The normal ordering is essential because \(\delta\) depends on \(L\).  Bell
polynomials give a finite coefficient formula.  Near the identity, explicit
universal coefficients reveal how derivatives of the outer logarithmic
polynomials enter.  A leading factor \(L^\beta\) in the inner argument creates
\(\log L\) and forces a larger iterated-log class.
\end{summarybox}

\chapter{Series reversal and compositional inverses}\label{ch:reversion}

\section{Three meanings of ``inverse''}
The word \emph{inverse} is used for several different operations:
\[
 -A\quad\text{(additive inverse)},
 \qquad
 \recip{A}=\frac1A\quad\text{(multiplicative inverse)},
\]
\[
 \rev{F}\quad\text{(compositional inverse)}.
\]
This chapter concerns the last operation.  It is also called \emph{series
reversion} or \emph{series reversal}: given \(F\), find \(G\) such that
\begin{equation}\label{eq:inverse-equations}
 F\compose G=X,
 \qquad
 G\compose F=X.
\end{equation}
It does not mean writing the coefficients in reverse order.

\section{The near-identity composition group}
Define
\begin{equation}\label{eq:near-id-group}
 \NearIdentityPolynomialLogGroup{\K}=
 \left\{
 X+\sum_{n\ge0}f_n(L)X^{-n}: f_n\in\K[L]
 \right\}.
\end{equation}
Every \(F\in\NearIdentityPolynomialLogGroup{\K}\) satisfies \(F/X=1+\LittleO(1)\), so the composition formulas of
\cref{ch:composition} apply.  The set is closed under composition, and \(X\)
is the identity element.

The key triangular fact is that a new inner coefficient first appears
unchanged at its own t-order.

\begin{lemma}[Leading effect of a correction]\label{lem:leading-correction}
Let \(F\in\NearIdentityPolynomialLogGroup{\K}\), let \(G\in\NearIdentityPolynomialLogGroup{\K}\), and let
\(H=C(L)X^{-N}\) with \(N\ge0\).  Then
\begin{equation}\label{eq:leading-correction}
 F(G+H)=F(G)+C(L)X^{-N}+\FormalBigOAt{t^{N+1}}{t}.
\end{equation}
\end{lemma}
\begin{proof}
The formal Taylor formula gives
\[
 F(G+H)=F(G)+(F'\compose G)H+\frac12(F''\compose G)H^2+\cdots.
\]
For \(F\in\NearIdentityPolynomialLogGroup{\K}\), \(F'=1+\FormalBigOAt{t}{t}\), while \(F''=\FormalBigOAt{t^2}{t}\).  Thus the
linear term is \(H+\FormalBigOAt{t^{N+1}}{t}\), and all higher Taylor terms have t-order
strictly greater than \(N\).
\end{proof}

\begin{theorem}[Existence and uniqueness of near-identity reversal]\label{thm:reversion-exists}
Every \(F\in\NearIdentityPolynomialLogGroup{\K}\) has a unique compositional inverse
\(\rev{F}\in\NearIdentityPolynomialLogGroup{\K}\).
\end{theorem}
\begin{proof}
Construct
\[
 G_N=X+A_0(L)+A_1(L)t+\cdots+A_N(L)t^N
\]
inductively so that
\[
 F\compose G_N=X+\FormalBigOAt{t^{N+1}}{t}.
\]
At order zero, choose \(A_0\) to cancel the t-zero residual.  Suppose the
coefficients through \(A_{N-1}\) have been chosen and
\[
 F\compose G_{N-1}=X+r_N(L)t^N+\FormalBigOAt{t^{N+1}}{t}.
\]
By \cref{lem:leading-correction}, replacing \(G_{N-1}\) by
\(G_N=G_{N-1}-r_Nt^N\) cancels the residual at order \(N\) without altering
lower orders.  Q-adic completeness produces a formal right inverse \(G\).

For uniqueness, if \(G\) and \(\widetilde G\) first differ by
\(C(L)t^N\), then \cref{lem:leading-correction} shows that
\(F\compose G\) and \(F\compose\widetilde G\) first differ by the same nonzero
block, a contradiction.  Finally,
\[
 F\compose(G\compose F)=(F\compose G)\compose F=F=F\compose X.
\]
The same leading-difference argument shows that left composition by \(F\) is
injective, hence \(G\compose F=X\).  Therefore the right inverse is two-sided.
\end{proof}

\section{Residual lifting: one block at a time}
The proof gives a particularly simple algorithm.

\begin{algorithm}[Triangular series reversal]\label{alg:triangular-reversion}
To compute \(G=\rev{F}\) through t-order \(N\):
\begin{enumerate}
\item Start with \(G_{-1}=X\).
\item For \(n=0,1,\ldots,N\), compute the residual
\[
 R_n=\bigl[F\compose G_{n-1}-X\bigr]_{<n+1}.
\]
\item Extract its t-order \(n\) coefficient \(r_n(L)=[t^n]R_n\).
\item Set
\[
 G_n=G_{n-1}-r_n(L)t^n.
\]
\end{enumerate}
Then \(F\compose G_N=X+\FormalBigOAt{t^{N+1}}{t}\).
\end{algorithm}

This algorithm is ideal for derivations and verification because it exposes the
residual that determines each new coefficient.  It uses repeated composition,
so a high-order implementation should switch to Newton iteration.

\section{Universal inverse coefficients}
Let
\begin{equation}\label{eq:F-four-blocks}
 F(X)=X+a(L)+\frac{b(L)}X+\frac{c(L)}{X^2}
      +\frac{d(L)}{X^3}+\FormalBigOAt{t^4}{t},
\end{equation}
and write
\begin{equation}\label{eq:G-four-blocks}
 \rev{F}(X)=X+A(L)+\frac{B(L)}X+\frac{C(L)}{X^2}
      +\frac{D(L)}{X^3}+\FormalBigOAt{t^4}{t}.
\end{equation}

\begin{theorem}[Reversal through fourth block]\label{thm:inverse-coefficients}
The coefficients in \cref{eq:G-four-blocks} are
\begin{align}
 A={}&-a,\label{eq:inverse-A}\\
 B={}&aa'-b,\label{eq:inverse-B}\\
 C={}&\frac12a^2a'-\frac12a^2a''-a(a')^2-ab+ab'+a'b-c,
 \label{eq:inverse-C}\\
 D={}&\frac13a^3a'-\frac12a^3a''+\frac16a^3a'''
 -\frac32a^2(a')^2+\frac32a^2a'a''\notag\\
 &{}-a^2b+\frac32a^2b'-\frac12a^2b''
 +a(a')^3+3aa'b-2aa'b'-aa''b\notag\\
 &{}-2ac+ac'-(a')^2b+a'c-b^2+bb'-d.
 \label{eq:inverse-D}
\end{align}
\end{theorem}
\begin{proof}
Substitute \cref{eq:F-four-blocks,eq:G-four-blocks} into the universal
composition coefficients \cref{eq:h0,eq:h1,eq:h2,eq:h3}.  Set
\(h_0=h_1=h_2=h_3=0\) and solve successively for \(A,B,C,D\).  The system is
triangular: \(A\) occurs only in the first equation, \(B\) first occurs linearly
in the second, and so on.  Straight simplification gives the displayed
formulas.
\end{proof}

\begin{checkpointbox}
The derivatives in \cref{eq:inverse-A,eq:inverse-B,eq:inverse-C,eq:inverse-D}
are derivatives of the \emph{original} coefficient functions \(a,b,c,d\) with
respect to \(L\).  No derivative of an unknown inverse coefficient is needed.
This makes the formulas directly evaluable.
\end{checkpointbox}

\section{A parameterized example}
Consider
\begin{equation}\label{eq:param-inverse-example}
 F(X)=X+\alpha L+\frac{\beta}{X}+\FormalBigOAt{t^2}{t},
\end{equation}
where \(\alpha,\beta\in\K\).  Here \(a=\alpha L\), \(a'=\alpha\), and
\(b=\beta\).  The inverse begins
\begin{equation}\label{eq:param-inverse-result}
\begin{split}
 \rev{F}(X)={}&X-\alpha L
 +\frac{\alpha^2L-\beta}{X}\\
 &+\frac{\alpha^3(\tfrac12L^2-L)-\alpha\beta L+\alpha\beta}{X^2}
 +\FormalBigOAt{t^3}{t}.
\end{split}
\end{equation}
A direct substitution into \cref{eq:param-inverse-example} verifies that the
residual is \(\FormalBigOAt{t^3}{t}\).

\section{The Lagrange--Bürmann reversion formula at infinity}
There is also a closed differential expression for all coefficients.  Write
\begin{equation}\label{eq:F-x-plus-h}
 F(X)=X+h(X),
 \qquad h(X)=\LittleO(X).
\end{equation}
For our core class, \(h\in\K[L][[t]]\): it begins at t-order zero, but is small
relative to \(X=t^{-1}\).

\begin{theorem}[Differential reversion formula]\label{thm:lagrange-infinity}
For \(F=X+h\in\NearIdentityPolynomialLogGroup{\K}\),
\begin{equation}\label{eq:lagrange-infinity}
 \rev{F}(X)
 =X+\sum_{m\ge1}\frac{(-1)^m}{m!}
 \left(\frac{\Differential}{\Differential X}\right)^{m-1}\!\bigl(h(X)^m\bigr).
\end{equation}
The m-th summand has t-valuation at least \(m-1\), so only finitely many
summands contribute at any fixed precision.
\end{theorem}
\begin{proof}
Introduce a formal homotopy parameter \(\tau\) and consider
\(F_\tau(Y)=Y+\tau h(Y)\).  The classical Lagrange--Bürmann formula for the inverse
of a perturbation of the identity gives
\[
 \CompositionalInverseOf{F_\tau}(X)
 =X+\sum_{m\ge1}\frac{(-\tau)^m}{m!}
 \left(\frac{\Differential}{\Differential X}\right)^{m-1}h(X)^m.
\]
For polynomial--logarithmic \(h\), each derivative raises t-order by at least
one, so the coefficient of any fixed t-power is a finite polynomial in \(\tau\).
Setting \(\tau=1\) is therefore formally legitimate and yields
\cref{eq:lagrange-infinity}.  The result is the unique inverse from
\cref{thm:reversion-exists}.
\end{proof}

\begin{remark}
Formula~\eqref{eq:lagrange-infinity} is often the fastest route to a symbolic
closed form for a family of inverse coefficients.  The residual-lifting
algorithm is usually the clearest route to a proof by induction.  Newton
iteration, developed next, is usually the fastest route to many coefficients.
\end{remark}

\section{Newton reversion}
Let \(G\) approximate \(\rev{F}\).  Newton's correction for the equation
\(F(G)=X\) is
\begin{equation}\label{eq:newton-reversion}
 G_{\mathrm{new}}
 =G-\frac{F\compose G-X}{F'\compose G}.
\end{equation}
The denominator is a unit because \(F'=1+\FormalBigOAt{t}{t}\).

\begin{theorem}[Precision doubling for reversion]\label{thm:newton-reversion}
Suppose
\[
 E=F\compose G-X,
 \qquad \val_t(E)\ge N\ge1.
\]
Then the Newton update \cref{eq:newton-reversion} satisfies
\begin{equation}\label{eq:newton-reversion-precision}
 \val_t(F\compose G_{\mathrm{new}}-X)\ge2N.
\end{equation}
\end{theorem}
\begin{proof}
Set
\[
 \Delta=-\frac{E}{F'\compose G}.
\]
Since the denominator is a unit, \(\val_t(\Delta)\ge N\).  Taylor expansion
around \(G\) gives
\[
 F(G+\Delta)-X
 =E+(F'\compose G)\Delta
 +\sum_{r\ge2}\frac{F^{(r)}\compose G}{r!}\Delta^r.
\]
The first two terms cancel.  Every remaining term contains \(\Delta^r\) with
r at least two, hence has t-valuation at least \(2N\); in fact the derivatives
of a near-identity series improve the bound further.  This proves
\cref{eq:newton-reversion-precision}.
\end{proof}

A practical schedule is:
\begin{enumerate}
\item determine the t-zero correction by residual lifting;
\item compute one more block if necessary so that the residual has positive
      t-valuation;
\item apply Newton updates at precisions \(2,4,8,\ldots\), truncating all
      intermediate operations.
\end{enumerate}

\section{Left and right residuals}
To verify a computed inverse \(G_N\), check both
\begin{equation}\label{eq:two-residuals}
 F\compose G_N-X=\FormalBigOAt{t^{N+1}}{t},
 \qquad
 G_N\compose F-X=\FormalBigOAt{t^{N+1}}{t}.
\end{equation}
In exact formal algebra, one side and uniqueness imply the other.  In software,
the two computations exercise different code paths and are therefore valuable
independent tests.

\section{Reversion away from the identity}
The near-identity group is not the only invertible asymptotic class.  The first
step in a more general problem is to invert the dominant monomial.

\subsection{Algebraic leading term}
If
\[
 F(Y)\sim cY^p,
 \qquad c\ne0,
 \qquad p>0,
\]
then a branch of the inverse begins
\[
 Y\sim c^{-1/p}X^{1/p}.
\]
Set \(Y=c^{-1/p}X^{1/p}(1+U)\), substitute, and solve for the small series
\(U\).  Fractional t-powers generally appear.

\subsection{Power times logarithm}
If
\begin{equation}\label{eq:power-log-leading}
 F(Y)\sim cY^p(\log Y)^\beta,
\end{equation}
then a first balance gives
\begin{equation}\label{eq:power-log-leading-inverse}
 Y\sim
 \left(\frac{p^\beta}{c}\right)^{1/p}
 X^{1/p}(\log X)^{-\beta/p}.
\end{equation}
Indeed, \(\log Y\sim p^{-1}\log X\).  The next correction generally contains
\(\log\log X/\log X\).  Thus the natural inverse class is larger than
\(\K[L]((X^{-1}))\).

\subsection{Branch information}
Formal reversal determines a family of possible expansions once roots and
logarithms are chosen.  Analytic inversion requires a domain on which \(F\) is
one-to-one and its derivative does not vanish.  Different choices of roots or
logarithm branches can produce different formal inverse sectors.

\section{A reversion checklist}
\begin{enumerate}
\item Distinguish compositional inversion from reciprocal formation.
\item Normalize the dominant term before assuming a near-identity form.
\item Choose a support class large enough for fractional powers or iterated
      logarithms created by the leading inversion.
\item Use residual lifting for transparent coefficient derivation.
\item Use Newton reversion for high order.
\item Verify at least one composition residual, preferably both.
\item Keep analytic branch choices separate from formal coefficient algebra.
\end{enumerate}

\begin{summarybox}
Near-identity polynomial--logarithmic series form a composition group.  Their
inverse exists uniquely because a correction \(C(L)X^{-N}\) first changes the
composition residual by exactly the same block.  Universal inverse formulas
through t-order three follow from the composition formulas.  The
Lagrange--Bürmann identity
\[
 \rev{(X+h)}=X+\sum_{m\ge1}\frac{(-1)^m}{m!}
 \left(\frac{\Differential}{\Differential X}\right)^{m-1}h^m
\]
provides a closed all-orders expression, while Newton's method doubles formal
precision.  Non-tangent leading terms must first be inverted at the monomial
level and may force fractional powers or iterated logarithms.
\end{summarybox}

\chapter{Lambert \texorpdfstring{\(W\)}{W} as the central example}\label{ch:lambert}

\section{Definition and branches}
Lambert's \(W\) function is the multivalued inverse of
\[
 w\longmapsto w\EulerE^w.
\]
Thus every branch satisfies
\begin{equation}\label{eq:Lambert-def}
 W_k(z)\EulerE^{W_k(z)}=z,
 \qquad k\in\IntegerNumbers.
\end{equation}
The principal branch is \(W_0\).  On the real interval \((-\EulerE^{-1},0)\), there
is a second real branch \(W_{-1}\).  The branch geometry, singularity at
\(-\EulerE^{-1}\), numerical evaluation, and standard asymptotic formulas are
developed in Corless et al. and in the NIST Digital Library of Mathematical
Functions \cite{corless-lambert,dlmf-lambert}.

For a chosen branch and a large complex argument, introduce
\begin{equation}\label{eq:xi-k}
 X=\xi_k:=\PrincipalLogarithm z+2\pi\ImaginaryUnit k,
 \qquad L=\log X,
\end{equation}
with compatible logarithm branches.  Taking a logarithm of
\cref{eq:Lambert-def} gives the asymptotic inverse equation
\begin{equation}\label{eq:W-log-equation}
 W+\log W=X.
\end{equation}
The same equation defines the Wright omega function once branch unwinding is
handled globally.  Formally, it is the inverse of \(Y\mapsto Y+\log Y\).

\section{The ansatz and its fixed-point equation}
Since \(W\sim X\), the first correction in \cref{eq:W-log-equation} is
\(-\log X=-L\).  Write
\begin{equation}\label{eq:W-ansatz}
 W=X-L+R,
 \qquad
 R=\sum_{n\ge1}\LambertPolynomial{n}(L)X^{-n}.
\end{equation}
Throughout this article the shared Lambert family is zero-based, with
\(\LambertPolynomial{0}(u)=-u\).  Equivalently,
\[
 W=X+\sum_{n\ge0}\LambertPolynomial{n}(L)X^{-n}.
\]
Then
\[
 \log W
 =L+\log\!\left(1-\frac{L-R}{X}\right).
\]
Substitution into \cref{eq:W-log-equation} yields
\begin{equation}\label{eq:W-fixed-point}
 R=-\log\bigl(1-t(L-R)\bigr),
 \qquad t=X^{-1}.
\end{equation}
This single equation determines all Lambert polynomials \(\LambertPolynomial{n}\).

\begin{proposition}[Triangular Lambert recurrence]\label{prop:Lambert-triangular}
Set \(R_{n-1}=\sum_{j=1}^{n-1}\LambertPolynomial{j}(L)t^j\).  Then
\begin{equation}\label{eq:Lambert-triangular-recurrence}
 \LambertPolynomial{n}(L)=[t^n]\left[-\log\bigl(1-t(L-R_{n-1})\bigr)\right].
\end{equation}
The right side involves only \(\LambertPolynomial{1},\ldots,\LambertPolynomial{n-1}\).
\end{proposition}
\begin{proof}
In \(t(L-R)\), the occurrence of \(\LambertPolynomial{n}t^n\) is shifted to t-order \(n+1\).
Therefore the coefficient of \(t^n\) on the right of
\cref{eq:W-fixed-point} depends only on earlier polynomials.  Matching that
coefficient with the left side gives \cref{eq:Lambert-triangular-recurrence}.
\end{proof}

\begin{workedbox}[title={The first four Lambert polynomials}]
From \cref{eq:W-fixed-point},
\[
 R=t(L-R)+\frac{t^2}{2}(L-R)^2
   +\frac{t^3}{3}(L-R)^3+\cdots.
\]
At t-order one,
\[
 \LambertPolynomial{1}=L.
\]
At t-order two,
\[
 \LambertPolynomial{2}=-\LambertPolynomial{1}+\frac{L^2}{2}
 =\frac{L(L-2)}2.
\]
At t-order three,
\[
 \LambertPolynomial{3}=-\LambertPolynomial{2}-L\LambertPolynomial{1}+\frac{L^3}{3}
 =\frac{L(2L^2-9L+6)}6.
\]
A further collection gives
\[
 \LambertPolynomial{4}=\frac{L(3L^3-22L^2+36L-12)}{12}.
\]
\end{workedbox}

\section{Derivation by series reversal}
Equation~\eqref{eq:W-log-equation} asks for the inverse of
\[
 F(Y)=Y+\log Y.
\]
Apply the differential reversion formula \cref{eq:lagrange-infinity} with
\(h(X)=\log X=L\):
\begin{equation}\label{eq:W-lagrange}
 W(\EulerE^X)
 =X+\sum_{m\ge1}\frac{(-1)^m}{m!}
 \left(\frac{\Differential}{\Differential X}\right)^{m-1}L^m.
\end{equation}
The term \(m=1\) is \(-L\).  For \(n=m-1\ge1\), comparison with
\cref{eq:W-ansatz} gives the compact operator formula
\begin{equation}\label{eq:Pn-operator}
 \LambertPolynomial{n}(L)=\frac{(-1)^{n+1}}{(n+1)!}
 \prod_{j=0}^{n-1}(\Dop-j)\,L^{n+1},
 \qquad \Dop=\frac{\Differential}{\Differential L}.
\end{equation}
This formula computes \(\LambertPolynomial{n}\) independently of all earlier polynomials.

\section{Stirling-cycle-number formula}
Let \(\stirlingone{n}{j}\) denote the unsigned Stirling number of the first kind,
the number of permutations of \(n\) elements with exactly \(j\) cycles.  The
signed Stirling numbers \(s(n,j)\) satisfy
\[
 \prod_{r=0}^{n-1}(z-r)=\sum_{j=0}^ns(n,j)z^j,
 \qquad
 s(n,j)=(-1)^{n-j}\stirlingone{n}{j}.
\]
Expanding the operator product in \cref{eq:Pn-operator} yields the standard
closed formula.

\begin{theorem}[Lambert polynomial formula]\label{thm:Lambert-Stirling}
For \(n\ge1\),
\begin{equation}\label{eq:Lambert-Stirling}
 \LambertPolynomial{n}(L)=
 \sum_{m=1}^{n}(-1)^{n+m}
 \stirlingone{n}{n-m+1}\frac{L^m}{m!}.
\end{equation}
Consequently,
\begin{equation}\label{eq:W-all-orders}
 W_k(z)\sim
 X-L+\sum_{n\ge1}\frac1{X^n}
 \sum_{m=1}^{n}(-1)^{n+m}
 \stirlingone{n}{n-m+1}\frac{L^m}{m!},
\end{equation}
where \(X=\PrincipalLogarithm z+2\pi\ImaginaryUnit k\) and \(L=\log X\).
\end{theorem}
\begin{proof}
From \cref{eq:Pn-operator},
\[
 \LambertPolynomial{n}(L)=\frac{(-1)^{n+1}}{(n+1)!}
 \sum_{j=0}^ns(n,j)\Dop^jL^{n+1}.
\]
The \(j=0\) term vanishes because the falling-factorial operator begins with
\(\Dop\), and
\[
 \Dop^jL^{n+1}=\frac{(n+1)!}{(n+1-j)!}L^{n+1-j}.
\]
Put \(m=n+1-j\).  Then \(j=n-m+1\), and the sign becomes
\[
 (-1)^{n+1}(-1)^{n-j}=(-1)^{n+m}.
\]
This gives \cref{eq:Lambert-Stirling}.  Substitution into the ansatz proves
\cref{eq:W-all-orders}.
\end{proof}

This formula was historically connected with the inversion of
\(y^\alpha\EulerE^y\) and associated Stirling numbers
\cite{jeffrey-inversion,corless-lambert}.  It is also the form recorded in
standard references such as the DLMF \cite{dlmf-lambert}.

\section{Structural properties of the Lambert polynomials}
Several facts are immediate from \cref{eq:Lambert-Stirling}.

\begin{proposition}\label{prop:Pn-properties}
For \(n\ge1\):
\begin{enumerate}
\item \(\LambertPolynomial{n}\) has degree \(n\) and is divisible by \(L\);
\item its leading term is \(L^n/n\);
\item its linear term is \((-1)^{n+1}L\);
\item the nonzero coefficients alternate in sign as the power of \(L\)
      increases;
\item \(n!\LambertPolynomial{n}(L)\in\IntegerNumbers[L]\).
\end{enumerate}
\end{proposition}
\begin{proof}
The term with \(m=n\) uses
\(\stirlingone{n}{1}=(n-1)!\) and gives \(L^n/n\).  The term with \(m=1\)
uses \(\stirlingone{n}{n}=1\).  Every summand has \(m\ge1\), proving
divisibility by \(L\), and its sign is \((-1)^{n+m}\).  Finally,
\(n!/m!\) is an integer for \(m\le n\), so multiplying
\cref{eq:Lambert-Stirling} by \(n!\) clears all denominators.
\end{proof}

\section{The expansion through tenth order}
The first ten coefficient polynomials are listed here for reference:
\begin{align*}
\LambertPolynomial{1}={}&L,\\
\LambertPolynomial{2}={}&\frac{L(L-2)}2,\\
\LambertPolynomial{3}={}&\frac{L(2L^2-9L+6)}6,\\
\LambertPolynomial{4}={}&\frac{L(3L^3-22L^2+36L-12)}{12},\\
\LambertPolynomial{5}={}&\frac{L(12L^4-125L^3+350L^2-300L+60)}{60},\\
\LambertPolynomial{6}={}&\frac{L}{120}\bigl(20L^5-274L^4+1125L^3
          -1700L^2+900L-120\bigr),\\
\LambertPolynomial{7}={}&\frac{L}{840}\bigl(120L^6-2058L^5+11368L^4-25725L^3\\[-0.2ex]
&\hspace{7.4em}{}+24500L^2-8820L+840\bigr),\\
\LambertPolynomial{8}={}&\frac{L}{2520}\bigl(315L^7-6534L^6+45962L^5-142149L^4\\[-0.2ex]
&\hspace{7.4em}{}+205800L^3-135240L^2+35280L-2520\bigr),\\
\LambertPolynomial{9}={}&\frac{L}{2520}\bigl(280L^8-6849L^7+59062L^6-235494L^5\\[-0.2ex]
&\hspace{7.4em}{}+471429L^4-476280L^3+229320L^2\\[-0.2ex]
&\hspace{7.4em}{}-45360L+2520\bigr),\\
\LambertPolynomial{10}={}&\frac{L}{10080}\bigl(1008L^9-28516L^8+293175L^7
 -1447360L^6\\[-0.2ex]
&\hspace{7.4em}{}+3770550L^5-5314932L^4+3969000L^3\\[-0.2ex]
&\hspace{7.4em}{}-1461600L^2+226800L-10080\bigr).
\end{align*}

\section{Formal asymptotics versus actual convergence}
Most asymptotic series in applications are divergent, so one should not assume
convergence merely from a formal derivation.  The Lambert expansion is unusually
well behaved.  The DLMF records absolute convergence of the standard double
series for sufficiently large \(\AbsoluteValue{z}\), and convergence on the principal real
branch for \(z\ge\EulerE\) \cite{dlmf-lambert}.  The formal derivation remains useful
outside a proved convergence region because it determines asymptotic
coefficients and branch structure.

Even for a convergent series, the absolute error of successive truncations need
not decrease monotonically at every small argument.  Cancellation can make one
partial sum better than the next.

\begin{table}[htbp]
\centering
\caption{Absolute error of the truncation
\(W_N=X-L+\sum_{n=1}^N\LambertPolynomial{n}(L)X^{-n}\) on the principal real branch.}
\label{tab:W-errors}
\begin{tabular}{@{}rccc@{}}
\toprule
\(N\) & \(z=10\) & \(z=10^3\) & \(z=10^{20}\) \\
\midrule
0  & \(2.770\times10^{-1}\) & \(2.745\times10^{-1}\) & \(8.482\times10^{-2}\) \\
2  & \(6.468\times10^{-3}\) & \(3.923\times10^{-3}\) & \(3.530\times10^{-6}\) \\
4  & \(3.255\times10^{-3}\) & \(1.196\times10^{-4}\) & \(1.247\times10^{-7}\) \\
6  & \(7.427\times10^{-4}\) & \(4.019\times10^{-6}\) & \(1.186\times10^{-10}\) \\
8  & \(1.115\times10^{-4}\) & \(1.183\times10^{-7}\) & \(6.340\times10^{-13}\) \\
10 & \(2.176\times10^{-6}\) & \(1.584\times10^{-9}\) & \(1.728\times10^{-15}\) \\
\bottomrule
\end{tabular}
\end{table}

\section{Residuals and a posteriori correction}
For a truncation \(W_N\), define the logarithmic residual
\begin{equation}\label{eq:W-residual}
 \mathcal R_N=W_N+\log W_N-X.
\end{equation}
Formal construction gives \(\mathcal R_N=\FormalBigOAt{t^{N+1}}{t}\).  Since the derivative
of \(w+\log w-X\) is \(1+1/w\), one Newton correction is
\begin{equation}\label{eq:W-newton-correction}
 W_N^{\mathrm{new}}
 =W_N-\frac{W_N+\log W_N-X}{1+1/W_N}.
\end{equation}
As a formal series, this doubles the t-adic precision.  Numerically, the
residual also provides an inexpensive a posteriori diagnostic.  On a domain
where \(\AbsoluteValue{1+1/w}\) is bounded away from zero, the true error is approximately
\(-\mathcal R_N/(1+1/W_N)\).

One may instead use the original residual \(W_N\EulerE^{W_N}-z\), but
\cref{eq:W-residual} is usually better scaled for large \(z\).

\section{The real branch \texorpdfstring{\(W_{-1}\)}{W-1} near zero}
Let \(x\to0^-\) and define
\begin{equation}\label{eq:Wm1-TL}
 T=\log\frac1{-x},
 \qquad L=\log T.
\end{equation}
Writing \(W_{-1}(x)=-u\) transforms
\(W\EulerE^W=x\) into
\[
 u-\log u=T.
\]
Equivalently, use the general branch formula with the real asymptotic
coordinates.  The result is
\begin{equation}\label{eq:Wm1-expansion}
 W_{-1}(x)
 \sim -T-L+\sum_{n\ge1}\frac{(-1)^n\LambertPolynomial{n}(L)}{T^n}.
\end{equation}
Thus
\begin{equation}\label{eq:Wm1-first}
\begin{split}
 W_{-1}(x)={}&-T-L-\frac{L}{T}
 +\frac{L(L-2)}{2T^2}\\
 &-\frac{L(2L^2-9L+6)}{6T^3}
 +\frac{L(3L^3-22L^2+36L-12)}{12T^4}+\cdots.
\end{split}
\end{equation}
The alternating powers of \(T^{-1}\) come from replacing the large logarithmic
coordinate \(X\) by \(-T\).

\begin{warningbox}
For \(x<0\), writing simply \(\log x\) hides a branch choice.  The real formula
\cref{eq:Wm1-expansion} avoids that ambiguity by using the positive quantity
\(T=\log(1/(-x))\).  Complex branch formulas should instead retain
\(X=\PrincipalLogarithm z+2\pi\ImaginaryUnit k\) explicitly.
\end{warningbox}

\section{The family \texorpdfstring{\(y+\alpha\log y=X\)}{y + alpha log y = X}}
Consider
\begin{equation}\label{eq:alpha-log-inversion}
 y+\alpha\log y=X,
 \qquad \alpha\ne0.
\end{equation}
Exponentiating gives \(y^\alpha\EulerE^y=\EulerE^X\), and an exact solution is
\begin{equation}\label{eq:alpha-exact-W}
 y=\alpha W\!\left(\frac{\EulerE^{X/\alpha}}{\alpha}\right),
\end{equation}
with compatible branches.  Direct reversion with \(h(X)=\alpha\log X\) gives a
particularly simple scaling law:
\begin{equation}\label{eq:alpha-expansion}
 y=X-\alpha L+
 \sum_{n\ge1}\frac{\alpha^{n+1}\LambertPolynomial{n}(L)}{X^n}.
\end{equation}
The first terms are
\begin{equation}\label{eq:alpha-first-terms}
\begin{split}
 y={}&X-\alpha L+\frac{\alpha^2L}{X}
 +\frac{\alpha^3L(L-2)}{2X^2}\\
 &+\frac{\alpha^4L(2L^2-9L+6)}{6X^3}+\cdots.
\end{split}
\end{equation}
This follows immediately from \cref{eq:W-lagrange}: the m-th reversion term is
multiplied by \(\alpha^m\).

\section{Perturbed Lambert equations}
Now consider
\begin{equation}\label{eq:perturbed-Lambert}
 F(y)=y+\alpha\log y+\frac{\beta}{y}+\frac{\gamma\log y}{y^2}=X.
\end{equation}
No new principle is required.  Place
\[
 a(L)=\alpha L,
 \qquad b(L)=\beta,
 \qquad c(L)=\gamma L
\]
in the universal inverse formulas of \cref{thm:inverse-coefficients}.  The
first blocks are
\begin{equation}\label{eq:perturbed-Lambert-inverse}
\begin{split}
 y={}&X-\alpha L+\frac{\alpha^2L-\beta}{X}\\
 &+\frac{\alpha^3(\tfrac12L^2-L)-\alpha\beta L+\alpha\beta-\gamma L}{X^2}
 +\FormalBigOAt{t^3}{t}.
\end{split}
\end{equation}
This illustrates the practical value of universal reversal formulas: once the
outer coefficient functions are identified, no fresh composition derivation is
needed.

\section{Branch-aware use in practice}
For complex computation, proceed in this order:
\begin{enumerate}
\item choose the Lambert branch index \(k\);
\item set \(X=\PrincipalLogarithm z+2\pi\ImaginaryUnit k\) using a specified base branch of \(\PrincipalLogarithm\);
\item choose the compatible branch of \(L=\log X\);
\item evaluate the polynomial--logarithmic truncation;
\item if desired, apply Newton or Halley iteration without crossing a branch
      cut.
\end{enumerate}
The formal coefficients are branch-independent; branch dependence enters
through \(X\) and \(L\).

\begin{summarybox}
Lambert \(W\) is the compositional inverse of \(Y+\log Y\) after setting
\(X=\PrincipalLogarithm z+2\pi\ImaginaryUnit k\).  Its coefficients satisfy a triangular fixed-point
recurrence, an all-orders differential formula, and the closed Stirling-number
identity
\[
 \LambertPolynomial{n}(L)=\sum_{m=1}^n(-1)^{n+m}
 \stirlingone{n}{n-m+1}\frac{L^m}{m!}.
\]
The same polynomials control every equation \(y+\alpha\log y=X\), with a factor
\(\alpha^{n+1}\).  Real \(W_{-1}\) near zero follows by replacing the large
coordinate with \(-T\).  Residuals and Newton corrections provide both formal
precision doubling and practical numerical refinement.
\end{summarybox}

\chapter{Closure, nested logarithms, and multiseries}\label{ch:closure}

\section{No single small ring is closed under everything}
The core ring \(\K[L]((t))\) is deliberately economical.  Its limitations are
mathematically informative:
\begin{itemize}
\item dividing by a nonconstant leading polynomial requires rational functions
      of \(L\);
\item composing with an inner factor \(L^\beta\) creates \(\log L\);
\item exponentiating a nonlinear polynomial such as \(L^2\) creates a new
      transmonomial;
\item reversing a leading term \(X^pL^\beta\) creates fractional powers and
      usually \(\log L\) corrections.
\end{itemize}
Rather than declaring these operations illegal, one builds a controlled
hierarchy of larger classes.

\section{A practical closure ladder}
Let
\[
 L_1=\log X,
 \qquad L_2=\log L_1,
 \qquad \ldots,
 \qquad L_{r+1}=\log L_r.
\]
A useful ladder is
\begin{equation}\label{eq:closure-ladder}
\begin{split}
 \K[L_1]((X^{-1}))
 &\subset \K(L_1)((X^{-1}))\\
 &\subset \K(L_1,L_2)((X^{-1}))\\
 &\subset \K(L_1,L_2,L_3)((X^{-1}))\subset\cdots.
\end{split}
\end{equation}
For finer asymptotic expansions, replace the rational-function coefficient
field by generalized Laurent or Hahn series in inverse powers of the successive
logs.  A typical finite-depth logarithmic monomial is
\begin{equation}\label{eq:log-monomial-tower}
 X^\rho L_1^{\alpha_1}L_2^{\alpha_2}\cdots L_r^{\alpha_r}.
\end{equation}
The dominance order reads from left to right: a change in the power of \(X\)
dominates every finite change in the powers of \(L_1,L_2,\ldots\); at equal
\(X\)-power, a change in \(L_1\) dominates every lower logarithmic level, and
so on.

\begin{scopebox}
The fields in \cref{eq:closure-ladder} are convenient coefficient models, not a
complete definition of logarithmic transseries.  Full closure requires support
conditions across all monomials \eqref{eq:log-monomial-tower}.  The general
transseries literature supplies those conditions.  For finite computations on
a fixed grid and fixed logarithmic depth, the displayed models are usually
sufficient.
\end{scopebox}

\section{Differentiating an iterated-log tower}
The derivatives are
\begin{equation}\label{eq:iterated-log-derivatives}
 \frac{\Differential L_1}{\Differential X}=\frac1X,
 \qquad
 \frac{\Differential L_2}{\Differential X}=\frac1{XL_1},
 \qquad
 \frac{\Differential L_3}{\Differential X}=\frac1{XL_1L_2},
\end{equation}
and in general
\begin{equation}\label{eq:Lj-derivative}
 \frac{\Differential L_j}{\Differential X}
 =\frac1{X L_1L_2\cdots L_{j-1}}.
\end{equation}
Thus a polynomial ring in \(L_1,\ldots,L_r\) is not closed under
differentiation once \(r\ge2\); negative powers of earlier logarithms appear.
The rational field \(\K(L_1,\ldots,L_r)\), or a Laurent-monomial field, repairs
this defect.

For a coefficient \(a(L_1,\ldots,L_r)\), the chain rule becomes
\begin{equation}\label{eq:tower-chain-operator}
 \frac{\Differential a}{\Differential X}
 =\frac1X\left(
 \partial_{L_1}
 +\frac1{L_1}\partial_{L_2}
 +\frac1{L_1L_2}\partial_{L_3}
 +\cdots
 +\frac1{L_1\cdots L_{r-1}}\partial_{L_r}
 \right)a.
\end{equation}
Define the tower derivation inside parentheses by \(\mathcal D_r\).  Then
\[
 \frac{\Differential}{\Differential X}\bigl(aX^{-n}\bigr)
 =\bigl(\mathcal D_r-n\bigr)a\,X^{-n-1},
\]
which is the exact analogue of \(\Delta_n=\Dop-n\).

\section{Composition depth}
Suppose
\[
 G(X)=\lambda X^pL_1^\beta(1+u).
\]
Then
\[
 \log G=pL_1+\beta L_2+\log\lambda+\log(1+u).
\]
Therefore an outer coefficient polynomial in \(\log Y\) becomes a polynomial in
\(L_1\) and \(L_2\).  If the inner leading factor also contains \(L_2^\gamma\),
then \(L_3\) appears.  In general, taking the logarithm of a monomial of depth
\(r\) can expose depth \(r+1\).

This gives a simple planning rule:
\begin{quote}
Before composition, factor the dominant inner monomial and take its logarithm.
Every logarithmic level visible there must be admitted in the output
coefficient class.
\end{quote}

\section{Changing the fast variable}
An expression with two logarithmic levels may become an ordinary
polynomial--logarithmic series after a change of dominant coordinate.  For
example,
\begin{equation}\label{eq:multiseries-form}
 \sum_{n\ge0}\frac{p_n(\log\log X)}{(\log X)^n}
\end{equation}
is polynomial--logarithmic in the new fast variable
\[
 T=\log X,
 \qquad \Lambda=\log T=\log\log X.
\]
Such nested expansions are often called \emph{multiseries}: one first expands
in a dominant small scale, then allows coefficients that themselves carry an
asymptotic expansion in a slower scale.  Algorithms for exp--log functions and
their inverses automatically discover these scales
\cite{rssv-exp-log,salvy-shackell}.

\section{Case study: inverting \texorpdfstring{\(y/\log y\)}{y/log y}}
Consider
\begin{equation}\label{eq:y-over-log-y}
 X=\frac{y}{\log y},
 \qquad X\to+\infty.
\end{equation}
This is not tangent to the identity.  Rearranging,
\[
 \log y=\frac yX,
 \qquad
 y=\EulerE^{y/X}.
\]
Multiplying by \(-1/X\) gives
\[
 -\frac yX\EulerE^{-y/X}=-\frac1X,
\]
so the large solution is exactly
\begin{equation}\label{eq:y-over-log-exact}
 y=-XW_{-1}(-X^{-1}).
\end{equation}
Set
\[
 T=\log X,
 \qquad \Lambda=\log T.
\]
Applying \cref{eq:Wm1-expansion} with \(x=-X^{-1}\) yields
\begin{equation}\label{eq:y-over-log-expansion}
 y\sim X\left[
 T+\Lambda+\sum_{n\ge1}\frac{(-1)^{n+1}\LambertPolynomial{n}(\Lambda)}{T^n}
 \right].
\end{equation}
The first terms are
\begin{equation}\label{eq:y-over-log-first}
\begin{split}
 y={}&X\Biggl(
 T+\Lambda+\frac{\Lambda}{T}
 -\frac{\Lambda(\Lambda-2)}{2T^2}\\
 &\hspace{5em}+
 \frac{\Lambda(2\Lambda^2-9\Lambda+6)}{6T^3}
 +\cdots\Biggr).
\end{split}
\end{equation}
This famous scale---\(X\log X\), then \(X\log\log X\), then inverse powers of
\(\log X\) with polynomial \(\log\log X\) coefficients---also appears in
asymptotic inversion problems surrounding prime-number estimates.

\begin{intuitionbox}
The operation is still ``polynomial--logarithmic,'' but the correct fast
variable is \(T=\log X\), not \(X\).  Choosing the wrong fast variable makes
the answer look as if it falls outside the method; choosing the correct one
reveals the same Lambert polynomials again.
\end{intuitionbox}

\section{Case study: a logarithmic leading prefactor}
Suppose
\begin{equation}\label{eq:invert-XlogX}
 X=y(\log y)^\beta,
 \qquad \beta\in\RealNumbers.
\end{equation}
A first approximation is
\[
 y_0=X(\log X)^{-\beta}.
\]
But
\[
 \log y_0=\log X-\beta\log\log X,
\]
so even the next correction contains \(L_2=\log\log X\).  Write
\[
 y=X L_1^{-\beta}(1+u),
\]
then
\[
 \log y=L_1-\beta L_2+\log(1+u).
\]
Substituting into \cref{eq:invert-XlogX} gives
\begin{equation}\label{eq:ulog-equation}
 1=(1+u)\left(
 1-\beta\frac{L_2}{L_1}
 +\frac{\log(1+u)}{L_1}
 \right)^\beta.
\end{equation}
Now the small parameter is \(L_1^{-1}\), and the coefficients are polynomials
in \(L_2\).  The first balance is
\begin{equation}\label{eq:u-first-logprefactor}
 u=\beta^2\frac{L_2}{L_1}
 +\BigOAt{L_2^2/L_1^2}{X\to+\infty}.
\end{equation}
Thus
\begin{equation}\label{eq:invert-XlogX-first}
 y=X L_1^{-\beta}
 \left(1+\beta^2\frac{L_2}{L_1}+\cdots\right).
\end{equation}
The appearance of \(L_2\) was predictable from the logarithm of the dominant
prefactor.

\section{Rational logarithmic coefficients versus asymptotic expansion in \texorpdfstring{\(L^{-1}\)}{L inverse}}
A rational function such as
\[
 \frac{L^2+1}{L-1}
\]
can be kept exactly in \(\K(L)\), or expanded for large \(L\):
\[
 \frac{L^2+1}{L-1}
 =L+1+\frac2{L-1}
 =L+1+\frac2L+\frac2{L^2}+\cdots.
\]
These are different representations:
\begin{itemize}
\item the rational form is exact and finite;
\item the \(L^{-1}\) expansion resolves a finer asymptotic scale but is
      infinite.
\end{itemize}
A multiseries implementation should not expand rational functions in \(L^{-1}\)
unless the requested scale requires it.  Premature expansion causes expression
swell and may obscure exact cancellation.

\section{Exponentials and genuinely new monomials}
If a coefficient contains a large nonlinear function of \(L\), exponentiation
leaves every finite logarithmic tower of rational coefficients.  For example,
\[
 \EulerE^{L^2}=X^{\log X},
 \qquad
 \EulerE^{X^{1/2}}.
\]
These are new transmonomials, not coefficient functions.  A full transseries
construction promotes exponentials of purely large transseries to monomials and
then permits well-based sums of them.  Polynomial--logarithmic arithmetic is
best viewed as the height-zero or low-height sector in which these larger
monomials do not yet occur.

\section{Closure tests for the four main operations}
\begin{longtable}{@{}p{0.16\textwidth}p{0.35\textwidth}p{0.39\textwidth}@{}}
\caption{Practical closure tests.}\label{tab:closure-tests}\\
\toprule
Operation & Test & Enlargement when the test fails \\
\midrule
\endfirsthead
\toprule
Operation & Test & Enlargement when the test fails \\
\midrule
\endhead
Multiplication & Are coefficient products in the chosen coefficient ring, and
is exponent support locally finite? & Enlarge the coefficient ring or support
grid. \\
\addlinespace
Division & Is the leading coefficient a unit? & Pass from \(\K[L]\) to
\(\K(L)\), or factor logarithmic monomials explicitly. \\
\addlinespace
Composition & Does \(\log G\) remain in the coefficient algebra, and do
substituted supports remain well based? & Add iterated logs, generalized powers,
or move to a full transseries class. \\
\addlinespace
Reversal & Can the dominant monomial be inverted inside the chosen monomial
group, and is the normalized remainder small? & Add roots, inverse log powers,
iterated logs, or exponential monomials as indicated by dominant balance. \\
\bottomrule
\end{longtable}

\section{The value of a deliberately small class}
A class need not be closed under every conceivable operation to be useful.
The ring \(\K[L]((t))\) has three strong virtues:
\begin{enumerate}
\item every coefficient operation is finite polynomial algebra;
\item t-adic truncation makes recursive algorithms elementary;
\item many important inverse expansions, including Lambert \(W\), stay inside
      it exactly.
\end{enumerate}
The right strategy is to begin in the smallest class that contains the problem,
perform closure tests before each operation, and enlarge only when a specific
new monomial is forced.

\begin{summarybox}
Division, composition, and reversal reveal a hierarchy of necessary
extensions: polynomial coefficients, rational log coefficients, inverse-log
series, iterated logarithms, generalized powers, and finally exponential
transmonomials.  Nested expansions often become ordinary polynomial--logarithmic
series after choosing the correct fast variable.  The inverse of \(y/\log y\),
for example, is controlled by the same Lambert polynomials in the variables
\(T=\log X\) and \(\Lambda=\log T\).
\end{summarybox}

\chapter{Power--log series near zero and Dulac conventions}\label{ch:dulac}

\section{The small-variable dictionary}
Many authors study logarithmic transseries as \(x\to0^+\) rather than
\(X\to+\infty\).  Set
\begin{equation}\label{eq:small-dictionary}
 X=x^{-1},
 \qquad
 \ell=\log\frac1x=\log X=L.
\end{equation}
Then
\[
 X^{-n}P(\log X)=x^nP(\ell).
\]
Thus the integer-exponent core ring becomes
\begin{equation}\label{eq:small-ring}
 \K[\ell]((x)).
\end{equation}
The same object has simply been re-centered at zero.  The notation
\(\ell=\log(1/x)\) is preferable to \(\log x\) on the positive real axis because
\(\ell\to+\infty\).

\section{General power--log and Dulac series}
A common generalized form is
\begin{equation}\label{eq:Dulac-series}
 \DulacFormalSeriesOf{f}(x)=\sum_{j\ge0}x^{\lambda_j}p_j(\ell),
 \qquad
 \lambda_0<\lambda_1<\lambda_2<\cdots\to+\infty,
\end{equation}
where the \(p_j\) are polynomials and the exponent support satisfies suitable
local-finiteness conditions.  Such expansions are called power--log series; in
dynamical contexts, closely related series are called Dulac series.  They occur
as asymptotic expansions of transition and return maps near hyperbolic
polycycles.  More general algebras allow real powers of \(\ell\) and iterated
logs \cite{mardesic-powerlog,mardesic-dulac,peran-logarithmic}.

The dominance order is
\[
 x^{\lambda_1}\ell^{m_1}\succ x^{\lambda_2}\ell^{m_2}
\]
when \(\lambda_1<\lambda_2\), or when the exponents are equal and
\(m_1>m_2\).  As at infinity, one power of \(x\) dominates every finite power
of the logarithm.

\section{Multiplication}
For
\[
 A=\sum_\lambda a_\lambda(\ell)x^\lambda,
 \qquad
 B=\sum_\mu b_\mu(\ell)x^\mu,
\]
the product coefficient at exponent \(\nu\) is
\begin{equation}\label{eq:Dulac-product}
 c_\nu(\ell)=\sum_{\lambda+\mu=\nu}
 a_\lambda(\ell)b_\mu(\ell).
\end{equation}
The support hypotheses must ensure that this sum is finite.  On a finitely
generated exponent grid, the calculation is the same sparse convolution used
in \cref{alg:mul}.

\section{Differentiation near zero}
Since \(\ell=-\log x\),
\[
 \frac{\Differential\ell}{\Differential x}=-\frac1x.
\]
Therefore
\begin{equation}\label{eq:Dulac-derivative}
 \frac{\Differential}{\Differential x}\bigl(P(\ell)x^\lambda\bigr)
 =\bigl(\lambda P(\ell)-P'(\ell)\bigr)x^{\lambda-1}.
\end{equation}
The shifted operator is now \(\lambda-\Dop\), with the opposite sign from the
large-variable convention.  Repeated differentiation gives
\begin{equation}\label{eq:Dulac-repeated-derivative}
 \frac{\Differential^r}{\Differential x^r}\bigl(P(\ell)x^\lambda\bigr)
 =\left(\prod_{j=0}^{r-1}(\lambda-j-\Dop)\right)
 P(\ell)x^{\lambda-r}.
\end{equation}
This formula is indispensable when using Lagrange reversion near zero.

\section{Units and division}
In \(\K[\ell][[x]]\), a series
\[
 A=a_0(\ell)+a_1(\ell)x+\cdots
\]
is a unit only if \(a_0\) is a nonzero constant.  The reciprocal recurrence is
identical to \cref{eq:reciprocal-recurrence}, with \(x\) replacing \(t\).
If \(a_0\) is a nonconstant polynomial in \(\ell\), rational logarithmic
coefficients are required.

For generalized leading exponent \(\lambda_0\), factor
\[
 A=x^{\lambda_0}a_{\lambda_0}(\ell)(1+U).
\]
The power of \(x\) is inverted directly, the logarithmic leading coefficient is
handled in the chosen coefficient field, and \((1+U)^{-1}\) is expanded
geometrically.

\section{Near-identity composition at zero}
Let
\[
 f(y)=y+h(y),
 \qquad h(y)=\LittleO(y),
\]
and
\[
 g(x)=x+H(x),
 \qquad H(x)=\LittleO(x).
\]
Then \(f\compose g\) is obtained by the formal Taylor series
\[
 f(x+H)=\sum_{r\ge0}\frac{f^{(r)}(x)}{r!}H^r.
\]
The smallness condition is now measured relative to \(x\).  For example,
\(H=x^2\ell\) is small because \(x\ell\to0\).

To track the logarithmic variable directly, write
\[
 g=x(1+u),
 \qquad u=H/x=\LittleO(1).
\]
Then
\begin{equation}\label{eq:small-log-compose}
 \log\frac1g
 =\ell-\log(1+u).
\end{equation}
The sign differs from the large-variable formula because
\(\ell=-\log x\).  An outer coefficient \(P(\log(1/y))\) must therefore be
Taylor-expanded at \(\ell\) with increment \(-\log(1+u)\).

\section{A normal-ordered formula near zero}
Suppose
\begin{equation}\label{eq:small-outer}
 f(y)=\sum_n f_n\!\left(\log\frac1y\right)y^n
\end{equation}
and
\[
 g(x)=\lambda x(1+u),
 \qquad \eta=\log(1+u).
\]
Then
\[
 \log\frac1g=\ell-\log\lambda-\eta,
 \qquad
 g^n=\lambda^nx^n\EulerE^{n\eta}.
\]
Consequently,
\begin{equation}\label{eq:small-normal-ordered}
 f\compose g
 =\sum_n\lambda^nx^n
 \mathopen{:}\EulerE^{-\eta(\Dop-n)}\mathclose{:}
 f_n(\ell-\log\lambda),
\end{equation}
where normal ordering again means that powers of \(\eta\) are not
differentiated.  This is the zero-centered counterpart of
\cref{eq:regular-composition-formula}.

\section{Series reversal near zero}
The Lagrange--Bürmann formula has exactly the same form:
\begin{equation}\label{eq:small-Lagrange}
 \rev{(x+h)}(x)
 =x+\sum_{m\ge1}\frac{(-1)^m}{m!}
 \left(\frac{\Differential}{\Differential x}\right)^{m-1}h(x)^m,
\end{equation}
provided \(h=\LittleO(x)\) and the support conditions make each truncation finite.
The differential rule \cref{eq:Dulac-repeated-derivative} turns every term into
a finite polynomial calculation.

\begin{workedbox}[title={Reversing \(x+a x^2\log(1/x)\)}]
Let
\[
 f(x)=x+a x^2\ell,
 \qquad \ell=\log(1/x).
\]
Here \(h(x)=a x^2\ell\).  The first Lagrange term is
\[
 -h=-a x^2\ell.
\]
The second is
\[
 \frac12\frac{\Differential}{\Differential x}h^2
 =\frac{a^2}{2}\frac{\Differential}{\Differential x}(x^4\ell^2)
 =a^2x^3(2\ell^2-\ell).
\]
The third term begins at x-order four.  Hence
\begin{equation}\label{eq:small-inverse-example}
 \rev f(x)=x-a x^2\ell+a^2x^3(2\ell^2-\ell)
 +\FormalBigO_x(x^4).
\end{equation}
Substitution into \(f\) cancels the x-squared and x-cubed blocks.
\end{workedbox}

\section{Composition can create an iterated logarithm}
Let
\[
 g(x)=x^p\ell^\beta.
\]
Then
\begin{equation}\label{eq:small-nested-log}
 \log\frac1{g(x)}
 =p\ell-\beta\log\ell.
\end{equation}
Thus an outer polynomial in \(\log(1/y)\) becomes a polynomial in \(\ell\) and
\(\log\ell\).  The same closure boundary seen at infinity reappears.  In Dulac
normal-form problems, iterated logarithms can be forced by repeated
normalization and inversion; modern logarithmic transseries frameworks are
designed to accommodate them.

\section{Parabolic and hyperbolic leading behavior}
A map
\[
 f(x)=x+\LittleO(x)
\]
is called tangent to the identity or parabolic in many dynamical settings.  Its
inverse is again tangent to the identity, and the triangular reversion method
applies.  A map
\[
 f(x)=\lambda x+\LittleO(x),
 \qquad \lambda\ne0,
\]
is first normalized by its linear factor.  A map
\[
 f(x)=x^\alpha+\text{higher terms},
 \qquad \alpha\ne1,
\]
has strongly non-tangent leading behavior; inversion or conjugacy may require
fractional exponents and iterated logs.  Formal normalization results for such
classes are an active part of the theory of logarithmic transseries
\cite{mardesic-powerlog,peran-logarithmic}.

\section{Translation errors to avoid}
When moving between \(X\to\infty\) and \(x\to0^+\):
\begin{enumerate}
\item remember that \(X^{-n}=x^n\);
\item use \(\ell=\log(1/x)=\log X\), not \(\log x\), unless all signs are
      changed consistently;
\item replace the shifted derivative \(\Dop-n\) by \(n-\Dop\) for
      \(x^nP(\ell)\);
\item reverse the sign of the logarithmic increment in composition, as in
      \cref{eq:small-log-compose};
\item translate analytic sectors: a sector at infinity maps to a punctured
      sector at zero under reciprocal coordinates.
\end{enumerate}

\begin{summarybox}
The large-variable and small-variable theories are equivalent under
\(X=1/x\) and \(\ell=\log(1/x)\).  Near zero, differentiation acts by
\(\lambda-\Dop\) on a block \(x^\lambda P(\ell)\), and composition changes
\(\ell\) by \(-\log(1+u)\).  Lagrange reversion remains valid and gives explicit
power--log inverses.  General Dulac supports require the same local-finiteness
care as generalized series at infinity, and composition with logarithmic
prefactors can force iterated logs.
\end{summarybox}

\chapter{Algorithms and implementation}\label{ch:implementation}

\section{A canonical sparse representation}
For a fixed coefficient domain, represent
\[
 A=\sum_na_n(L)t^n
\]
as a finite map
\[
 n\longmapsto a_n(L)
\]
containing only the blocks needed below the current cutoff.  For example,
\[
 X-L+\frac{L}{X}+\frac{L(L-2)}{2X^2}
\]
is stored as
\[
 \{-1\mapsto1,\ 0\mapsto-L,\ 1\mapsto L,\
 2\mapsto L(L-2)/2\}.
\]
A sparse map is preferable to a dense list when generalized exponents or large
gaps occur.  For the integer t-adic core, either representation works.

Every constructor and arithmetic operation should enforce four invariants:
\begin{enumerate}
\item coefficient expressions are normalized in the chosen domain;
\item zero coefficients are removed;
\item no exponent at or beyond the active cutoff is stored;
\item the coefficient ring is recorded, so an unintended transition from
      \(\K[L]\) to \(\K(L)\) is visible.
\end{enumerate}

\section{Precision as part of the type}
A truncated formal series is not merely an expression; it is an expression plus
a guarantee:
\[
 A_N=[A]_{<N},
 \qquad A=A_N+\FormalBigOAt{t^N}{t}.
\]
Every operation should compute the strongest guaranteed output precision from
the input precisions.  If both inputs are known modulo \(t^N\) and have
nonnegative valuation, their sum and product are known modulo \(t^N\).  If
negative exponents occur, product precision requires care because an unknown
high block of one factor can combine with a negative block of the other.

A safe implementation either:
\begin{itemize}
\item stores an explicit lower valuation and requests enough guard terms from
      each operand; or
\item uses a lazy series interface that computes coefficients on demand.
\end{itemize}
In the common near-identity calculations, the negative part is fixed and
finite---usually only the block \(t^{-1}\)---so guard precision is easy to
calculate.

\section{Core pseudocode}
The following routines assume integer t-exponents and exact coefficient
arithmetic.

\subsection{Multiplication}
\begin{lstlisting}[style=pseudostyle]
function multiply(A, B, N):
    C := empty map
    for (i, ai) in A:
        for (j, bj) in B:
            if i + j < N:
                C[i+j] := C.get(i+j, 0) + ai*bj
    normalize every coefficient of C
    delete zero blocks
    return C
\end{lstlisting}

\subsection{Reciprocal of a unit}
\begin{lstlisting}[style=pseudostyle]
function reciprocal(A, N):
    require A has valuation 0
    a0 := A[0]
    require a0 is a unit in the coefficient ring
    B[0] := 1/a0
    for n from 1 to N-1:
        s := 0
        for k from 1 to n:
            s := s + A.get(k,0)*B.get(n-k,0)
        B[n] := -s/a0
    return normalized B
\end{lstlisting}

\subsection{Formal logarithm}
\begin{lstlisting}[style=pseudostyle]
function log1p(U, N):
    require valuation(U) >= 1
    S := 0
    P := 1
    for m from 1 to N-1:
        P := truncate(P*U, N)
        S := S + (-1)^(m+1)*P/m
    return truncate(S, N)
\end{lstlisting}

\subsection{Differentiation}
\begin{lstlisting}[style=pseudostyle]
function derivative_X(A, N):
    B := empty map
    for (n, p(L)) in A:
        if n+1 < N:
            B[n+1] := derivative_L(p) - n*p
    return normalized B
\end{lstlisting}

\section{Normal-ordered composition as an algorithm}
For \(G=\lambda X(1+u)\), first compute
\(\delta=\log(1+u)\).  Then process each outer block independently.

\begin{lstlisting}[style=pseudostyle]
function compose_regular(F, lambda, u, N):
    delta := log1p(u, required_guard_precision)
    H := 0
    for (n, f_n(L)) in F that can contribute below N:
        phi := f_n(L + log(lambda))
        jet := phi
        delta_power := 1
        for r from 0 while n+r < N:
            term := lambda^(-n) * t^n
                    * delta_power * jet / r!
            H := truncate(H + term, N)
            delta_power := truncate(delta_power*delta, N-n)
            jet := derivative_L(jet) - n*jet
    return H
\end{lstlisting}
The variable \texttt{jet} stores \((\Dop-n)^r\phi\).  The algorithm never
constructs an ambiguous differential-operator exponential.

For a near-identity inner series
\[
 G=X+A_0+A_1t+A_2t^2+\cdots,
\]
form
\begin{equation}\label{eq:implementation-u}
 u=qG-1=A_0t+A_1t^2+A_2t^3+\cdots
\end{equation}
and call the regular composition routine with \(\lambda=1\).

\section{Triangular and Newton reversion}
The triangular algorithm is almost definitionally simple:
\begin{lstlisting}[style=pseudostyle]
function reverse_triangular(F, N):
    G := X
    for n from 0 to N-1:
        R := truncate(compose(F,G) - X, n+1)
        r_n := coefficient(R, t^n)
        G := G - r_n*t^n
    return G
\end{lstlisting}
At the end of the n-th pass, the residual vanishes below t-order \(n+1\).

A Newton implementation uses
\begin{lstlisting}[style=pseudostyle]
function reverse_newton(F, G, N):
    R  := truncate(compose(F,G) - X, N)
    J  := truncate(compose(derivative_X(F),G), N)
    dG := truncate(R / J, N)
    return truncate(G - dG, N)
\end{lstlisting}
The current approximation \(G\) should already be correct to positive t-order.
If it is correct modulo \(t^M\), request the new result modulo
\(t^{\min(2M,N)}\).

\section{Coefficient-domain design}
The series layer should be parameterized by its coefficient domain.  Useful
choices include:
\begin{enumerate}
\item exact polynomials \(\K[L]\);
\item rational functions \(\K(L)\);
\item polynomials or rational functions in \(L_1,\ldots,L_r\);
\item sparse Laurent monomials in the logarithmic tower;
\item symbolic expressions with a canonical simplifier.
\end{enumerate}
The first two choices are especially robust.  A completely general symbolic
expression domain is convenient but can hide zero tests, introduce accidental
branch transformations, and cause expression swell.

A good interface asks the coefficient domain to provide:
\[
 +,\quad -,\quad \times,\quad \text{unit test},\quad \text{inverse of a unit},
 \quad \Dop,\quad \text{substitution }L\mapsto L+c,
\]
and a reliable zero test.

\section{Bell-polynomial acceleration}
The direct composition algorithm builds powers \(\delta^r\).  If many outer
blocks share the same inner argument, precompute
\[
 C_{m,r}=[t^m]\frac{\delta^r}{r!}
 =\frac1{m!}\ExponentialPartialBellPolynomial{m}{r}
 (1!\delta_1,2!\delta_2,\ldots).
\]
Then \cref{eq:composition-bell} becomes a matrix-like contraction between the
Bell table and shifted derivative jets.  This avoids recomputing the same
powers for every outer block.

For moderate \(N\), repeated truncated multiplication is simpler.  For large
\(N\), Bell recurrences, relaxed series, or fast composition algorithms become
attractive.

\section{A tested SymPy prototype}
The following compact implementation covers sparse addition, multiplication,
reciprocal formation, \(\log(1+u)\), differentiation, regular composition,
near-identity composition, triangular reversal, and Lambert polynomial
generation.  The exponent \(n\) in the dictionary represents
\(t^n=X^{-n}\).  It is intended as transparent reference code rather than a
high-performance library.

\begin{lstlisting}[style=pythonstyle]
from __future__ import annotations
from typing import Dict
import sympy as sp

L = sp.Symbol("L")
Series = Dict[int, sp.Expr]


def clean(a: Series, cutoff: int | None = None) -> Series:
    out: Series = {}
    for n, p in a.items():
        if cutoff is not None and n >= cutoff:
            continue
        p = sp.cancel(sp.expand(p))
        if p != 0:
            out[n] = p
    return out


def add(a: Series, b: Series, cutoff: int) -> Series:
    out = dict(a)
    for n, p in b.items():
        if n < cutoff:
            out[n] = out.get(n, sp.S.Zero) + p
    return clean(out, cutoff)


def scale(a: Series, c: sp.Expr, cutoff: int) -> Series:
    return clean({n: c*p for n, p in a.items()}, cutoff)


def shift_t(a: Series, s: int, cutoff: int) -> Series:
    return clean({n+s: p for n, p in a.items()
                  if n+s < cutoff}, cutoff)


def mul(a: Series, b: Series, cutoff: int) -> Series:
    out: Series = {}
    for i, ai in a.items():
        for j, bj in b.items():
            n = i + j
            if n < cutoff:
                out[n] = out.get(n, sp.S.Zero) + ai*bj
    return clean(out, cutoff)


def reciprocal_unit(a: Series, cutoff: int) -> Series:
    a0 = sp.cancel(a.get(0, sp.S.Zero))
    if a0 == 0:
        raise ValueError("t^0 coefficient must be nonzero")
    b: Series = {0: 1/a0}
    for n in range(1, cutoff):
        s = sum((a.get(k, 0)*b.get(n-k, 0)
                 for k in range(1, n+1)), sp.S.Zero)
        b[n] = -s/a0
    return clean(b, cutoff)


def log1p(u: Series, cutoff: int) -> Series:
    if u and min(u) <= 0:
        raise ValueError("u must have positive t-valuation")
    out: Series = {}
    power: Series = {0: sp.S.One}
    for m in range(1, cutoff):
        power = mul(power, u, cutoff)
        c = sp.Rational((-1)**(m+1), m)
        out = add(out, scale(power, c, cutoff), cutoff)
    return clean(out, cutoff)


def derivative_x(a: Series, cutoff: int) -> Series:
    out: Series = {}
    for n, p in a.items():
        if n+1 < cutoff:
            out[n+1] = sp.diff(p, L) - n*p
    return clean(out, cutoff)


def compose_regular(f: Series, u: Series,
                    cutoff: int, lam=sp.S.One) -> Series:
    """Truncate F(lam*X*(1+u)) below t**cutoff."""
    delta = log1p(u, cutoff + 2)
    out: Series = {}
    for n, poly in f.items():
        poly = sp.sympify(poly)
        phi = sp.expand(poly.subs(L, L + sp.log(lam)))
        jet = phi
        delta_power: Series = {0: sp.S.One}
        for r in range(max(0, cutoff-n)):
            coeffs = {j: p*jet for j, p in delta_power.items()}
            term = shift_t(coeffs, n, cutoff)
            term = scale(term,
                         lam**(-n)/sp.factorial(r), cutoff)
            out = add(out, term, cutoff)
            delta_power = mul(delta_power, delta, cutoff-n)
            jet = sp.diff(jet, L) - n*jet
    return clean(out, cutoff)


def near_identity_u(g: Series, cutoff: int) -> Series:
    """For G=X+..., return u=G/X-1=t*G-1."""
    u = shift_t(g, 1, cutoff + 1)
    u[0] = u.get(0, 0) - 1
    return clean(u, cutoff + 1)


def compose_near_identity(f: Series, g: Series,
                          cutoff: int) -> Series:
    return compose_regular(f, near_identity_u(g, cutoff),
                           cutoff)


def reverse_triangular(f: Series, cutoff: int) -> Series:
    identity: Series = {-1: sp.S.One}
    g: Series = dict(identity)
    for n in range(0, cutoff):
        fg = compose_near_identity(f, g, n+1)
        residual = add(fg, scale(identity, -1, n+1), n+1)
        g[n] = sp.expand(g.get(n, 0)-residual.get(n, 0))
        g = clean(g, cutoff)
    return g


def lambert_polynomial(n: int) -> sp.Expr:
    S1 = sp.functions.combinatorial.numbers.stirling
    return sp.expand(sum(
        (-1)**(n+m)*S1(n, n-m+1, kind=1, signed=False)
        * L**m/sp.factorial(m)
        for m in range(1, n+1)
    ))


# Formal self-test: inverse of F(X)=X+log X.
F: Series = {-1: 1, 0: L}
G = reverse_triangular(F, 7)
expected: Series = {-1: 1, 0: -L}
expected.update({n: lambert_polynomial(n)
                 for n in range(1, 7)})
assert clean(G, 7) == clean(expected, 7)
assert add(compose_near_identity(F, G, 7), {-1: -1}, 7) == {}
assert add(compose_near_identity(G, F, 7), {-1: -1}, 7) == {}
\end{lstlisting}
The final assertions check the Stirling formula and both composition residuals
through t-order six.

\section{Testing strategy}
Formal-series software should be tested by identities, not only by matching a
few printed coefficients.

\subsection{Ring identities}
For random low-degree polynomial blocks and safe cutoffs, verify
\[
 A(B+C)=AB+AC,
 \qquad
 (AB)C=A(BC).
\]
Use exact coefficients so every comparison is symbolic.

\subsection{Reciprocal and quotient residuals}
Verify
\[
 A\recip{A}=1+\FormalBigOAt{t^N}{t},
 \qquad
 A(C/A)=C+\FormalBigOAt{t^N}{t}.
\]
Include tests whose leading coefficient is a nonconstant polynomial and confirm
that the polynomial-only domain rejects them while the rational domain accepts
them.

\subsection{Composition identities}
Check
\[
 F\compose X=F,
 \qquad
 X\compose G=G,
\]
\[
 (F\compose G)\compose H=F\compose(G\compose H),
\]
and the chain rule \cref{eq:formal-chain-rule}.  Associativity tests are
particularly good at detecting missing guard terms.

\subsection{Reversal identities}
For a random near-identity series, compute \(G=\rev F\) and check both residuals
in \cref{eq:two-residuals}.  Then perturb the highest known coefficient and
confirm that the residual changes at exactly the predicted order.

\subsection{Lambert regression tests}
Generate \(\LambertPolynomial{n}\) by three independent methods:
\begin{enumerate}
\item the fixed-point recurrence \cref{eq:Lambert-triangular-recurrence};
\item the operator formula \cref{eq:Pn-operator};
\item the Stirling sum \cref{eq:Lambert-Stirling}.
\end{enumerate}
Agreement gives strong protection against sign and indexing errors.

\section{Expression swell and normalization policy}
Expanding every polynomial product immediately makes zero detection reliable
but may create large intermediate expressions.  Factoring every expression
keeps some patterns compact but makes addition and zero tests expensive.  A
balanced policy is:
\begin{enumerate}
\item keep coefficient polynomials expanded and collected in \(L\);
\item keep external parameters in rational factored form when possible;
\item cancel rational functions after division but avoid unnecessary partial
      fractions;
\item truncate before expensive normalization;
\item expose optional factorization only for presentation.
\end{enumerate}
The printed Lambert polynomials, for example, are best presented with a common
factor \(L\), but an implementation may store them expanded.

\section{High-order improvements}
For hundreds or thousands of t-orders, the following improvements matter:
\begin{itemize}
\item Newton reciprocal and reversion instead of coefficientwise recurrences;
\item FFT-based multiplication in t and in \(L\);
\item relaxed or online power-series algorithms;
\item baby-step/giant-step composition;
\item memoization of shifted derivative jets;
\item Bell-polynomial or exponential-generating-function recurrences;
\item modular arithmetic and rational reconstruction for exact integer
      coefficient problems.
\end{itemize}
The foundational formulas do not change.  Fast algorithms reorganize the same
finite coefficient dependencies.

\section{Reproducibility record}
For a serious symbolic computation, record:
\begin{enumerate}
\item the formal variable, fast scale, and logarithmic tower;
\item the coefficient domain and branch conventions;
\item the support grid and truncation order;
\item whether the truncation is blockwise in t or monomialwise in t and \(L\);
\item the residual identity used for verification;
\item the exact software version and simplification assumptions.
\end{enumerate}
This information is often more important than the final list of coefficients,
because it determines whether the displayed remainder is a formal valuation
tail or an analytic estimate with an explicit limiting regime.

\begin{summarybox}
A polynomial--logarithmic series can be implemented as a sparse map from
t-exponents to normalized coefficient polynomials.  Precision and coefficient
domain must be explicit.  Multiplication, reciprocal, logarithm,
differentiation, normal-ordered composition, and triangular reversal all reduce
to finite loops.  Newton methods accelerate high-order work, while residual
identities and independent Lambert polynomial formulas provide strong tests.
\end{summarybox}

\chapter{Extended worked calculations}\label{ch:worked}

\section{A complete quotient calculation}
Let
\begin{align*}
 N={}&1+(2L-1)t+(L^2+3)t^2+(L^3-L)t^3+\FormalBigOAt{t^4}{t},\\
 D={}&1+(L+1)t+(2L-1)t^2+L^2t^3+\FormalBigOAt{t^4}{t}.
\end{align*}
Seek \(Q=N/D=\sum_{n=0}^3d_nt^n+\FormalBigOAt{t^4}{t}\).  Since the leading coefficient of
\(D\) is one, \cref{eq:quotient-recurrence} becomes
\[
 d_n=N_n-\sum_{k=1}^nD_kd_{n-k}.
\]
Thus
\[
 d_0=1,
\]
\[
 d_1=(2L-1)-(L+1)=L-2,
\]
\begin{align*}
 d_2&=(L^2+3)-(L+1)(L-2)-(2L-1)\\
 &=6-L,
\end{align*}
and
\begin{align*}
 d_3={}&(L^3-L)-(L+1)(6-L)\\
 &{}-(2L-1)(L-2)-L^2\\
={}&L^3-2L^2-L-8.
\end{align*}
Therefore
\begin{equation}\label{eq:worked-quotient}
 \frac ND=1+(L-2)t+(6-L)t^2
 +(L^3-2L^2-L-8)t^3+\FormalBigOAt{t^4}{t}.
\end{equation}
Multiplying \cref{eq:worked-quotient} by \(D\) reproduces \(N\) through
t-order three.

\section{Logarithm followed by exponential}
Take
\[
 A=1+Lq+(2-L)t^2+(L^2+1)t^3+\FormalBigOAt{t^4}{t}.
\]
Set \(U=A-1\).  Since
\[
 \log(1+U)=U-\frac{U^2}{2}+\frac{U^3}{3}+\FormalBigOAt{t^4}{t},
\]
we find
\begin{equation}\label{eq:worked-log}
\begin{split}
 \log A={}&Lt+
 \left(2-L-\frac{L^2}{2}\right)t^2\\
 &+\left(2L^2-L-2+\frac{L^3}{3}\right)t^3
 +\FormalBigOAt{t^4}{t}.
\end{split}
\end{equation}
Let the right side be \(V\).  Expanding
\[
 \EulerE^V=1+V+\frac{V^2}{2}+\frac{V^3}{6}+\FormalBigOAt{t^4}{t}
\]
returns the original \(A\).  This round-trip test simultaneously checks
multiplication, truncation, logarithm, and exponential routines.

\section{Composition is not commutative}
Let
\[
 F(X)=X+a(L),
 \qquad
 G(X)=X+A(L).
\]
From \cref{eq:h0,eq:h1},
\[
 F\compose G
 =X+(a+A)+Aa'X^{-1}+\FormalBigOAt{t^2}{t},
\]
while
\[
 G\compose F
 =X+(a+A)+aA'X^{-1}+\FormalBigOAt{t^2}{t}.
\]
Hence the first commutator block is
\begin{equation}\label{eq:composition-commutator}
 F\compose G-G\compose F
 =(Aa'-aA')X^{-1}+\FormalBigOAt{t^2}{t}.
\end{equation}
For \(a=L^2\) and \(A=L\), this becomes
\[
 F\compose G-G\compose F=L^2X^{-1}+\FormalBigOAt{t^2}{t}.
\]
The noncommutativity is created entirely by the change in the logarithmic
argument.

\begin{remark}
Expression~\eqref{eq:composition-commutator} resembles the Lie bracket of the
vector fields \(a(L)\partial_L\) and \(A(L)\partial_L\).  This is not an
accident: groups of near-identity substitutions have infinitesimal Lie
algebras, and logarithmic coefficient derivatives encode the first nonlinear
interaction.
\end{remark}

\section{A four-block composition with nontrivial coefficients}
Let
\[
 F(Y)=Y+(\log^2Y-1)+\frac{2\log Y}{Y}+\frac1{Y^2}
 +\FormalBigOAt{Y^{-3}}{Y^{-1}}
\]
and
\[
 G(X)=X+(1-L)+\frac{L^2}{X}+\FormalBigOAt{t^2}{t}.
\]
Thus
\[
 a=L^2-1,
 \quad b=2L,
 \quad c=1,
 \qquad
 A=1-L,
 \quad B=L^2,
 \quad C=D=0.
\]
Using \cref{thm:near-composition-coeffs},
\begin{align*}
 h_0&=L^2-L,\\
 h_1&=-L^2+4L,\\
 h_2&=L^3+5L^2-7L+4,\\
 h_3&=\frac43L^4-L^3-8L^2+\frac{41}{3}L-6.
\end{align*}
Therefore
\begin{equation}\label{eq:worked-composition-four}
\begin{split}
 F(G(X))={}&X+(L^2-L)+\frac{-L^2+4L}{X}\\
 &+\frac{L^3+5L^2-7L+4}{X^2}\\
 &+\frac{\tfrac43L^4-L^3-8L^2+\tfrac{41}{3}L-6}{X^3}
 +\FormalBigOAt{t^4}{t}.
\end{split}
\end{equation}
A direct expansion of \(\log G\), \(G^{-1}\), and \(G^{-2}\) gives the same
answer but requires substantially more intermediate algebra.

\section{Reversing a perturbed logarithmic map}
Consider
\begin{equation}\label{eq:worked-F-perturbed}
 F(X)=X+L+\frac{L^2+1}{X}+\FormalBigOAt{t^2}{t}.
\end{equation}
Here
\[
 a=L,
 \qquad b=L^2+1,
 \qquad c=d=0.
\]
Substitution into \cref{thm:inverse-coefficients} gives
\begin{align*}
 A={}&-L,\\
 B={}&-L^2+L-1,\\
 C={}&-L^3+\frac72L^2-2L+1,\\
 D={}&-2L^4+\frac{25}{3}L^3-\frac{21}{2}L^2+6L-2.
\end{align*}
Hence
\begin{equation}\label{eq:worked-inverse-perturbed}
\begin{split}
 \rev F(X)={}&X-L+\frac{-L^2+L-1}{X}\\
 &+\frac{-L^3+\tfrac72L^2-2L+1}{X^2}\\
 &+\frac{-2L^4+\tfrac{25}{3}L^3-\tfrac{21}{2}L^2+6L-2}{X^3}
 +\FormalBigOAt{t^4}{t}.
\end{split}
\end{equation}
This example shows that the Lambert polynomials are not a universal answer by
themselves: a lower-order perturbation of the original map couples into every
subsequent inverse coefficient.

\section{A constant and logarithmic shift}
Let
\begin{equation}\label{eq:constant-log-shift}
 F(Y)=Y+c+\alpha\log Y.
\end{equation}
Set
\[
 M=L+\frac c\alpha
\]
when \(\alpha\ne0\).  Since \(c+\alpha L=\alpha M\) and
\(\Differential M/\Differential X=1/X\), the Lagrange formula gives
\begin{equation}\label{eq:constant-log-shift-inverse}
 \rev F(X)
 =X-\alpha M+
 \sum_{n\ge1}\frac{\alpha^{n+1}\LambertPolynomial{n}(M)}{X^n}.
\end{equation}
Equivalently,
\[
 \rev F(X)=X-c-\alpha L
 +\frac{\alpha^2(L+c/\alpha)}X+\cdots.
\]
A constant in the original perturbation shifts the logarithmic polynomial
argument in every inverse coefficient.

\section{Reversing \texorpdfstring{\(X-\log X\)}{X-log X}}
The inverse of
\[
 F(Y)=Y-\log Y
\]
solves \(u-\log u=X\), the real equation behind \(W_{-1}\).  Put
\(\alpha=-1\) in \cref{eq:alpha-expansion}:
\begin{equation}\label{eq:inverse-X-minus-logX}
\begin{split}
 \rev F(X)={}&X+L+\frac{L}{X}
 -\frac{L(L-2)}{2X^2}\\
 &+\frac{L(2L^2-9L+6)}{6X^3}
 -\frac{L(3L^3-22L^2+36L-12)}{12X^4}+\cdots.
\end{split}
\end{equation}
Multiplying by \(-1\) and setting \(X=T\) recovers the real
\(W_{-1}\) expansion in \cref{eq:Wm1-first}.

\section{Newton reversion in action}
For \(F(Y)=Y+\log Y\), start with
\[
 G_1=X-L+\frac{L}{X}.
\]
As calculated earlier,
\[
 F(G_1)-X=\left(L-\frac{L^2}{2}\right)X^{-2}+\FormalBigOAt{t^3}{t}.
\]
Also
\[
 F'(G_1)=1+\frac1{G_1}=1+X^{-1}+\FormalBigOAt{t^2}{t}.
\]
The Newton correction is therefore
\begin{align*}
 -\frac{F(G_1)-X}{F'(G_1)}
 &=-\left(L-\frac{L^2}{2}\right)X^{-2}+\FormalBigOAt{t^3}{t}\\
 &=\frac{L(L-2)}{2X^2}+\FormalBigOAt{t^3}{t},
\end{align*}
which inserts \(\LambertPolynomial{2}/X^2\).  If the residual and denominator are carried to
higher order before division, the same Newton step determines several further
blocks simultaneously.

\section{A numerical inverse check}
Take
\[
 F(y)=y+\log y+\frac{1}{y}
\]
and solve \(F(y)=100\).  The expansion
\cref{eq:perturbed-Lambert-inverse} with \(\alpha=\beta=1\), \(\gamma=0\)
gives
\[
 y\approx X-L+\frac{L-1}{X}
 +\frac{\tfrac12L^2-2L+1}{X^2},
 \qquad X=100,
 \quad L=\log100.
\]
Numerically this gives approximately
\[
 y\approx95.43112.
\]
Substituting back produces a residual of order \(10^{-5}\); one Newton
correction for the original scalar equation reduces it quadratically.  The
formal expansion supplies a high-quality branch-aware starting value, while
ordinary numerical iteration supplies final precision.

\section{A small-variable reversal checked by composition}
Return to
\[
 f(x)=x+a x^2\ell,
 \qquad
 g(x)=x-a x^2\ell+a^2x^3(2\ell^2-\ell).
\]
To verify \(f(g(x))\), first note
\[
 \log\frac1{g(x)}
 =\ell-\log\left(1-a x\ell+a^2x^2(2\ell^2-\ell)\right).
\]
Through x-order one inside the logarithm,
\[
 \log\frac1g=\ell+a x\ell+\FormalBigO_x(x^2).
\]
Also
\[
 g^2=x^2-2a x^3\ell+\FormalBigO_x(x^4).
\]
Therefore
\begin{align*}
 f(g)&=g+a g^2\log(1/g)\\
 &=x-a x^2\ell+a^2x^3(2\ell^2-\ell)
   +a x^2\ell+a^2x^3(\ell-2\ell^2)
   +\FormalBigO_x(x^4)\\
 &=x+\FormalBigO_x(x^4).
\end{align*}
Every displayed correction cancels.

\section{How to organize a hand calculation}
For a nontrivial composition or reversal, use the following worksheet.
\begin{enumerate}
\item Write the target precision and the exact meaning of the remainder.
\item Factor the inner series by its dominant monomial.
\item Compute only the needed terms of \(\delta=\log(1+u)\).
\item Make a dependency table: which outer blocks and derivative orders can
      reach each t-order?
\item Keep each t-block as a polynomial in \(L\) until the end.
\item Normalize and collect after every t-order, not after every elementary
      multiplication.
\item Verify the final residual before simplifying the presentation.
\end{enumerate}
This method is less error-prone than expanding one enormous symbolic expression
and asking a computer algebra system to sort it afterward.

\begin{summarybox}
The examples in this chapter illustrate the full workflow: quotient recurrence,
logarithm/exponential round trips, noncommutative composition, universal
composition coefficients, perturbed inverse formulas, Lambert-type shifts,
Newton correction, numerical refinement, and small-variable reversion.
Residual checks turn each printed expansion into a verifiable calculation.
\end{summarybox}

\chapter{From formal series to analytic asymptotics}\label{ch:analytic}

\section{Two questions that should not be conflated}
A formal calculation answers:
\begin{quote}
If an expansion of the chosen form exists, what must its coefficients be?
\end{quote}
An analytic theorem answers:
\begin{quote}
Does an actual function have that expansion, on which domain, with what
remainder, and on which branch?
\end{quote}
Formal algebra often determines coefficients uniquely and efficiently, but it
does not by itself prove existence, convergence, or branch validity.  This
chapter explains the standard bridges in the polynomial--logarithmic setting.

\section{An ordered asymptotic scale}
On the positive real axis, the monomials
\[
 X^{-n}L^m,
 \qquad n\in\IntegerNumbers,
 \quad m\in\NaturalNumbers,
\]
form an asymptotic scale under the lexicographic order
\cref{eq:lex-order}.  More explicitly,
\[
 X^{-n}L^{m-1}=\LittleO(X^{-n}L^m),
\]
and, for every fixed \(r,s\),
\[
 X^{-(n+1)}L^r=\LittleO(X^{-n}L^s).
\]
The second fact follows from
\(L^{r-s}/X\to0\).  Thus a finite polynomial t-block is asymptotically
well separated from the next t-block.

If logarithmic degrees grow with t-order, a blockwise statement such as
\cref{eq:analytic-remainder} is usually clearer than enumerating all monomials.
For sectorial complex asymptotics, replace real inequalities by uniform bounds
on a sector where \(\log X\) is single-valued.

\section{Addition and multiplication of actual expansions}
Suppose
\[
 f=A_N+R_f,
 \qquad
 g=B_N+R_g,
\]
where \(A_N,B_N\) are finite polynomial--logarithmic sums and the remainders
are controlled at the next t-order.  Then
\[
 fg=A_NB_N+A_NR_g+B_NR_f+R_fR_g.
\]
Polynomial--logarithmic factors grow more slowly than any positive power of
\(X\), so the cross terms remain at the predicted lower order after allowing a
possibly larger logarithmic power in the bound.  This justifies Cauchy
multiplication of Poincaré-type expansions.

The formal coefficient convolution is therefore not merely symbolic: under
standard uniform remainder hypotheses, it gives the analytic product
expansion.

\section{Reciprocals and quotients}
Suppose
\[
 f(X)=a_0+\LittleO(1),
 \qquad a_0\ne0.
\]
Then \(f\) is bounded away from zero for sufficiently large \(X\), and
\[
 \frac1f=\frac1{a_0}\frac1{1+(f-a_0)/a_0}.
\]
The analytic geometric series is valid once
\(|(f-a_0)/a_0|<1\), and its coefficient expansion agrees with the formal
reciprocal recurrence.  If the leading coefficient is \(a_0(L)\), one needs a
region in which that function is nonzero and a coefficient class that can
represent \(1/a_0(L)\).

A formal unit test thus mirrors the analytic condition: the leading block must
not vanish in the asymptotic region.

\section{Logarithms and powers}
If \(u(X)=\LittleO(1)\), then for sufficiently large positive \(X\), or uniformly in
a suitable complex sector,
\[
 \log(1+u)=u-\frac{u^2}{2}+\frac{u^3}{3}-\cdots.
\]
The analytic Taylor remainder after \(M\) terms is
\(\BigOAt{u^{M+1}}{u\to0}\) when \(u\) stays in a compact subset avoiding the branch cut and
\(-1\).  Since positive t-valuation implies \(u^M\) gains at least M t-orders,
formal truncation and analytic Taylor truncation agree.

For \(A=cX^\rho(1+u)\), branch choices enter through \(\log c\),
\(X^\rho\), and \(\log X\).  The small-unit coefficients are universal once
those choices are fixed.

\section{Composition with a regular inner function}
Let \(G(X)=\lambda X(1+u(X))\), where \(u=\LittleO(1)\), and suppose the outer
function has a uniform expansion in a region containing the image of \(G\).
Then
\[
 \log G=L+\log\lambda+\log(1+u)
\]
and
\[
 G^{-n}=\lambda^{-n}X^{-n}(1+u)^{-n}
\]
have ordinary analytic Taylor expansions.  Substituting them into finitely many
outer blocks and estimating the outer remainder proves the corresponding
finite part of \cref{eq:regular-composition-formula}.

The important analytic hypotheses are geometric:
\begin{enumerate}
\item \(G(X)\) must tend to infinity in the outer function's asymptotic
      domain;
\item logarithm branches must remain consistent;
\item the perturbation \(u\) must be uniformly small;
\item the outer remainder estimate must remain uniform after substitution.
\end{enumerate}
Formal support control alone does not provide these domain facts.

\section{Residual-to-error transfer for inverse problems}
A formal approximation is especially valuable when its functional residual can
be converted into a rigorous error bound.

\begin{proposition}[Real residual bound]\label{prop:residual-error}
Let \(F\) be continuously differentiable and strictly monotone on an interval
containing an exact solution \(y_*\) of \(F(y_*)=X\) and an approximation
\(y_N\).  If
\[
 \AbsoluteValue{F'(y)}\ge m>0
\]
throughout the interval between \(y_N\) and \(y_*\), then
\begin{equation}\label{eq:residual-error-bound}
 \AbsoluteValue{y_N-y_*}\le\frac{\AbsoluteValue{F(y_N)-X}}{m}.
\end{equation}
\end{proposition}
\begin{proof}
The mean value theorem gives
\[
 F(y_N)-F(y_*)=F'(\xi)(y_N-y_*)
\]
for some intermediate \(\xi\).  Take absolute values and use the lower bound on
\(\AbsoluteValue{F'}\).
\end{proof}

For \(F(y)=y+\log y\) on \(y>0\),
\[
 F'(y)=1+\frac1y>1.
\]
Hence the absolute residual \(\AbsoluteValue{y_N+\log y_N-X}\) is already an upper bound for
the absolute inverse error.  This makes the logarithmic Lambert residual
particularly convenient.

\section{Asymptotic correctness of near-identity reversal}
The next statement explains why formal reversion coefficients are inherited by
an actual inverse.

\begin{theorem}[Transfer of a formal inverse expansion]\label{thm:analytic-inverse-transfer}
Suppose a real \(C^1\) function \(F\) is eventually strictly increasing and
\[
 F(y)=y+\LittleO(y),
 \qquad
 F'(y)=1+\LittleO(1)
 \qquad (y\to+\infty).
\]
Let \(G_N(X)\) be a finite polynomial--logarithmic expression such that
\begin{equation}\label{eq:inverse-residual-analytic}
 F(G_N(X))-X
 =\BigOAt{X^{-N-1}L^M}{X\to+\infty}
\end{equation}
for some \(M\), and suppose \(G_N(X)\to+\infty\).  Then the eventual inverse
\(G=\rev{F}\) satisfies
\begin{equation}\label{eq:inverse-error-transfer}
 G(X)-G_N(X)
 =\BigOAt{X^{-N-1}L^M}{X\to+\infty}.
\end{equation}
\end{theorem}
\begin{proof}
Because \(F'(y)=1+\LittleO(1)\), there is an eventual interval on which
\(\AbsoluteValue{F'(y)}\ge1/2\).  Both \(G(X)\) and \(G_N(X)\) lie in that interval for large
\(X\).  Apply \cref{prop:residual-error} with \(m=1/2\) and use
\cref{eq:inverse-residual-analytic}.
\end{proof}

This theorem reduces analytic asymptotic inversion to two tasks: prove an
actual inverse exists in the domain, and prove the formal residual has the
stated analytic size.

\section{Complex inverse problems}
In the complex plane, a lower bound on \(\AbsoluteValue{F'}\) along a line segment does not by
itself guarantee global injectivity.  Typical tools include:
\begin{itemize}
\item the analytic inverse function theorem on a chosen local domain;
\item Rouché's theorem to count a unique zero of \(F(y)-X\) in a disk;
\item Newton--Kantorovich estimates;
\item contraction mappings for a rearranged fixed-point equation;
\item sectorial continuation with explicit branch cuts.
\end{itemize}
Once a unique nearby branch is established, a residual estimate again controls
the error.  The formal coefficients remain the same; only domain and branch
selection change.

\section{Convergent, asymptotic, and divergent series}
Three situations should be distinguished.

\subsection{Convergent expansion}
The infinite series converges to the function on a nonempty domain.  The large
Lambert expansion has such a domain, although one often uses it asymptotically
rather than summing blindly.

\subsection{Divergent Poincaré expansion}
For every fixed \(N\), truncation after \(N\) terms gives the claimed remainder,
but the infinite series diverges.  The best numerical accuracy may occur at an
\(X\)-dependent optimal truncation order.

\subsection{Transseries with exponentially small sectors}
After all algebraic and logarithmic orders, the remainder may be controlled by
new scales such as \(\EulerE^{-X}\).  Different analytic sectors can then differ by
exponentially small terms, leading to Stokes phenomena and resurgence.  These
phenomena motivate full transseries but are not needed for the algebraic
operations developed here.

\begin{warningbox}
The label ``transseries'' does not imply divergence, and the label
``asymptotic expansion'' does not imply convergence.  Convergence is a theorem
about a particular function and domain, not a formal property of the notation.
\end{warningbox}

\section{Numerical evaluation strategy}
For a branch-aware inverse approximation:
\begin{enumerate}
\item evaluate the dominant and logarithmic coordinates with sufficient guard
      precision;
\item sum the formal terms in an order that minimizes cancellation;
\item compute a functional residual;
\item estimate the error from a derivative bound when available;
\item apply one or two Newton or Halley corrections;
\item verify that the iteration remains on the intended branch.
\end{enumerate}
At very large \(X\), evaluating a polynomial \(\LambertPolynomial{n}(L)\) by Horner's rule is
stable and efficient.  At moderate \(X\), the residual is more informative than
the apparent size of the last term.

\section{Branch cuts and symbolic simplification}
Identities such as
\[
 \log(ab)=\log a+\log b,
 \qquad
 (ab)^\alpha=a^\alpha b^\alpha
\]
are not globally valid for complex principal branches.  Formal monomial
factorization assumes a chosen logarithmic chart.  A symbolic implementation
should therefore distinguish:
\begin{itemize}
\item formal logarithms with declared additive branch constants;
\item principal numerical logarithms;
\item algebraic simplifications valid only under positivity assumptions.
\end{itemize}
For real \(X>0\) and positive leading constants, the distinction is often
invisible.  For Lambert branches it is essential.

\section{Documenting a remainder}
A useful final formula should state:
\begin{enumerate}
\item the limit, such as \(X\to+\infty\) or \(z\to\infty\) in a sector;
\item the chosen branch of every logarithm and inverse;
\item the formal truncation order, stated with \(\FormalBigOAt{t^N}{t}\);
\item any analytic remainder, stated separately with its limit, domain, and
      logarithmic factor;
\item whether the series is known to converge or is only asymptotic.
\end{enumerate}
For example,
\[
 W_0(z)=X-L+\sum_{n=1}^N\frac{\LambertPolynomial{n}(L)}{X^n}
 +\BigOAt{L^{N+1}/X^{N+1}}{z\to+\infty},
\]
with \(X=\log z\), \(L=\log X\), and \(z\to+\infty\), is an analytic
statement.  The notation
\[
 W=X-L+\sum_{n=1}^N\LambertPolynomial{n}(L)t^n+\FormalBigOAt{t^{N+1}}{t}
\]
is a formal block statement.

\begin{summarybox}
Formal algebra determines coefficient identities; analytic asymptotics require
domain, branch, and remainder control.  Standard Taylor estimates justify the
arithmetic operations when the small perturbations are uniformly small.
Functional residuals transfer directly to inverse errors when the derivative
is bounded away from zero.  Near-identity inverses with \(F'=1+\LittleO(1)\) inherit
the formal remainder order.  Convergence, divergence, and exponentially small
completion are separate analytic possibilities.
\end{summarybox}

\chapter{Pitfalls, diagnostics, and a decision procedure}\label{ch:pitfalls}

\section{Confusing the Lambert variables}
The large-argument Lambert expansion uses two nested coordinates:
\[
 X=\PrincipalLogarithm z+2\pi\ImaginaryUnit k,
 \qquad
 L=\log X.
\]
It is not an expansion directly in \(1/z\).  Replacing \(X\) by \(z\) produces
an approximation of the wrong scale.  On the principal positive branch,
\[
 W(z)\sim\log z-\log\log z,
\]
not \(z-\log z\).

A safe notation policy is to reserve \(z\) for the original Lambert argument,
\(X\) for its branch-adjusted logarithm, and \(L\) for the logarithm of \(X\).

\section{Treating \texorpdfstring{\(L\)}{L} as independent of
\texorpdfstring{\(X\)}{X}}
The symptoms are missing derivative terms in composition and incorrect
reversion coefficients.  The cure is to use
\[
 \frac{\Differential}{\Differential X}\bigl(a(L)X^{-n}\bigr)
 =(a'-na)X^{-n-1}
\]
or to substitute \(L=\log X\) before asking a computer algebra system to
differentiate.

In composition, remember that
\[
 L\longmapsto\log G(X),
\]
not merely \(L\).  A correct implementation never substitutes only the t-part
of an inner series.

\section{Assuming the polynomial coefficient ring is a field}
The reciprocal of \(L+t\) begins with \(L^{-1}\), so it is not in
\(\K[L]((t))\).  If a division unexpectedly produces denominators in \(L\),
this may be mathematically necessary rather than a simplification failure.

Perform the unit test before recurrence:
\begin{quote}
Is the first nonzero t-coefficient a nonzero constant in the chosen coefficient
ring?
\end{quote}
If not, enlarge the domain or keep the factor unexpanded.

\section{Truncating without guard terms}
Suppose an outer series contains \(t^{-1}\).  To compute a product modulo
\(t^N\), the other factor may be needed through t-order \(N\), not merely
\(N-1\), because \(t^{-1}t^N=t^{N-1}\).  Similar shifts occur in composition.

The remedy is valuation arithmetic.  If a factor has lowest exponent \(r\), an
unknown coefficient at order \(M\) can affect a product as low as \(M+r\).
Request enough input precision so that every omitted contribution lands at or
beyond the output cutoff.

\section{Using an unordered operator exponential}
The compact composition symbol
\[
 \mathopen{:}\EulerE^{\delta(\Dop-n)}\mathclose{:}
\]
is normal ordered.  Removing the colons changes the meaning when
\(\delta=\delta(L,t)\), because \(\Dop\delta\ne0\).  A naive symbolic operator
exponential then creates spurious derivatives of \(\delta\).

When in doubt, expand the definition
\[
 \sum_{r\ge0}\frac{\delta^r}{r!}(\Dop-n)^rf.
\]
It is unambiguous and finite at fixed t-precision.

\section{Composing in the wrong order}
\(F\compose G\) means \(F(G(X))\).  Since composition is noncommutative, a
silent reversal changes coefficients.  In code, use names such as
\texttt{compose(outer, inner)} and test against
\[
 X\compose G=G,
 \qquad
 F\compose X=F.
\]
For inverse verification, check both \(F\compose G\) and \(G\compose F\).

\section{Forgetting the relative smallness condition}
An additive perturbation \(H\) need not be numerically small; it must be small
relative to the dominant inner argument.  For
\[
 G=X-L,
\]
the correction \(-L\) tends to infinity in magnitude, but
\[
 \frac{-L}{X}\to0.
\]
Thus \(G=X(1-L/X)\) is a valid near-identity substitution.  Conversely,
\(G=X+X/2\) has a bounded relative perturbation but not a vanishing one; it
should be normalized as \((3/2)X\), not treated with a small-unit logarithm.

\section{Missing an iterated logarithm}
If the inner leading monomial contains \(L^\beta\), then
\[
 \log G=\cdots+\beta\log L+\cdots.
\]
Trying to force the result back into \(\K[L]((t))\) will either fail or silently
replace \(\log L\) by an unrelated symbol.  Factor the dominant inner monomial
before composition and inspect its logarithm.  This predicts the necessary
logarithmic depth.

\section{Misreading a formal remainder}
The statement
\[
 A=B+\FormalBigOAt{t^N}{t}
\]
means coefficient agreement below t-order \(N\).  It does not by itself assert
an analytic bound with a constant.  Conversely,
\[
 A(X)-B(X)=\BigOAt{X^{-N}L^M}{X\to+\infty}
\]
is an analytic statement and may hide a nontrivial logarithmic exponent.
Always label the formal t-adic remainder distinctly.

\section{Applying the large expansion near a branch point}
The large-argument Lambert expansion is not the correct local coordinate near
\(z=-\EulerE^{-1}\), where \(W\) has a square-root branch point.  There the natural
small variable is
\[
 p=\sqrt{2(\EulerE z+1)},
\]
and the expansion is a Puiseux series in \(p\), not a polynomial--logarithmic
series in \(1/\log z\).  An asymptotic method must begin with the local dominant
balance appropriate to the limit under study.

\section{Assuming more terms always improve a numerical approximation}
Even a convergent series may have nonmonotone partial errors, and a divergent
asymptotic series eventually worsens.  Use the functional residual, not only
the last term, to judge accuracy.  If the residual starts increasing, truncate
earlier or switch to Newton refinement.

\section{Losing branch constants during logarithm splitting}
For complex quantities,
\[
 \PrincipalLogarithm(ab)=\PrincipalLogarithm a+\PrincipalLogarithm b+2\pi\ImaginaryUnit m
\]
for an integer depending on the continuation path.  In Lambert calculations,
those constants are precisely what distinguish branch indices.  Keep
\(X=\PrincipalLogarithm z+2\pi\ImaginaryUnit k\) intact rather than simplifying it through identities
that assume principal branches compose globally.

\section{A symptom--cause--repair table}
\begin{longtable}{@{}p{0.29\textwidth}p{0.30\textwidth}p{0.31\textwidth}@{}}
\caption{Common diagnostic patterns.}\label{tab:diagnostics}\\
\toprule
Symptom & Likely cause & Repair \\
\midrule
\endfirsthead
\toprule
Symptom & Likely cause & Repair \\
\midrule
\endhead
Unexpected \(L^{-1}\) terms in a quotient & Nonconstant leading coefficient &
Use \(\K(L)\) or factor the leading log monomial. \\
\addlinespace
Inverse residual fails one order too early & Missing guard term or off-by-one
cutoff & Track valuations of all negative exponents. \\
\addlinespace
Composition lacks derivatives of coefficient polynomials & \(L\) treated as
independent & Replace \(L\) by \(\log G\) or use shifted jets. \\
\addlinespace
Spurious derivatives of the inner perturbation & Ordinary instead of
normal-ordered operator exponential & Use the explicit series in
\cref{eq:normal-ordered-def}. \\
\addlinespace
A \(\log\log X\) term appears & Inner dominant monomial contains a power of
\(\log X\) & Enlarge logarithmic depth. \\
\addlinespace
Left and right inverse checks disagree & Composition order bug or insufficient
precision & Test identity substitutions and add guard terms. \\
\addlinespace
Numerical iteration jumps branches & Branch-inconsistent initial coordinates &
Retain \(k\), branch cuts, and continuation path explicitly. \\
\addlinespace
Coefficient expressions grow explosively & Premature expansion or no
truncation & Truncate early; normalize blockwise; delay optional factoring. \\
\bottomrule
\end{longtable}

\section{A decision procedure for a new problem}
Given an implicit or inverse asymptotic problem, proceed as follows.

\subsection*{Step 1: identify the limit and dominant monomial}
Determine whether the variable tends to zero, infinity, a finite regular point,
or a branch point.  Solve the leading balance before choosing a series class.

\subsection*{Step 2: choose the fast and slow coordinates}
For Lambert \(W(z)\) at infinity, the fast coordinate is \(X=\PrincipalLogarithm z\), not
\(z\).  For the inverse of \(y/\log y\), the correction series is fast in
\(T=\log X\), with slow variable \(\log T\).

\subsection*{Step 3: predict closure}
Factor dominant monomials.  Inspect reciprocals and logarithms of those
factors.  Add rational log coefficients, roots, or iterated logs before the
coefficient calculation begins.

\subsection*{Step 4: choose the operation}
Use convolution for products, a triangular recurrence or Newton for division,
normal-ordered shifted derivatives for composition, and residual lifting,
Lagrange--Bürmann, or Newton for reversal.

\subsection*{Step 5: compute with explicit precision}
Write every intermediate result with a cutoff.  Derive a dependency bound so
that omitted terms cannot leak below it.

\subsection*{Step 6: verify formally}
Use a product or composition residual.  Prefer an independent formula when one
exists.

\subsection*{Step 7: justify analytically}
State the domain, branch, and remainder.  Transfer the residual to an error
bound if possible.

\begin{summarybox}
Most mistakes arise from one of five sources: an undersized coefficient field,
forgetting that \(L=\log X\), inadequate guard precision, ambiguous branch
algebra, or confusion between formal and analytic remainders.  A reliable
workflow begins with dominant balance and coordinate choice, performs closure
tests, computes in a valued algebra, verifies a residual, and only then assigns
analytic meaning.
\end{summarybox}

\chapter{Conclusion and further directions}\label{ch:conclusion}
Polynomial--logarithmic transseries provide a remarkably effective middle
ground.  They retain the finite, coefficientwise algorithms of ordinary power
series while encoding the slow logarithmic scales forced by asymptotic
inversion.  Their basic arithmetic is elementary, but composition reveals the
essential transserial mechanisms: the logarithmic variable moves, shifted
derivatives appear, support and branch conditions matter, and one operation
can force a larger monomial hierarchy.

The central formulas may be summarized in four lines.  Multiplication is
\[
 [t^n](AB)=\sum_{i+j=n}a_i(L)b_j(L).
\]
A reciprocal unit satisfies
\[
 b_0=a_0^{-1},
 \qquad
 b_n=-a_0^{-1}\sum_{k=1}^na_kb_{n-k}.
\]
Regular composition with \(G=\lambda X(1+u)\) is
\[
 F\compose G
 =\sum_n\lambda^{-n}X^{-n}
 \mathopen{:}\EulerE^{\log(1+u)(\Dop-n)}\mathclose{:}
 f_n(L+\log\lambda).
\]
Near-identity reversal is
\[
 \rev{(X+h)}
 =X+\sum_{m\ge1}\frac{(-1)^m}{m!}
 \left(\frac{\Differential}{\Differential X}\right)^{m-1}h^m.
\]
Everything else in the core theory is organized bookkeeping, closure analysis,
and verification around these identities.

Lambert \(W\) shows the method at its best.  Its coefficient polynomials can be
generated recursively, differentially, or combinatorially; all branches share
the same polynomials; the real lower branch follows by a sign change in the
large coordinate; and generalized equations \(y+\alpha\log y=X\) merely scale
the same family.  The example also teaches a broader lesson: the correct
asymptotic variable may itself be a logarithm of the original argument.

Several directions extend naturally from this foundation.
\begin{enumerate}
\item \textbf{Automatic scale discovery.}  Algorithms can infer whether the
next scale should be \(X^{-1}\), \((\log X)^{-1}\), \(\log\log X\), a root, or
an exponential monomial, following the exp--log multiseries tradition.
\item \textbf{Generalized supports.}  Hahn-series and grid-based methods allow
real exponents and several generators while preserving finite coefficient
dependencies.
\item \textbf{Iteration and conjugacy.}  Near-identity composition groups lead
to iterative logarithms, fractional iteration, formal normal forms, and Fatou
coordinates.
\item \textbf{Differential equations.}  The shifted derivative operators make
polynomial--logarithmic ansätze effective for resonant asymptotic ODEs.
\item \textbf{Exponentially small completion.}  When algebraic--logarithmic
orders are exhausted, full transseries add exponential sectors, Stokes data,
and resurgent structure.
\item \textbf{Certified symbolic computation.}  Exact coefficient domains,
formal valuations, residual certificates, and proof assistants can turn high-
order asymptotic computations into checkable mathematical artifacts.
\end{enumerate}

The practical message is simple: choose the asymptotic coordinates carefully,
work in the smallest class that is genuinely closed for the next operation,
and make the residual part of the answer.  With those habits, multiplication,
division, composition, and series reversal become systematic rather than
mysterious.

\appendix

\chapter{Bell polynomials and coefficient extraction}\label{app:bell}

\section{Definition}
The partial exponential Bell polynomials
\(\ExponentialPartialBellPolynomial{n}{k}(x_1,x_2,\ldots)\) are defined by
\begin{equation}\label{eq:bell-definition}
 \ExponentialPartialBellPolynomial{n}{k}(x_1,x_2,\ldots)
 =\sum
 \frac{n!}{m_1!m_2!\cdots}
 \prod_{j\ge1}\left(\frac{x_j}{j!}\right)^{m_j},
\end{equation}
where the sum runs over all nonnegative integer sequences satisfying
\[
 \sum_{j\ge1}m_j=k,
 \qquad
 \sum_{j\ge1}jm_j=n.
\]
Only \(x_1,\ldots,x_{n-k+1}\) can occur.  The conventions are
\[
 \ExponentialPartialBellPolynomial{0}{0}=1,
 \qquad
 \ExponentialPartialBellPolynomial{n}{0}=0\quad(n>0).
\]

Their exponential generating function is
\begin{equation}\label{eq:bell-egf}
 \frac1{k!}\left(\sum_{j\ge1}x_j\frac{t^j}{j!}\right)^k
 =\sum_{n\ge k}\ExponentialPartialBellPolynomial{n}{k}(x_1,x_2,\ldots)\frac{t^n}{n!}.
\end{equation}
This identity is exactly what is needed for powers of a t-series.

\section{First values}
The first nonzero partial Bell polynomials are
\begin{align*}
 \ExponentialPartialBellPolynomial{1}{1}&=x_1,\\
 \ExponentialPartialBellPolynomial{2}{1}&=x_2,
 &\ExponentialPartialBellPolynomial{2}{2}&=x_1^2,\\
 \ExponentialPartialBellPolynomial{3}{1}&=x_3,
 &\ExponentialPartialBellPolynomial{3}{2}&=3x_1x_2,
 &\ExponentialPartialBellPolynomial{3}{3}&=x_1^3,\\
 \ExponentialPartialBellPolynomial{4}{1}&=x_4,
 &\ExponentialPartialBellPolynomial{4}{2}&=4x_1x_3+3x_2^2,\\
 \ExponentialPartialBellPolynomial{4}{3}&=6x_1^2x_2,
 &\ExponentialPartialBellPolynomial{4}{4}&=x_1^4.
\end{align*}
A useful recurrence is
\begin{equation}\label{eq:bell-recurrence}
 \ExponentialPartialBellPolynomial{n}{k}
 =\sum_{j=1}^{n-k+1}
 \binom{n-1}{j-1}x_j\ExponentialPartialBellPolynomial{n-j}{k-1}.
\end{equation}
It follows by singling out the block containing a distinguished labeled
element in the combinatorial interpretation, or directly from the generating
function.

\section{Powers of a t-series}
Let
\[
 U(t)=\sum_{j\ge1}u_jt^j.
\]
Set \(x_j=j!u_j\) in \cref{eq:bell-egf}.  Then
\[
 \frac{U(t)^k}{k!}
 =\sum_{n\ge k}\frac1{n!}
 \ExponentialPartialBellPolynomial{n}{k}(1!u_1,2!u_2,\ldots)t^n.
\]
Therefore
\begin{equation}\label{eq:bell-U-power}
 [t^n]U^k
 =\frac{k!}{n!}\ExponentialPartialBellPolynomial{n}{k}(1!u_1,2!u_2,\ldots).
\end{equation}

\section{Logarithm and exponential coefficients}
For \(U\) of positive valuation,
\[
 \log(1+U)=\sum_{k\ge1}\frac{(-1)^{k+1}}kU^k.
\]
Hence
\begin{equation}\label{eq:bell-log-coeff}
 [t^n]\log(1+U)
 =\frac1{n!}\sum_{k=1}^n(-1)^{k+1}(k-1)!
 \ExponentialPartialBellPolynomial{n}{k}(1!u_1,2!u_2,\ldots).
\end{equation}
Similarly, if \(V=\sum_{j\ge1}v_jt^j\),
\begin{equation}\label{eq:bell-exp-coeff}
 [t^n]\EulerE^V
 =\frac1{n!}\sum_{k=0}^n
 \ExponentialPartialBellPolynomial{n}{k}(1!v_1,2!v_2,\ldots).
\end{equation}
The sum over \(k\) is the complete exponential Bell polynomial.

\section{Composition coefficients revisited}
For regular composition, let
\[
 \delta=\sum_{j\ge1}\delta_jt^j.
\]
The normal-ordered block
\[
 \mathopen{:}\EulerE^{\delta(\Dop-n)}\mathclose{:}f
 =\sum_{r\ge0}\frac{\delta^r}{r!}(\Dop-n)^rf
\]
has coefficient
\begin{equation}\label{eq:bell-normal-block}
 [t^m]\left(
 \mathopen{:}\EulerE^{\delta(\Dop-n)}\mathclose{:}f
 \right)
 =\frac1{m!}\sum_{r=0}^m
 \ExponentialPartialBellPolynomial{m}{r}(1!\delta_1,2!\delta_2,\ldots)
 (\Dop-n)^rf.
\end{equation}
After multiplying by \(\lambda^{-n}t^n\) and summing over outer blocks, one
recovers \cref{eq:composition-bell}.

\section{Faà di Bruno's formula}
Bell polynomials also express ordinary higher derivatives of a composite:
\begin{equation}\label{eq:faa-di-bruno}
 \frac{\Differential^n}{\Differential x^n}f(g(x))
 =\sum_{k=1}^nf^{(k)}(g(x))
 \ExponentialPartialBellPolynomial{n}{k}\bigl(g'(x),g''(x),\ldots,g^{(n-k+1)}(x)\bigr).
\end{equation}
The composition formula in this article is a close formal relative: Bell
polynomials collect repeated powers of the logarithmic increment, while the
shifted derivative \(\Dop-n\) simultaneously handles the outer coefficient and
its power prefactor.

\chapter{Deriving the universal composition formulas}\label{app:universal-derivation}

\section{Inner expansions}
Let
\[
 G=X+A+Bq+Ct^2+Dt^3+\FormalBigOAt{t^4}{t}
 =X(1+u),
\]
where
\[
 u=Aq+Bt^2+Ct^3+Dt^4+\FormalBigOAt{t^5}{t}.
\]
The logarithmic increment is
\begin{equation}\label{eq:appendix-delta}
\begin{split}
 \delta={}&At+\left(B-\frac{A^2}{2}\right)t^2
 +\left(C-AB+\frac{A^3}{3}\right)t^3\\
 &+\left(D-AC-\frac{B^2}{2}+A^2B-\frac{A^4}{4}\right)t^4
 +\FormalBigOAt{t^5}{t}.
\end{split}
\end{equation}
The reciprocal powers needed through t-order three are
\begin{align}
 G^{-1}={}&t-At^2+(A^2-B)t^3
 +(-A^3+2AB-C)t^4+\FormalBigOAt{t^5}{t},\label{eq:Ginv-app}\\
 G^{-2}={}&t^2-2At^3+(3A^2-2B)t^4+\FormalBigOAt{t^5}{t},\label{eq:Ginv2-app}\\
 G^{-3}={}&t^3-3At^4+\FormalBigOAt{t^5}{t}.\label{eq:Ginv3-app}
\end{align}

\section{Outer coefficient Taylor expansions}
For a coefficient function \(a\),
\begin{equation}\label{eq:a-taylor-app}
 a(L+\delta)
 =a+a'\delta+\frac12a''\delta^2+\frac16a'''\delta^3
 +\FormalBigOAt{t^4}{t}.
\end{equation}
For \(b,c,d\), fewer derivatives are needed after multiplication by
\(G^{-1},G^{-2},G^{-3}\):
\begin{align*}
 b(L+\delta)&=b+b'\delta+\frac12b''\delta^2+\FormalBigOAt{t^3}{t},\\
 c(L+\delta)&=c+c'\delta+\FormalBigOAt{t^2}{t},\\
 d(L+\delta)&=d+\FormalBigOAt{t}{t}.
\end{align*}
Substitute \cref{eq:appendix-delta} and collect.  For example,
\begin{align*}
 a(L+\delta)
={}&a+Aa't\\
&+\left[
 \left(B-\frac{A^2}{2}\right)a'
 +\frac{A^2}{2}a''
 \right]t^2\\
&+\Biggl[
 \left(C-AB+\frac{A^3}{3}\right)a'
 +A\left(B-\frac{A^2}{2}\right)a''
 +\frac{A^3}{6}a'''
 \Biggr]t^3
 +\FormalBigOAt{t^4}{t}.
\end{align*}

\section{Collecting the outer series}
Write
\[
 F(G)=G+a(\log G)+b(\log G)G^{-1}
 +c(\log G)G^{-2}+d(\log G)G^{-3}+\cdots.
\]
At t-order zero, only \(G\) and \(a\) contribute, giving
\[
 h_0=A+a.
\]
At t-order one, contributions are
\[
 B+Aa'+b,
\]
which gives \cref{eq:h1}.  At t-order two, use
\cref{eq:Ginv-app,eq:Ginv2-app} and the Taylor expansions to obtain
\[
 h_2=-\frac12A^2a'+\frac12A^2a''-Ab+Ab'+Ba'+C+c.
\]
At t-order three, the contributions are:
\begin{itemize}
\item \(D\) from \(G\);
\item the t-cubic Taylor block of \(a(L+\delta)\);
\item the product of \(b(L+\delta)\) with \(G^{-1}\) through t-cubic order;
\item the product of \(c(L+\delta)\) with \(G^{-2}\) through t-cubic order;
\item \(d\) from \(d(L+\delta)G^{-3}\).
\end{itemize}
Their sum simplifies to \cref{eq:h3}.  This derivation also explains the
highest derivatives that can occur at each order.

\section{Solving for the inverse}
Set \(h_j=0\) successively.  The equations have the schematic form
\[
 h_n=\text{known expression in }A,\ldots,A_{n-1}+A_n+f_n,
\]
so each new inverse block enters with coefficient one.  This is the explicit
version of \cref{lem:leading-correction}.  Solving the first four equations
produces \cref{eq:inverse-A,eq:inverse-B,eq:inverse-C,eq:inverse-D}.

\chapter{Exercises with solutions}\label{app:exercises}

\section{Exercises}

\begin{exercise}[Product]
Compute through t-order three:
\[
 (1+Lt+t^2)(1-Lt+(L^2+2)t^2+Lt^3).
\]
\end{exercise}

\begin{exercise}[Reciprocal]
Find \((1+(L+1)t+Lt^2)^{-1}\) through t-order three.
\end{exercise}

\begin{exercise}[Differentiation]
Differentiate
\[
 A=X^2+L^3+\frac{L^2-1}{X^2}
\]
blockwise using shifted derivatives.
\end{exercise}

\begin{exercise}[A logarithm]
Compute
\[
 \log\left(1+(L-1)t+(L+2)t^2\right)
\]
through t-order two.
\end{exercise}

\begin{exercise}[Composition]
Let
\[
 F(Y)=Y+(\log Y)^2,
 \qquad
 G(X)=X-L.
\]
Compute \(F\compose G\) through t-order two.
\end{exercise}

\begin{exercise}[Closure]
Explain why \(F(Y)=\log Y\) composed with
\(G(X)=X^2(\log X)^3\) does not remain in \(\K[L]((X^{-1}))\).
Identify the smallest new logarithmic symbol forced by the composition.
\end{exercise}

\begin{exercise}[Reversion]
Use the universal formulas to reverse
\[
 F(X)=X+2L+\frac3X+\FormalBigOAt{t^2}{t}
\]
through t-order two.
\end{exercise}

\begin{exercise}[Lagrange formula]
Apply \cref{eq:lagrange-infinity} to \(F(X)=X+\alpha\log X\) and prove that the
coefficient of \(X^{-n}\) is \(\alpha^{n+1}\LambertPolynomial{n}(L)\).
\end{exercise}

\begin{exercise}[Lambert polynomial]
Compute \(\LambertPolynomial{4}(L)\) from the Stirling formula, using
\[
 \stirlingone{4}{4}=1,
 \quad \stirlingone{4}{3}=6,
 \quad \stirlingone{4}{2}=11,
 \quad \stirlingone{4}{1}=6.
\]
\end{exercise}

\begin{exercise}[Lower Lambert branch]
Starting from \(u-\log u=T\), derive the first three correction blocks for
\(u\), and hence for \(W_{-1}(x)=-u\) as \(x\to0^-\).
\end{exercise}

\begin{exercise}[Small-variable inverse]
Find the inverse of
\[
 f(x)=x+a x^2\log(1/x)+b x^3
\]
through x-order three.
\end{exercise}

\begin{exercise}[Residual bound]
Let \(F(y)=y+\log y\) for \(y>0\).  Show that any positive approximation
\(y_N\) to the solution of \(F(y)=X\) satisfies
\[
 \AbsoluteValue{y_N-y_*}\le\AbsoluteValue{y_N+\log y_N-X}.
\]
\end{exercise}

\section{Solutions}

\subsection*{Solution 1}
The t coefficient vanishes.  The t-squared coefficient is
\[
 (L^2+2)-L^2+1=3.
\]
The t-cubic coefficient is
\[
 L+L(L^2+2)-L=L^3+2L.
\]
Thus the product is
\[
 1+3t^2+(L^3+2L)t^3+\FormalBigOAt{t^4}{t}.
\]

\subsection*{Solution 2}
Let the reciprocal coefficients be \(b_n\).  Then
\[
 b_0=1,
 \qquad b_1=-(L+1),
\]
\[
 b_2=(L+1)^2-L=L^2+L+1,
\]
and
\begin{align*}
 b_3&=-\bigl((L+1)b_2+Lb_1\bigr)\\
 &=-L^3-L^2-L-1.
\end{align*}
Therefore
\[
 \frac1{1+(L+1)t+Lt^2}
 =1-(L+1)t+(L^2+L+1)t^2
 -(L^3+L^2+L+1)t^3+\FormalBigOAt{t^4}{t}.
\]

\subsection*{Solution 3}
The derivative of \(X^2\) is \(2X\).  For \(L^3\), use the n-zero block:
\[
 \frac{\Differential}{\Differential X}L^3=3L^2X^{-1}.
\]
For \((L^2-1)X^{-2}\), use \(\Dop-2\):
\[
 \bigl(2L-2(L^2-1)\bigr)X^{-3}
 =(-2L^2+2L+2)X^{-3}.
\]
Hence
\[
 A'=2X+\frac{3L^2}{X}
 +\frac{-2L^2+2L+2}{X^3}.
\]

\subsection*{Solution 4}
Put \(U=(L-1)t+(L+2)t^2\).  Then
\[
 \log(1+U)=U-\frac12U^2+\FormalBigOAt{t^3}{t},
\]
so
\[
 \log(1+U)
 =(L-1)t+
 \left(L+2-\frac{(L-1)^2}{2}\right)t^2+\FormalBigOAt{t^3}{t}.
\]

\subsection*{Solution 5}
Here \(a(L)=L^2\), \(A=-L\), and all other displayed coefficient functions
vanish.  The universal formulas give
\[
 h_0=L^2-L,
 \qquad
 h_1=Aa'=-2L^2,
\]
\[
 h_2=-\frac12A^2a'+\frac12A^2a''
 =-L^3+L^2.
\]
Thus
\[
 F(G)=X+(L^2-L)-\frac{2L^2}{X}
 +\frac{-L^3+L^2}{X^2}+\FormalBigOAt{t^3}{t}.
\]

\subsection*{Solution 6}
\[
 \log G=2L+3\log L.
\]
The new symbol is \(L_2=\log L=\log\log X\).  Therefore the output lies in a
two-log-level coefficient class, not in \(\K[L]((t))\).

\subsection*{Solution 7}
Use \(a=2L\), \(b=3\), \(c=0\).  Then
\[
 A=-2L,
 \qquad
 B=aa'-b=4L-3.
\]
Also
\begin{align*}
 C&=\frac12a^2a'-a(a')^2-ab+a'b\\
 &=4L^2-8L-6L+6\\
 &=4L^2-14L+6.
\end{align*}
Hence
\[
 \rev F(X)=X-2L+\frac{4L-3}{X}
 +\frac{4L^2-14L+6}{X^2}+\FormalBigOAt{t^3}{t}.
\]

\subsection*{Solution 8}
In \cref{eq:lagrange-infinity}, the m-th term is
\[
 \frac{(-1)^m\alpha^m}{m!}
 \left(\frac{\Differential}{\Differential X}\right)^{m-1}L^m.
\]
It has exactly t-order \(m-1\).  Put \(n=m-1\) and compare with
\cref{eq:Pn-operator}; the coefficient is \(\alpha^{n+1}\LambertPolynomial{n}(L)\).

\subsection*{Solution 9}
Formula~\eqref{eq:Lambert-Stirling} gives
\begin{align*}
 \LambertPolynomial{4}(L)
 &=-\stirlingone{4}{4}L
 +\stirlingone{4}{3}\frac{L^2}{2!}
 -\stirlingone{4}{2}\frac{L^3}{3!}
 +\stirlingone{4}{1}\frac{L^4}{4!}\\
 &=-L+3L^2-\frac{11}{6}L^3+\frac14L^4\\
 &=\frac{L(3L^3-22L^2+36L-12)}{12}.
\end{align*}

\subsection*{Solution 10}
Reversion of \(u-\log u=T\) is the case \(\alpha=-1\):
\[
 u=T+L+\frac{L}{T}-\frac{L(L-2)}{2T^2}
 +\frac{L(2L^2-9L+6)}{6T^3}+\cdots.
\]
Therefore \(W_{-1}(x)=-u\) with
\(T=\log(1/(-x))\), giving \cref{eq:Wm1-first}.

\subsection*{Solution 11}
Let \(h(x)=a x^2\ell+b x^3\).  The first reversion term is
\[
 -h=-a x^2\ell-bx^3.
\]
Only the leading part \(a x^2\ell\) of \(h\) contributes to x-order three in
the second term:
\[
 \frac12(h^2)'=a^2x^3(2\ell^2-\ell)+\FormalBigO_x(x^4).
\]
Thus
\[
 \rev f(x)=x-a x^2\ell+
 \bigl(a^2(2\ell^2-\ell)-b\bigr)x^3
 +\FormalBigO_x(x^4).
\]

\subsection*{Solution 12}
Since \(F'(y)=1+1/y>1\) for \(y>0\), apply
\cref{prop:residual-error} with \(m=1\).

\chapter{Notation and quick reference}\label{app:notation}

\begin{longtable}{@{}p{0.22\textwidth}p{0.68\textwidth}@{}}
\toprule
Symbol & Meaning \\
\midrule
\endfirsthead
\toprule
Symbol & Meaning \\
\midrule
\endhead
\(X\) & Large asymptotic variable. \\
\(t=X^{-1}\) & Small t-adic variable. \\
\(L=\log X\) & First slow logarithmic variable. \\
\(L_j\) & Iterated logarithm, \(L_{j+1}=\log L_j\). \\
\(\PolynomialLogLaurentRing{\K}\) & Polynomial--logarithmic Laurent ring \(\K[L]((t))\). \\
\(\RationalLogLaurentField{\K}\) & Rational-log coefficient field \(\K(L)((t))\). \\
\(\val_t A\) & Lowest t-exponent of a nonzero series. \\
\([A]_{<N}\) & Truncation to t-exponents below \(N\). \\
\(\FormalBigOAt{t^N}{t}\) & Formal t-adic remainder of valuation at least \(N\). \\
\(\Dop\) & Derivative \(\Differential/\Differential L\). \\
\(\Delta_n\) & Shifted derivative \(\Dop-n\). \\
\(\recip A\) & Multiplicative inverse \(1/A\). \\
\(\rev F\) & Compositional inverse of \(F\). \\
\(F\compose G\) & Composition \(F(G(X))\). \\
\(\delta\) & Logarithmic increment \(\log(1+u)\) for
\(G=\lambda X(1+u)\). \\
\(\mathopen{:}\EulerE^{\delta(\Dop-n)}\mathclose{:}\) & Normal-ordered shifted
Taylor operator. \\
\(\ExponentialPartialBellPolynomial{n}{k}\) & Partial exponential Bell polynomial. \\
\(\stirlingone{n}{k}\) & Unsigned Stirling number of the first kind. \\
\(\LambertPolynomial{n}(L)\) & n-th Lambert polynomial in the expansion of \(W\). \\
\(X=\xi_k\) & Branch coordinate \(\PrincipalLogarithm z+2\pi\ImaginaryUnit k\) for \(W_k(z)\). \\
\bottomrule
\end{longtable}


\backmatter

\chapter*{Source note and further reading}
\addcontentsline{toc}{chapter}{Source note and further reading}

This article was developed from the collection of transseries papers, lecture
notes, and tutorial drafts supplied with the project, supplemented by the
primary literature and authoritative reference sources listed below.  The local
materials were especially useful for the general Hahn--transseries viewpoint
and for Edgar's treatment of composition.  The online sources were used to
verify bibliographic data, branch conventions, the all-orders Lambert
expansion, and the computer-algebra literature on functional inversion.

The presentation is deliberately narrower than a full construction of the
field of logarithmic--exponential transseries.  Its main object is the
computable sector generated by inverse powers of a large variable and
polynomials (or controlled extensions) in one or more logarithms.  Readers who
want the complete ordered differential-field theory should begin with
\cite{edgar-beginners,vdhoeven-book,adh-book}; readers chiefly interested in
algorithmic asymptotics and inverse functions should also consult
\cite{salvy-fast,salvy-shackell,rssv-exp-log}.

\begin{thebibliography}{99}

\bibitem{adh-book}
M.~Aschenbrenner, L.~van den Dries, and J.~van der Hoeven,
\emph{Asymptotic Differential Algebra and Model Theory of Transseries},
Annals of Mathematics Studies, vol.~195, Princeton University Press,
Princeton, 2017.
\newblock DOI: \href{https://doi.org/10.1515/9781400885411}{10.1515/9781400885411}.

\bibitem{corless-lambert}
R.~M. Corless, G.~H. Gonnet, D.~E.~G. Hare, D.~J. Jeffrey, and D.~E. Knuth,
``On the Lambert $W$ function,''
\emph{Advances in Computational Mathematics} \textbf{5} (1996), 329--359.
\newblock DOI: \href{https://doi.org/10.1007/BF02124750}{10.1007/BF02124750}.

\bibitem{dlmf-lambert}
NIST Digital Library of Mathematical Functions,
\S4.13, ``Lambert $W$-Function,'' release 1.2.7, June 15, 2026.
\newblock \url{https://dlmf.nist.gov/4.13} (accessed August 31, 2026).

\bibitem{edgar-beginners}
G.~A. Edgar,
``Transseries for beginners,''
\emph{Real Analysis Exchange} \textbf{35} (2009/2010), no.~2, 253--310.
\newblock arXiv: \href{https://arxiv.org/abs/0801.4877}{0801.4877}.

\bibitem{edgar-composition}
G.~A. Edgar,
``Transseries: composition, recursion, and convergence,''
preprint, 2009.
\newblock arXiv: \href{https://arxiv.org/abs/0909.1259}{0909.1259}.

\bibitem{jeffrey-inversion}
D.~J. Jeffrey, R.~M. Corless, D.~E.~G. Hare, and D.~E. Knuth,
``On the inversion of $y^\alpha \EulerE^y$ in terms of associated Stirling
numbers,''
\emph{Comptes Rendus de l'Acad\'emie des Sciences, S\'erie I -- Math\'ematique}
\textbf{320} (1995), no.~12, 1449--1452.
\newblock arXiv: \href{https://arxiv.org/abs/math/9512230}{math/9512230}.

\bibitem{mardesic-dulac}
P.~Marde\v{s}i\'c, M.~Resman, J.-P.~Rolin, and V.~\v{Z}upanovi\'c,
``Length of epsilon-neighborhoods of orbits of Dulac maps,''
preprint, 2016.
\newblock arXiv: \href{https://arxiv.org/abs/1606.02581}{1606.02581}.

\bibitem{mardesic-powerlog}
P.~Marde\v{s}i\'c, M.~Resman, J.-P.~Rolin, and V.~\v{Z}upanovi\'c,
``Normal forms and embeddings for power-log transseries,''
\emph{Advances in Mathematics} \textbf{303} (2016), 888--953.
\newblock DOI:
\href{https://doi.org/10.1016/j.aim.2016.08.030}{10.1016/j.aim.2016.08.030}.

\bibitem{peran-logarithmic}
D.~Peran,
``Normalizations of strongly hyperbolic logarithmic transseries and complex
Dulac germs,''
\emph{Analysis and Mathematical Physics} \textbf{13} (2023), article~66.
\newblock DOI:
\href{https://doi.org/10.1007/s13324-023-00830-w}{10.1007/s13324-023-00830-w}.

\bibitem{rssv-exp-log}
D.~Richardson, B.~Salvy, J.~Shackell, and J.~van der Hoeven,
``Asymptotic expansions of exp-log functions,''
in Y.~N. Lakshman (ed.), \emph{Proceedings of ISSAC '96}, ACM Press,
New York, 1996, pp.~309--313.
\newblock DOI:
\href{https://doi.org/10.1145/236869.237089}{10.1145/236869.237089}.

\bibitem{salvy-fast}
B.~Salvy,
``Fast computation of some asymptotic functional inverses,''
\emph{Journal of Symbolic Computation} \textbf{17} (1994), no.~3,
227--236.
\newblock DOI:
\href{https://doi.org/10.1006/jsco.1994.1014}{10.1006/jsco.1994.1014}.

\bibitem{salvy-shackell}
B.~Salvy and J.~Shackell,
``Symbolic asymptotics: multiseries of inverse functions,''
\emph{Journal of Symbolic Computation} \textbf{27} (1999), no.~6,
543--563.
\newblock DOI:
\href{https://doi.org/10.1006/jsco.1999.0281}{10.1006/jsco.1999.0281}.

\bibitem{vdhoeven-book}
J.~van der Hoeven,
\emph{Transseries and Real Differential Algebra},
Lecture Notes in Mathematics, vol.~1888, Springer, Berlin--Heidelberg, 2006.
\newblock DOI:
\href{https://doi.org/10.1007/3-540-35590-1}{10.1007/3-540-35590-1}.

\end{thebibliography}

\end{document}
